
\documentclass[12pt,a4paper]{article}

\usepackage[utf8]{inputenc}
\usepackage[T1]{fontenc}
\usepackage[brazil]{babel}
\usepackage{geometry}
\usepackage{titlesec}
\usepackage{setspace}
\usepackage{longtable}
\usepackage{array}
\usepackage{hyperref}
\usepackage{xcolor}
\usepackage{amsmath}
\usepackage{amsfonts}
\usepackage{mathrsfs}
\usepackage{siunitx}
\usepackage{pgfplots}
\pgfplotsset{compat=1.18}
\usepackage{float}

\sisetup{
    output-decimal-marker = {,},
    per-mode = symbol
}

\newcommand{\resposta}[1]{\noindent\textbf{Resposta:} #1}


\geometry{margin=2.5cm}
\setstretch{1.3}

\title{\textbf{Resolução dos Exercícios dos Capítulos 31, 32, 33 e 34}}
\author{
Ana Alicy \\
Kaylane Castro \\
Dávyla Maria \\
Rodrigo Albuquerque \\
Fábio Rodrigues \\
Mateus Valentim \\
Cicero Rodrigues \\
Guilherme Araújo Floriano \\
Arthur Roberto da Silva
}
\date{}

\begin{document}

% ================= CAPA =================

\begin{titlepage}
    \centering

    {\Large Universidade Federal do Ceará \par}
    \vspace{1.5cm}

    {\Large Ondas e Antenas \par}
    \vspace{1.5cm}

    {\huge \textbf{Resolução dos Exercícios do Livro Fundamentos da Física - Halliday e Resnick, Vol 3 e 4, 8 edição. Capítulos 31, 32, 33 e 34} \par}
    \vspace{2cm}

    {\Large Professor: Joel Ramiro de Castro \par}
    \vspace{2cm}

    \textbf{Integrantes:}

    \vspace{0.5cm}

    \begin{flushleft}
    Ana Alicy \\
    Kaylane Castro \\
    Dávyla Maria \\
    Rodrigo Albuquerque \\
    Fábio Rodrigues \\
    Mateus Valentim \\
    Cicero Rodrigues \\
    Guilherme Araújo Floriano \\
    Arthur Roberto da Silva
    \end{flushleft}

    \vfill

    {\large Quixadá -- CE \\
    2026}

\end{titlepage}

\tableofcontents
\newpage


% ================= CAPÍTULO 31 =================

\section{Capítulo 31}

\subsection{Questão 7 — Ana Alicy}

Na Fig. 31-28, $R = \SI{14,0}{\ohm}$, $C = \SI{6,20}{\micro\farad}$ e $L = \SI{54,0}{\milli\henry}$ e a fonte ideal tem uma força eletromotriz $\mathscr{E} = \SI{34,0}{\volt}$. A chave é mantida na posição $a$ por um longo tempo antes de ser colocada na posição $b$. Determine (a) a frequência e (b) a amplitude da corrente no circuito depois que a chave é colocada na posição $b$.

\vspace{0.2cm}


\textbf{Solução:}
\begin{quote}

\noindent \textbf{1. Entendendo a Física do Circuito}

\vspace{0.2cm}
\noindent \textbf{Momento 1 (Chave em $a$):} \\
Quando a chave fica na posição $a$ por muito tempo, o circuito ativo é formado apenas pela fonte $\mathscr{E}$, o resistor $R$ e o capacitor $C$. A corrente flui até carregar totalmente o capacitor. Quando ele está totalmente carregado, a corrente cessa. Nesse momento, a diferença de potencial no capacitor é máxima e igual à da fonte de tensão.
\begin{itemize}
    \item Tensão máxima no capacitor ($V_C$): \SI{34,0}{\volt}
\end{itemize}

\vspace{0.2cm}
\noindent \textbf{Momento 2 (Chave em $b$):} \\
Ao mudar a chave para $b$, a bateria e o resistor ficam desconectados. O circuito agora é formado apenas pelo capacitor $C$ (totalmente carregado) e o indutor $L$. Tem a formação de um circuito LC ideal, que começará a oscilar eletromagneticamente.

\vspace{0.4cm}
\noindent \textbf{(a) Qual a frequência da corrente no circuito?} \\
Em um circuito LC, a frequência natural de oscilação ($f$) depende apenas da indutância ($L$) e da capacitância ($C$), e é dada pela fórmula:
\[
f = \frac{1}{2\pi \sqrt{LC}}
\]

\noindent Organizando os dados no Sistema Internacional (SI):
\begin{itemize}
    \item $L = \SI{54,0}{\milli\henry} = 54,0 \times 10^{-3}\ \text{H}$
    \item $C = \SI{6,20}{\micro\farad} = 6,20 \times 10^{-6}\ \text{F}$
\end{itemize}

\noindent Substituindo na equação:
\[
f = \frac{1}{2\pi \sqrt{(54,0 \times 10^{-3}) \cdot (6,20 \times 10^{-6})}}
\]
\[
f = \frac{1}{2\pi \sqrt{334,8 \times 10^{-9}}}
\]
\[
f = \frac{1}{2\pi \cdot (5,786 \times 10^{-4})}
\]
\[
f \approx \frac{1}{3,635 \times 10^{-3}} \approx \SI{275,06}{\hertz}
\]

\noindent Arredondando para três algarismos significativos (padrão dos dados da questão):

\resposta{A frequência é de \SI{275}{\hertz}.}
\vspace{0.2cm}

\noindent \textbf{(b) Qual a amplitude da corrente no circuito?} \\
A amplitude da corrente ($I$) é a corrente máxima que passa pelo indutor. Para descobrir isso, usa o Princípio da Conservação de Energia. Em um circuito LC, a energia total do sistema fica alternando entre energia elétrica (armazenada no capacitor) e energia magnética (armazenada no indutor).

\vspace{0.2cm}
\noindent A energia elétrica máxima inicial no capacitor é igual à energia magnética máxima no indutor:
\[
U_{E,\text{máx}} = U_{B,\text{máx}}
\]
\[
\frac{1}{2} C \mathscr{E}^2 = \frac{1}{2} L I^2
\]

\noindent Simplificando o fator $\frac{1}{2}$ de ambos os lados e isolando a amplitude da corrente ($I$):
\[
I^2 = \frac{C}{L} \mathscr{E}^2
\]
\[
I = \mathscr{E} \sqrt{\frac{C}{L}}
\]

\noindent Substituindo os valores numéricos:
\[
I = 34,0 \cdot \sqrt{\frac{6,20 \times 10^{-6}}{54,0 \times 10^{-3}}}
\]
\[
I = 34,0 \cdot \sqrt{1,148 \times 10^{-4}}
\]
\[
I = 34,0 \cdot (0,01071)
\]
\[
I \approx \SI{0,364}{\ampere}
\]

\noindent Arredondando para três algarismos significativos:

\resposta{A amplitude da corrente é de \SI{0,364}{\ampere} (ou \SI{364}{\milli\ampere}).}
\end{quote}
\vspace{0.2cm}


\subsection{Questão 10 — Kaylane Castro}

\subsection*{Enuciado}
Considere um circuito de malha única formado por vários indutores, resistores e capacitores ligados em série. Deseja-se mostrar que, independentemente da ordem desses elementos, o circuito possui o mesmo comportamento de um circuito $LCR$ simples.

\vspace{0.3cm}

Como o circuito possui apenas uma malha, a corrente $i(t)$é a mesma em todos os elementos.

Aplicando a Lei das Malhas de Kirchhoff:

\[
\sum \Delta V = 0
\]

As quedas de tensão em cada elemento são:

\[
\Delta V_L = L_j \frac{di}{dt}
\]

\[
\Delta V_R = R_j i
\]

\[
\Delta V_C = \frac{q}{C_j}
\]

Logo, somando todas as tensões da malha:

\[
-\sum_j L_j \frac{di}{dt}
-\sum_j R_j i
-\sum_j \frac{q}{C_j}
=0
\]

Multiplicando a equação por $-1$:

\[
\left( \sum_j L_j \right)\frac{di}{dt}
+
\left( \sum_j R_j \right)i
+
\left( \sum_j \frac{1}{C_j} \right) q
=0
\]

Definindo os parâmetros equivalentes:

\[
L = \sum_j L_j
\]

\[
R = \sum_j R_j
\]

\[
\frac{1}{C} = \sum_j \frac{1}{C_j}
\]

A equação passa a ser:

\[
L\frac{di}{dt}+Ri+\frac{q}{C}=0
\]

Sabendo que:

\[
i=\frac{dq}{dt}
\]

então:

\[
\frac{di}{dt}=\frac{d^2q}{dt^2}
\]

Substituindo na equação anterior:

\[
L\frac{d^2q}{dt^2}
+
R\frac{dq}{dt}
+
\frac{1}{C}q
=0
\]

Essa é exatamente a equação diferencial de um circuito $LCR$ série simples.

Além disso, como soma de resistores, indutores e capacitores equivalentes não depende da ordem em que os elementos aparecem no circuito, conclui-se que o comportamento do circuito também não depende da sequência dos componentes.

Portanto, qualquer circuito de malha única formado por elementos $L$, $C$ e $R$ em série é equivalente a um circuito $LCR$ simples. 



\subsection{Questão 11 — Dávyla Maria}
Na Fig. 31.10, R = 14,0 Ω, C = 6,20 μF, L = 54,0 mH e a fonte ideal tem uma força eletromotriz = 34,0 V. A chave é mantida na posição a por um longo tempo e depois é colocada na posição b. Determine (a) a frequência e (b) a amplitude das oscilações resultantes.

Resolução:


a)

\subsection{Questão 12 — Rodrigo Albuquerque}
A frequência de oscilação de um circuito LC é dada por

\[
f=\frac{1}{2\pi\sqrt{LC}}
\]

Dados do problema:

\[
L=10\,mH=10\times10^{-3}\,H
\]

\[
C_1=5,0\,\mu F=5\times10^{-6}\,F
\]

\[
C_2=2,0\,\mu F=2\times10^{-6}\,F
\]

\subsection*{(a) Menor frequência}

A menor frequência ocorre quando a capacitância é máxima, então os capacitores ficam em paralelo.

\[
C_{eq}=C_1+C_2
\]

\[
C_{eq}=5+2=7\,\mu F
\]

Substituindo:

\[
f=\frac{1}{2\pi\sqrt{(10\times10^{-3})(7\times10^{-6})}}
\]

\[
f\approx600\,Hz
\]

\[
f\approx600\,Hz
\]

\subsection*{(b) Segunda menor frequência}

Usando apenas o capacitor de \(5\,\mu F\):

\[
f=\frac{1}{2\pi\sqrt{(10\times10^{-3})(5\times10^{-6})}}
\]

\[
f\approx712\,Hz
\]

\[
f\approx712\,Hz
\]

\subsection*{(c) Segunda maior frequência}

Usando apenas o capacitor de \(2\,\mu F\):

\[
f=\frac{1}{2\pi\sqrt{(10\times10^{-3})(2\times10^{-6})}}
\]

\[
f\approx1125\,Hz
\]

\[
f\approx1125\,Hz
\]

\subsection*{(d) Maior frequência}

A maior frequência acontece com a menor capacitância equivalente, ou seja, capacitores em série.

\[
\frac{1}{C_{eq}}=\frac{1}{C_1}+\frac{1}{C_2}
\]

\[
C_{eq}=\frac{C_1C_2}{C_1+C_2}
\]

\[
C_{eq}=\frac{5\cdot2}{5+2}
\]

\[
C_{eq}=\frac{10}{7}\,\mu F
\]

\[
C_{eq}\approx1,43\,\mu F
\]

Substituindo na fórmula:

\[
f=\frac{1}{2\pi\sqrt{(10\times10^{-3})(1,43\times10^{-6})}}
\]

\[
f\approx1330\,Hz
\]

\[
f\approx1330\,Hz
\]
\subsection{Questão 13 — Fábio Rodrigues}
\textbf{Dados:}

\textbf{Dados:}

\[
C = 1,0 \times 10^{-9}\,\text{F}
\]

\[
L = 3,0 \times 10^{-3}\,\text{H}
\]

\[
V_{\max} = 3,0\,\text{V}
\]

\textbf{(a) Carga máxima do capacitor}

Usamos:

\[
Q_{\max} = CV_{\max}
\]

Substituindo:

\[
Q_{\max} = (1,0 \times 10^{-9})(3,0)
\]

\[
Q_{\max} = 3,0 \times 10^{-9}\,\text{C}
\]

\textbf{Resultado:}

\[
\boxed{Q_{\max} = 3,0\,\text{nC}}
\]

\textbf{(b) Corrente máxima do circuito}

Num circuito LC, a energia oscila entre capacitor e indutor.

A energia máxima no capacitor é igual à energia máxima no indutor:

\[
\frac{1}{2}CV_{\max}^{2} = \frac{1}{2}LI_{\max}^{2}
\]

Isolando \(I_{\max}\):

\[
I_{\max} = V_{\max}\sqrt{\frac{C}{L}}
\]

Substituindo:

\[
I_{\max} = 3,0\sqrt{\frac{1,0 \times 10^{-9}}{3,0 \times 10^{-3}}}
\]

\[
I_{\max} = 3,0\sqrt{3,33 \times 10^{-7}}
\]

\[
I_{\max} = 3,0(5,77 \times 10^{-4})
\]

\[
I_{\max} \approx 1,73 \times 10^{-3}\,\text{A}
\]

\textbf{Resultado:}

\[
\boxed{I_{\max} \approx 1,7\,\text{mA}}
\]

\textbf{(c) Energia máxima armazenada no campo magnético do indutor}

A energia máxima do indutor é:

\[
U_{\max} = \frac{1}{2}LI_{\max}^{2}
\]

Mas é mais rápido usar a energia do capacitor:

\[
U_{\max} = \frac{1}{2}CV_{\max}^{2}
\]

Substituindo:

\[
U_{\max} = \frac{1}{2}(1,0 \times 10^{-9})(3,0)^2
\]

\[
U_{\max} = 0,5 \times 10^{-9} \times 9
\]

\[
U_{\max} = 4,5 \times 10^{-9}\,\text{J}
\]

\textbf{Resultado:}

\[
\boxed{U_{\max} = 4,5 \times 10^{-9}\,\text{J}}
\]

\subsection{Questão 14 — Mateus Valentim}
\section*{Resolução}

A frequência de ressonância de um circuito LC é dada por

\[
f=\frac{1}{2\pi\sqrt{LC}}.
\]

Deseja-se que a frequência varie linearmente com o ângulo de rotação \(\theta\), de modo que

\[
f(0^\circ)=2\times10^5\ \text{Hz}
\]

e

\[
f(180^\circ)=4\times10^5\ \text{Hz}.
\]

Como a dependência é linear, escrevemos

\[
f(\theta)=a\theta+b.
\]

O coeficiente angular é

\[
a=\frac{4\times10^5-2\times10^5}{180}
=\frac{2\times10^5}{180}.
\]

Como \(f(0)=2\times10^5\), temos

\[
b=2\times10^5.
\]

Portanto,

\[
f(\theta)=2\times10^5+\frac{2\times10^5}{180}\theta,
\]

ou, de forma equivalente,

\[
f(\theta)=2\times10^5\left(1+\frac{\theta}{180}\right).
\]

Agora isolamos a capacitância na expressão da frequência:

\[
f=\frac{1}{2\pi\sqrt{LC}}
\]

\[
f^2=\frac{1}{4\pi^2LC}
\]

\[
C=\frac{1}{4\pi^2Lf^2}.
\]

Substituindo \(L=1,0\,\text{mH}=10^{-3}\,\text{H}\), obtemos

\[
C(\theta)=
\frac{1}
{4\pi^2(10^{-3})[f(\theta)]^2}.
\]

Finalmente, substituindo a expressão de \(f(\theta)\),

\[
\boxed{
C(\theta)=
\frac{1}
{4\pi^2(10^{-3})
\left[
2\times10^5
\left(1+\frac{\theta}{180}\right)
\right]^2}
}.
\]

Para verificar os extremos:

\paragraph{Para \(\theta=0^\circ\):}

\[
f=2\times10^5\ \text{Hz}
\]

\[
C(0)=
\frac{1}
{4\pi^2(10^{-3})(2\times10^5)^2}
\]

\[
C(0)\approx6.33\times10^{-10}\,\text{F}
=633\,\text{pF}.
\]

\paragraph{Para \(\theta=180^\circ\):}

\[
f=4\times10^5\ \text{Hz}
\]

\[
C(180)=
\frac{1}
{4\pi^2(10^{-3})(4\times10^5)^2}
\]

\[
C(180)\approx1.58\times10^{-10}\,\text{F}
=158\,\text{pF}.
\]

Logo, a capacitância decresce de aproximadamente \(633\,\text{pF}\) para \(158\,\text{pF}\) à medida que o botão gira de \(0^\circ\) a \(180^\circ\).

\subsection{Questão 15 — Cicero Rodrigues}

\normalfont Um circuito LC oscilante formado por um capacitor de 1,0 nF e um indutor de 3,0 mH tem uma tensão máxima de 3,0 V. Determine \textbf{(a)} a carga máxima do capacitor, \textbf{(b)} a corrente máxima do circuito e \textbf{(c)} a energia máxima armazenada no campo magnético do indutor.

\textbf{Resolução:}
\vspace{0.3cm}

\textbf{Dados:}
\[C = 1,0 nF \Rightarrow 1,0 \times 10^{-9}\;F\]
\[L = 3,0 mH \Rightarrow 3,0 \times 10^{-3}\;H\]
\[V_{\max} = 3,0\;V\]

\vspace{0.2cm}
\textbf{(a) Carga máxima do capacitor}
A carga máxima é dada por:
\[Q_{\max} = C \cdot V_{\max}\]
Substituindo os valores:
\[Q_{\max} = (1,0 \times 10^{-9})(3,0)\]
\[Q_{\max} = 3,0 \times 10^{-9}\;C\]
\[\underline{Q_{\max} = 3,0\;nC}\]

\vspace{0.2cm}
\textbf{(b) Corrente máxima do circuito}
A energia máxima no capacitor é igual à energia máxima no indutor:
\[\frac{1}{2}CV_{\max}^{2} = \frac{1}{2}LI_{\max}^{2}\]
Isolando \(I_{\max}\):
\[I_{\max} = V_{\max}\sqrt{\frac{C}{L}}\]
Substituindo os valores:
\[I_{\max} = 3,0 \sqrt{\frac{1,0 \times 10^{-9}}{3,0 \times 10^{-3}}}\]
\[I_{\max} = 3,0 \sqrt{3,33 \times 10^{-7}}\]
\[I_{\max} \approx 1,73 \times 10^{-3}\;A\]
\[\underline{I_{\max} \approx 1,73\;mA}\]

\vspace{0.2cm}
\textbf{(c) Energia máxima armazenada no indutor}
A energia máxima é:
\[U_{\max} = \frac{1}{2}CV_{\max}^{2}\]
Substituindo os valores:
\[U_{\max} = \frac{1}{2}(1,0 \times 10^{-9})(3,0)^2\]
\[U_{\max} = 4,5 \times 10^{-9}\;J\]
\[\underline{U_{\max} = 4,5\;nJ}\]

\subsection{Questão 16 — Guilherme Araújo Floriano}

Para um circuito, a frequência angular natural é $\omega=(L_1C_1)^{-1/2}$. Para o segundo circuito, também $\omega=(L_2C_2)^{-1/2}$. Quando os dois circuitos são ligados em série, os elementos equivalentes ficam
\[
L_{\mathrm{eq}}=L_1+L_2,\qquad C_{\mathrm{eq}}=\frac{C_1C_2}{C_1+C_2}.
\]

Assim,
\[
\omega' = \left(L_{\mathrm{eq}}C_{\mathrm{eq}}\right)^{-1/2}
= \left[(L_1+L_2)\frac{C_1C_2}{C_1+C_2}\right]^{-1/2}
= \left[\frac{L_1C_1C_2 + L_2C_1C_2}{C_1+C_2}\right]^{-1/2}
\]

Simplificando, utilzando o fato de que $(L_1C_1)^{-1/2} = (L_2C_2)^{-1/2}$:

\[
\omega' = \left[\frac{L_1C_1C_2 + L_1C_1C_1}{C_1+C_2}\right]^{-1/2}
=  \left[{L_1C_1}\right]^{-1/2}*\left[\frac{C_1+C_2}{C_1+C_2}\right]^{-1/2}
\]

Como $\omega = {LC}^{-1/2}$ para cada circuito individual, o resultado simplifica para $\omega'=\omega$.

\subsection{Questão 17 — Arthur Roberto da Silva}
Esta questão aborda o comportamento de um circuito oscilante $LC$. Com base no termo fornecido para a fase da corrente, identificamos a frequência angular $\omega = 2500\text{ rad/s}$ e a amplitude da corrente $I = 1,60\text{ A}$ a partir dos dados do problema original associado.

\subsection*{(a) Determinação do instante de tempo $t$}

A expressão geral para a corrente variante no tempo é dada por $i(t) = I \sin(\omega t + \phi_0)$. Comparando com os dados do circuito onde a fase assume o valor máximo na condição considerada:
\begin{equation*}
\omega t + \phi_0 = 2500t + 0,680 = \frac{\pi}{2}
\end{equation*}

Isolando a variável de tempo $t$:
\begin{align*}
2500t &= \frac{\pi}{2} - 0,680 \\
2500t &\approx 1,5708 - 0,680 \\
2500t &= 0,8908 \\
t &= \frac{0,8908}{2500} \approx 3,56 \times 10^{-4}\text{ s}
\end{align*}

\subsection*{(b) Cálculo da Indutância ($L$)}

A frequência angular de ressonância em um circuito ideal $LC$ relaciona-se com os componentes através da expressão $\omega = (LC)^{-1/2} = \frac{1}{\sqrt{LC}}$. Sabendo que a capacitância do circuito é $C = 64,0\ \mu\text{F} = 64,0 \times 10^{-6}\text{ F}$:
\begin{equation*}
L = \frac{1}{\omega^2 C}
\end{equation*}

Substituindo os valores numéricos:
\begin{align*}
L &= \frac{1}{(2500\text{ rad/s})^2 \cdot (64,0 \times 10^{-6}\text{ F})} \\
L &= \frac{1}{(6,25 \times 10^6) \cdot (64,0 \times 10^{-6})} \\
L &= \frac{1}{400} = 2,50 \times 10^{-3}\text{ H} = 2,50\text{ mH}
\end{align*}

\subsection*{(c) Cálculo da Energia Total ($U$)}

A energia total armazenada no circuito em qualquer instante corresponde à energia máxima suportada pelos componentes, podendo ser calculada no momento em que toda a energia do sistema se encontra concentrada no indutor sob a forma de campo magnético (quando a corrente atinge sua amplitude máxima $I = 1,60\text{ A}$):
\begin{equation*}
U = \frac{1}{2} L I^2
\end{equation*}

Substituindo os valores obtidos nos itens anteriores:
\begin{align*}
U &= \frac{1}{2} \cdot (2,50 \times 10^{-3}\text{ H}) \cdot (1,60\text{ A})^2 \\
U &= 1,25 \times 10^{-3} \cdot 2,56 \\
U &= 3,20 \times 10^{-3}\text{ J} = 3,20\text{ mJ}
\end{align*}

\subsection*{Respostas Finais}
\begin{itemize}
    \item \textbf{(a)} $3,56 \times 10^{-4}\text{ s}$
    \item \textbf{(b)} $2,50 \times 10^{-3}\text{ H}$ \quad (ou $2,50\text{ mH}$)
    \item \textbf{(c)} $3,20 \times 10^{-3}\text{ J}$ \quad (ou $3,20\text{ mJ}$)
\end{itemize}

\subsection{Questão 18 — Arthur Roberto da Silva}
Esta questão aborda a distribuição de energia em um circuito oscilante $LC$ ideal em um determinado instante de tempo. O enunciado original estabelece que, nesse instante, $75,0\%$ da energia total do circuito está armazenada no campo magnético do indutor.

\subsection*{(a) Razão entre a carga instantânea e a carga máxima ($q/Q$)}

A energia total $U$ do circuito pode ser expressa em termos da carga máxima $Q$ no capacitor. Em um instante qualquer, a energia armazenada no campo elétrico do capacitor $U_E$ depende da carga instantânea $q$:
\begin{equation*}
U = \frac{Q^2}{2C} \quad \text{e} \quad U_E = \frac{q^2}{2C}
\end{equation*}

Se $75,0\%$ da energia total está no indutor, a porcentagem de energia restante contida no campo elétrico do capacitor é:
\begin{equation*}
100\% - 75,0\% = 25,0\%
\end{equation*}

Montando a razão entre a energia do campo elétrico e a energia total:
\begin{equation*}
\frac{U_E}{U} = \frac{\frac{q^2}{2C}}{\frac{Q^2}{2C}} = \left(\frac{q}{Q}\right)^2 = 0,250
\end{equation*}

Extraindo a raiz quadrada em ambos os lados para obter a razão das cargas:
\begin{equation*}
\frac{q}{Q} = \sqrt{0,250} = 0,500
\end{equation*}

\subsection*{(b) Razão entre a corrente instantânea e a corrente máxima ($i/I$)}

A energia total $U$ também pode ser expressa em termos da corrente máxima $I$ que atravessa o indutor. A energia magnética instantânea $U_B$ acumulada no indutor depende da corrente instantânea $i$:
\begin{equation*}
U = \frac{L I^2}{2} \quad \text{e} \quad U_B = \frac{L i^2}{2}
\end{equation*}

O enunciado afirma diretamente que a energia magnética instantânea corresponde a $75,0\%$ da energia total do sistema:
\begin{equation*}
\frac{U_B}{U} = \frac{\frac{L i^2}{2}}{\frac{L I^2}{2}} = \left(\frac{i}{I}\right)^2 = 0,750
\end{equation*}

Extraindo a raiz quadrada para obter a razão das correntes:
\begin{equation*}
\frac{i}{I} = \sqrt{0,750} \approx 0,866
\end{equation*}


\subsection{Questão 19 — Guilherme Araújo Floriano}

\textbf{(a)} A energia total é a soma das energias no indutor e no capacitor ($U = U_E + U_B$):
\[
U = \frac{q^2}{2C} + \frac{i^2L}{2}
= \frac{(3.80\times10^{-6}\,\mathrm{C})^2}{2(7.80\times10^{-6}\,\mathrm{F})}
 + \frac{(9.20\times10^{-3}\,\mathrm{A})^2(25.0\times10^{-3}\,\mathrm{H})}{2}
= 1.98\times10^{-6}\,\mathrm{J}.
\]

\textbf{(b)} No máximo carregamento do capacitor,
\[
U = \frac{Q^2}{2C}
\quad\Rightarrow\quad
Q = \sqrt{2CU}
= 5.56\times10^{-6}\,\mathrm{C}.
\]

\textbf{(c)} No máximo de corrente,
\[
U = \frac{I^2L}{2}
\quad\Rightarrow\quad
I = \sqrt{\frac{2U}{L}}
= 1.26\times10^{-2}\,\mathrm{A}.
\]

\textbf{(d)} Se $q(t=0)=q_0$, então $q_0=Q\cos\varphi$. Logo,
\[
\varphi = \cos^{-1}\!\left(\frac{q_0}{Q}\right)=\pm 46.9^\circ.
\]
Para que a carga esteja aumentando em $t=0$, escolhe-se $\varphi=-46.9^\circ$, já que, como $\frac{\mathrm{d}q}{\mathrm{d}t(0)} = -\omega Q \sin{\phi}$, o seno tem que ser negativo para tornar o resultado positivo

\textbf{(e)} Para que a derivada seja negativa e $\sin\varphi$ seja positiva, toma-se
\[
\varphi=+46.9^\circ.
\]

\subsection{Questão 20 — Cicero Rodrigues}


\textbf{20. }\normalfont Em um circuito LC oscilante no qual C = 4,00 $\mu$F, a diferença de potencial máxima entre os terminais do capacitor durante as oscilações é 1,50 V e a corrente máxima no indutor é 50,0 mA. Determine \textbf{(a)} a indutância L e \textbf{(b)} a frequência das oscilações. \textbf{(c)} Qual é o tempo necessário para que a carga do capacitor aumente de zero até o valor máximo?

\textbf{Resolução:}
\vspace{0.3cm}

\textbf{Dados:}
\[C = 4,00\,\mu F \Rightarrow 4,00 \times 10^{-6}\;F\]
\[V_{\max} = 1,50\;V\]
\[I_{\max} = 50,0\,mA \Rightarrow 50,0 \times 10^{-3}\;A\]

\vspace{0.2cm}
\textbf{(a) Indutância \(L\)}
A energia elétrica máxima no capacitor é igual à energia magnética máxima no indutor:
\[\frac{1}{2}CV_{\max}^{2} = \frac{1}{2}LI_{\max}^{2}\]
Isolando \(L\):
\[L = \frac{CV_{\max}^{2}}{I_{\max}^{2}}\]
Substituindo os valores:
\[L = \frac{(4,00 \times 10^{-6})(1,50)^2}{(50,0 \times 10^{-3})^2}\]
\[L = \frac{(4,00 \times 10^{-6})(2,25)}{2500 \times 10^{-6}}\]
\[L = \frac{9,00 \times 10^{-6}}{2500 \times 10^{-6}}\]
\[L = 0,0036\;H\]
\[\underline{L = 3,60\;mH}\]

\vspace{0.2cm}
\textbf{(b) Frequência das oscilações}
A frequência natural do circuito LC é dada por:
\[f = \frac{1}{2\pi \sqrt{LC}}\]
Substituindo os valores:
\[f = \frac{1}{2\pi \sqrt{(3,60 \times 10^{-3})(4,00 \times 10^{-6})}}\]
\[f = \frac{1}{2\pi \sqrt{14,4 \times 10^{-9}}}\]
\[f = \frac{1}{2\pi (1,20 \times 10^{-4})}\]
\[f = \frac{1}{2,4\pi \times 10^{-4}}\]
\[f \approx 1326\;Hz\]
\[\underline{f \approx 1,33\;kHz}\]

\vspace{0.2cm}
\textbf{(c) Tempo para a carga aumentar de zero ao máximo}
O tempo necessário para que o capacitor vá de totalmente descarregado (zero) até totalmente carregado (máximo) corresponde a um quarto do ciclo completo (\(T/4\)).
O período \(T\) é o inverso da frequência:
\[T = \frac{1}{f} = 2\pi (1,20 \times 10^{-4}) \approx 7,54 \times 10^{-4}\;s\]
Portanto, o tempo \(t\) é:
\[t = \frac{T}{4}\]
\[t = \frac{7,54 \times 10^{-4}}{4}\]
\[t \approx 1,88 \times 10^{-4}\;s\]
\[\underline{t \approx 0,188\;ms}\]

\subsection{Questão 21 — Mateus Valentim}
\paragraph{Resolução}

A frequência angular de um circuito \(LC\) é

\[
\omega=\frac{1}{\sqrt{LC}}.
\]

Substituindo os valores fornecidos,

\[
\omega
=
\frac{1}
{\sqrt{(3,00\times10^{-3})(2,70\times10^{-6})}}
=
1,11\times10^{4}\,\mathrm{rad/s}.
\]

Como \(q(0)=0\) e \(i(0)=2,00\,\mathrm{A}\), a carga e a corrente podem ser escritas como

\[
q(t)=Q_{\max}\sin(\omega t),
\]

\[
i(t)=\omega Q_{\max}\cos(\omega t).
\]

No instante \(t=0\),

\[
i(0)=I_{\max}=\omega Q_{\max},
\]

logo,

\[
Q_{\max}
=
\frac{I_{\max}}{\omega}
=
\frac{2,00}{1,11\times10^{4}}
=
1,80\times10^{-4}\,\mathrm{C}.
\]

Portanto,

\[
\boxed{Q_{\max}=1,80\times10^{-4}\,\mathrm{C}
=180\,\mu\mathrm{C}}.
\]

A energia armazenada no capacitor é

\[
U_C=\frac{q^2}{2C}.
\]

Assim,

\[
\frac{dU_C}{dt}
=
\frac{q}{C}\frac{dq}{dt}
=
\frac{qi}{C}.
\]

Substituindo as expressões de \(q(t)\) e \(i(t)\),

\[
\frac{dU_C}{dt}
=
\frac{\omega Q_{\max}^{2}}{C}
\sin(\omega t)\cos(\omega t).
\]

Utilizando a identidade trigonométrica

\[
\sin x\,\cos x
=
\frac{1}{2}\sin(2x),
\]

obtém-se

\[
\frac{dU_C}{dt}
=
\frac{\omega Q_{\max}^{2}}{2C}
\sin(2\omega t).
\]

O valor máximo ocorre quando

\[
\sin(2\omega t)=1,
\]

ou seja,

\[
2\omega t=\frac{\pi}{2}.
\]

Portanto,

\[
t
=
\frac{\pi}{4\omega}
=
\frac{\pi}{4(1,11\times10^{4})}
=
7,07\times10^{-5}\,\mathrm{s}.
\]

Assim,

\[
\boxed{
t=7,07\times10^{-5}\,\mathrm{s}
=
70,7\,\mu\mathrm{s}
}.
\]

Nesse instante, a taxa máxima de armazenamento de energia é

\[
\left(\frac{dU_C}{dt}\right)_{\max}
=
\frac{\omega Q_{\max}^{2}}{2C}.
\]

Substituindo os valores numéricos,

\[
\left(\frac{dU_C}{dt}\right)_{\max}
=
\frac{(1,11\times10^{4})
      (1,80\times10^{-4})^{2}}
           {2(2,70\times10^{-6})}
           =
           66,7\,\mathrm{W}.
           \]

           Portanto,

           \[
           \boxed{
           \left(\frac{dU_C}{dt}\right)_{\max}
           =
           66,7\,\mathrm{W}
           }.
           \]

           \[
           \boxed{
           \begin{aligned}
           (a)\;&Q_{\max}=1,80\times10^{-4}\,\mathrm{C}=180\,\mu\mathrm{C},\\[2mm]
           (b)\;&t=7,07\times10^{-5}\,\mathrm{s}=70,7\,\mu\mathrm{s},\\[2mm]
           (c)\;&\left(\frac{dU_C}{dt}\right)_{\max}=66,7\,\mathrm{W}.
           \end{aligned}
           }
           \]

\subsection{Questão 22 — Fábio Rodrigues}
\textbf{Dados}

\[
C = 4,00\,\mu F = 4,00 \times 10^{-6}\,F
\]

\[
V_{\max} = 1,50\,V
\]

\[
I_{\max} = 50,0\,mA = 5,00 \times 10^{-2}\,A
\]

\textbf{(a) Indutância \(L\)}

Em um circuito LC:

\[
\frac{1}{2}CV_{\max}^{2} = \frac{1}{2}LI_{\max}^{2}
\]

Isolando \(L\):

\[
L = \frac{CV_{\max}^{2}}{I_{\max}^{2}}
\]

Substituindo:

\[
L = \frac{(4,00 \times 10^{-6})(1,50)^2}{(5,00 \times 10^{-2})^2}
\]

\[
L = \frac{(4,00 \times 10^{-6})(2,25)}{2,50 \times 10^{-3}}
\]

\[
L = 3,60 \times 10^{-3}\,H
\]

\[
\boxed{L = 3,60\,mH}
\]

\textbf{(b) Frequência das oscilações}

A frequência de um circuito LC é:

\[
f = \frac{1}{2\pi\sqrt{LC}}
\]

Substituindo:

\[
f = \frac{1}{2\pi\sqrt{(3,60 \times 10^{-3})(4,00 \times 10^{-6})}}
\]

\[
f \approx 1,33 \times 10^{3}\,Hz
\]

\[
\boxed{f \approx 1,33\,kHz}
\]

\textbf{(c) Tempo para a carga ir de zero ao máximo}

Esse intervalo corresponde a:

\[
\frac{T}{4}
\]

Como:

\[
T = \frac{1}{f}
\]

Então:

\[
t = \frac{1}{4f}
\]

Substituindo:

\[
t = \frac{1}{4(1,33 \times 10^{3})}
\]

\[
t \approx 1,88 \times 10^{-4}\,s
\]

\[
\boxed{t \approx 188\,\mu s}
\]

\subsection{Questão 23 — Rodrigo Albuquerque}

Use a regra das malas para obter a equação diferencial de um circuito $LC$ (Eq. 31-11).  


Aplicando a regra das malhas no circuito LC:

\[
V_L+V_C=0
\]

A tensão no indutor é:

\[
V_L=L\frac{di}{dt}
\]

e a tensão no capacitor é:

\[
V_C=\frac{q}{C}
\]

Substituindo na equação da malha:

\[
L\frac{di}{dt}+\frac{q}{C}=0
\]

Como a corrente é a variação da carga no tempo:

\[
i=\frac{dq}{dt}
\]

então:

\[
\frac{di}{dt}=\frac{d^2q}{dt^2}
\]

Substituindo:

\[
L\frac{d^2q}{dt^2}+\frac{q}{C}=0
\]

Dividindo toda a equação por \(L\):

\[
\frac{d^2q}{dt^2}+\frac{1}{LC}q=0
\]

Essa é a equação diferencial do circuito LC.


\subsection{Questão 24 — Dávyla Maria}
Inserir resolução aqui.

\subsection{Questão 25 — Kaylane Castro}


\subsection*{Enunciado}

Que resistência $R$ deve ser ligada em série com uma indutância
$L = 220\,\text{mH}$ e uma capacitância
$C = 12,0\,\mu\text{F}$
para que a carga máxima do capacitor caia para $99,0\%$
do valor inicial após $50,0$ ciclos?
(Suponha que $\omega' \approx \omega$.)

\subsection*{Dados do problema}

\[
L = \num{220}\text{ mH} = \num{0,22}\text{ H}
\]

\[
C = \num{12,0}\,\mu\text{F} = \num{12,0}\times10^{-6}\text{ F}
\]

Após 50 ciclos, a carga máxima cai para:

\[
q_{\text{máx}} = 0{,}99Q
\]

\vspace{0.3cm}

Para um circuito $RLC$ fracamente amortecido, a amplitude da carga varia segundo:

\[
q_{\text{máx}}(t)=Qe^{-\frac{Rt}{2L}}
\]

Como o decaimento ocorre após 50 períodos:

\[
t=50T
\]

Substituindo:

\[
0{,}99Q = Qe^{-\frac{R(50T)}{2L}}
\]

Cancelando $Q$:

\[
0{,}99=e^{-\frac{25RT}{L}}
\]

Agora calculamos o período de oscilação.

Sabendo que:

\[
\omega \approx \frac{1}{\sqrt{LC}}
\]

e que:

\[
T=\frac{2\pi}{\omega}
\]

então:

\[
T=2\pi\sqrt{LC}
\]

Substituindo os valores:

\[
T=2\pi\sqrt{(0{,}22)(12,0\times10^{-6})}
\]

\[
T=2\pi\sqrt{2{,}64\times10^{-6}}
\]

\[
T\approx 0{,}01021\text{ s}
\]

Voltando para a equação do amortecimento:

\[
0{,}99=e^{-\frac{25RT}{L}}
\]

Aplicando logaritmo natural:

\[
\ln(0{,}99)=-\frac{25RT}{L}
\]

Isolando $R$:

\[
R=-\frac{L\ln(0{,}99)}{25T}
\]

Substituindo os valores:

\[
R=-\frac{(0{,}22)\ln(0{,}99)}{25(0{,}01021)}
\]

Como:

\[
\ln(0{,}99)\approx -0{,}01005
\]

temos:

\[
R=-\frac{(0{,}22)(-0{,}01005)}{0{,}25525}
\]

\[
R\approx 0{,}00866\,\Omega
\]

\subsection*{Resposta}

\[
\boxed{R \approx 8{,}66\times10^{-3}\,\Omega}
\]

ou

\[
\boxed{R \approx 8{,}66\text{ m}\Omega}
\]


\subsection{Questão 26 — Ana Alicy}

Um circuito de uma única malha é formado por um resistor de \SI{7,20}{\ohm}, um indutor de \SI{12,0}{\henry} e um capacitor de \SI{3,20}{\micro\farad}. Inicialmente, o capacitor possui uma carga de \SI{6,20}{\micro\coulomb} e a corrente é zero. Calcule a carga do capacitor após $N$ ciclos completos (a) para $N = 5$; (b) para $N = 10$; (c) para $N = 100$.
\vspace{0.8cm}

\textbf{Solução:}
\begin{quote}
\noindent \textbf{1. A Matemática do Decaimento}

\vspace{0.2cm}
\noindent A carga no capacitor em um circuito RLC varia com o tempo segundo a equação:
\[
q(t) = Q_0 e^{-Rt / 2L} \cos(\omega' t + \phi)
\]

\noindent O termo que interessa é o fator de amplitude (a parte antes do cosseno), pois ele dita qual será a carga máxima $Q$ ao final de um número $N$ de ciclos completos:
\[
Q_N = Q_0 e^{-Rt / 2L}
\]

\noindent Onde $t$ é o tempo total que se passou. Como precisa saber a carga após $N$ ciclos, o tempo total será o número de ciclos multiplicado pelo período de um ciclo ($T$):
\[
t = N \cdot T
\]

\vspace{0.4cm}
\noindent \textbf{2. Calculando o Período ($T$)}

\vspace{0.2cm}
\noindent O período é o tempo de um ciclo completo. Primeiro, precisa achar a frequência angular natural do circuito ($\omega_0$). 

\noindent Com $L = \SI{12,0}{\henry}$ e $C = \SI{3,20}{\micro\farad}$ (ou $3,20 \times 10^{-6}\ \text{F}$):
\[
\omega_0 = \frac{1}{\sqrt{LC}}
\]
\[
\omega_0 = \frac{1}{\sqrt{12,0 \cdot (3,20 \times 10^{-6})}}
\]
\[
\omega_0 = \frac{1}{\sqrt{38,4 \times 10^{-6}}}
\]
\[
\omega_0 \approx \SI{161,37}{\radian\per\second}
\]


\noindent Agora, calcula o período $T$:
\[
T = \frac{2\pi}{\omega_0}
\]
\[
T = \frac{2\pi}{161,37}
\]
\[
T \approx \SI{0,03894}{\second}
\]

\vspace{0.4cm}
\noindent \textbf{3. Calculando o Fator de Decaimento}

\vspace{0.2cm}
\noindent Agora vendo quanto a carga decai por ciclo. O expoente da nossa fórmula de amplitude para um tempo $t = N \cdot T$ fica:
\[
\text{Expoente} = -\frac{R}{2L} \cdot N \cdot T
\]
\[
\text{Expoente} = -\left(\frac{7,20}{2 \cdot 12,0}\right) \cdot N \cdot 0,03894
\]
\[
\text{Expoente} = -(0,300) \cdot N \cdot 0,03894
\]
\[
\text{Expoente} \approx -0,01168 \cdot N
\]

\noindent Então, a fórmula final prática para a carga máxima após $N$ ciclos é:
\[
Q_N = (\SI{6,20}{\micro\coulomb}) \cdot e^{-0,01168 \cdot N}
\]

\vspace{0.4cm}
\noindent \textbf{Respostas}

\vspace{0.2cm}
\noindent Agora é só substituir o valor de $N$ para cada alternativa da questão:

\vspace{0.2cm}
\noindent \textbf{(a) Para $N = 5$ ciclos:}
\[
Q_5 = 6,20 \cdot e^{-0,01168 \cdot 5}
\]
\[
Q_5 = 6,20 \cdot e^{-0,0584}
\]
\[
Q_5 = 6,20 \cdot 0,9433
\]

\vspace{0.2cm}
\noindent \textbf{(b) Para $N = 10$ ciclos:}
\[
Q_{10} = 6,20 \cdot e^{-0,01168 \cdot 10}
\]
\[
Q_{10} = 6,20 \cdot e^{-0,1168}
\]
\[
Q_{10} = 6,20 \cdot 0,8898
\]


\vspace{0.2cm}
\noindent \textbf{(c) Para $N = 100$ ciclos:}
\[
Q_{100} = 6,20 \cdot e^{-0,01168 \cdot 100}
\]
\[
Q_{100} = 6,20 \cdot e^{-1,168}
\]
\[
Q_{100} = 6,20 \cdot 0,3110
\]

\end{quote}
\vspace{0.2cm}


\subsection{Questão 32 — Ana Alicy}

Um gerador de corrente alternada com uma força eletromotriz $\mathscr{E} = \mathscr{E}_m \text{ sen } \omega_d t$, onde $\mathscr{E}_m = \SI{25,0}{\volt}$ e $\omega_d = \SI{377}{\radian\per\second}$, é ligado a um capacitor de \SI{4,15}{\micro\farad}. (a) Qual é o valor máximo da corrente? (b) Qual é a força eletromotriz do gerador no instante em que a corrente é máxima? (c) Qual é a corrente quando a força eletromotriz é \SI{-12,5}{\volt} e está aumentando em valor absoluto?
\vspace{0.8cm}

\textbf{Solução:}
\begin{quote}

\noindent \textbf{Organizando os dados fornecidos:}
\begin{itemize}
    \item Força eletromotriz (tensão): $\mathscr{E}(t) = 25,0 \text{ sen}(377t)\ \text{V}$
    \item Amplitude da tensão ($\mathscr{E}_m$): \SI{25,0}{\volt}
    \item Frequência angular ($\omega_d$): \SI{377}{\radian\per\second}
    \item Capacitância ($C$): \SI{4,15}{\micro\farad} = $4,15 \times 10^{-6}\ \text{F}$
\end{itemize}

\vspace{0.4cm}
\noindent \textbf{(a) Qual é o valor máximo da corrente?}

\vspace{0.2cm}
\noindent A corrente máxima ($I_m$) é a amplitude da tensão dividida pela reatância capacitiva ($X_C$):
\[
I_m = \frac{\mathscr{E}_m}{X_C}
\]

\noindent Sabendo que a reatância capacitiva é $X_C = \frac{1}{\omega_d C}$. Substituindo isso na equação:
\[
I_m = \mathscr{E}_m \cdot \omega_d \cdot C
\]

\noindent Agora, inserindo os valores:
\[
I_m = (25,0) \cdot (377) \cdot (4,15 \times 10^{-6})
\]
\[
I_m = 39113,75 \times 10^{-6}\ \text{A}
\]
\[
I_m \approx \SI{0,0391}{\ampere}
\]


\vspace{0.4cm}
\noindent \textbf{(b) Qual é a força eletromotriz quando a corrente é máxima?}

\vspace{0.2cm}
\noindent resolvendo  conceitualmente e matematicamente. A equação da tensão é uma função seno: $\mathscr{E}(t) = \mathscr{E}_m \text{ sen}(\omega_d t)$.

\noindent Como a corrente no capacitor é a derivada da carga (que depende da tensão), ela é uma função cosseno: $i(t) = I_m \cos(\omega_d t)$.

\noindent Quando a corrente atinge o seu valor máximo, sabemos que $\cos(\omega_d t) = \pm 1$. Para que o cosseno seja $\pm 1$, o ângulo $\omega_d t$ deve ser $0, \pi, 2\pi$, etc. E nesses exatos ângulos, a função seno ($\text{sen}$) vale zero.

\noindent (b) A força eletromotriz é \SI{0}{\volt}. (Sempre que a corrente atinge o pico em um capacitor ideal, a tensão cruza o zero, e vice-versa).
\end{quote}

\vspace{0.2cm}

\subsection*{Questão 33 — Kaylane Castro}


{Enunciado}

Um gerador de corrente alternada tem uma força eletromotriz

\[
\mathcal{E}(t)=\mathcal{E}_m\sin(\omega_d t-\pi/4),
\]

onde $\mathcal{E}_m = 30,0\,\text{V}$ e
$\omega_d = 350\,\text{rad/s}$.

A corrente produzida no circuito ao qual o gerador está ligado é

\[
i(t)=I\sin(\omega_d t-3\pi/4),
\]

onde $I=620\,\text{mA}$.

Determine:

\begin{itemize}
    \item[(a)] o instante em que a força eletromotriz atinge pela primeira vez o valor máximo;
    \item[(b)] o instante em que a corrente atinge pela primeira vez o valor máximo;
    \item[(c)] se o circuito contém um capacitor, um indutor ou um resistor;
    \item[(d)] o valor da capacitância, indutância ou resistência.
\end{itemize}

\subsection*{Dados}

\[
\mathcal{E}(t)=\mathcal{E}_m\sin(\omega_d t-\pi/4)
\]

\[
\mathcal{E}_m = 30,0\,\text{V}
\]

\[
\omega_d = 350\,\text{rad/s}
\]

\[
i(t)=I\sin(\omega_d t-3\pi/4)
\]

\[
I = 620\,\text{mA} = 0,620\,\text{A}
\]


\subsection*{(a) Instante do primeiro valor máximo da fem}

A função seno atinge valor máximo quando:

\[
\sin(\theta)=1
\]

ou seja:

\[
\theta = \frac{\pi}{2}
\]

Então:

\[
\omega_d t - \frac{\pi}{4} = \frac{\pi}{2}
\]

\[
\omega_d t = \frac{\pi}{2}+\frac{\pi}{4}
\]

\[
\omega_d t = \frac{3\pi}{4}
\]

Isolando $t$:

\[
t=\frac{3\pi}{4\omega_d}
\]

Substituindo $\omega_d = 350\,\text{rad/s}$:

\[
t=\frac{3\pi}{4(350)}
\]

\[
t=\frac{3\pi}{1400}
\]

\[
t\approx 6,73\times10^{-3}\,\text{s}
\]

\[
\boxed{t_{\mathcal{E}} \approx 6,73\,\text{ms}}
\]



\subsection*{(b) Instante do primeiro valor máximo da corrente}

Para a corrente:

\[
\omega_d t - \frac{3\pi}{4} = \frac{\pi}{2}
\]

\[
\omega_d t = \frac{\pi}{2}+\frac{3\pi}{4}
\]

\[
\omega_d t = \frac{5\pi}{4}
\]

Isolando $t$:

\[
t=\frac{5\pi}{4\omega_d}
\]

Substituindo os valores:

\[
t=\frac{5\pi}{4(350)}
\]

\[
t=\frac{5\pi}{1400}
\]

\[
t\approx 11,2\times10^{-3}\,\text{s}
\]

\[
\boxed{t_i \approx 11,2\,\text{ms}}
\]


\subsection*{(c) Identificação do elemento}

As fases são:

\[
\phi_{\mathcal{E}}=-\frac{\pi}{4}
\]

\[
\phi_i=-\frac{3\pi}{4}
\]

O ângulo de fase do circuito é:

\[
\phi=\phi_{\mathcal{E}}-\phi_i
\]

\[
\phi=
-\frac{\pi}{4}
-
\left(
-\frac{3\pi}{4}
\right)
\]

\[
\phi=\frac{\pi}{2}
\]

Isso significa que a tensão adianta a corrente em $90^\circ$.

Esse comportamento é característico de um circuito puramente indutivo.

Portanto, o elemento é um:

\[
\boxed{\text{Indutor}}
\]


\subsection*{(d) Cálculo da indutância}

Para um indutor:

\[
X_L=\omega_d L
\]

e também:

\[
X_L=\frac{\mathcal{E}_m}{I}
\]

Então:

\[
\omega_d L = \frac{\mathcal{E}_m}{I}
\]

Isolando $L$:

\[
L=\frac{\mathcal{E}_m}{I\omega_d}
\]

Substituindo os valores:

\[
L=
\frac{30,0}
{(0,620)(350)}
\]

\[
L=
\frac{30,0}{217}
\]

\[
L\approx 0,138\,\text{H}
\]

\[
\boxed{L\approx 138\,\text{mH}}
\]



\subsection{Questão 34 — Dávyla Maria}
Inserir resolução aqui.

\subsection{Questão 44 — Rodrigo Albuquerque}





\subsection{Questão 45 — Fábio Rodrigues}
\textbf{Dados}

\[
R = 5,00\,\Omega
\]

\[
C = 20,0\,\mu F = 2,00 \times 10^{-5}\,F
\]

\[
L = 1,00\,H
\]

\[
\varepsilon_m = 30,0\,V
\]

\textbf{(a) Frequência angular de ressonância}

Para um circuito RLC:

\[
\omega_0 = \frac{1}{\sqrt{LC}}
\]

Substituindo:

\[
\omega_0 = \frac{1}{\sqrt{(1,00)(2,00 \times 10^{-5})}}
\]

\[
\omega_0 = \frac{1}{4,47 \times 10^{-3}}
\]

\[
\omega_0 \approx 223,6\,rad/s
\]

\[
\boxed{\omega_0 \approx 224\,rad/s}
\]

\textbf{(b) Corrente máxima}

Na ressonância:

\[
I_{\max} = \frac{\varepsilon_m}{R}
\]

Substituindo:

\[
I_{\max} = \frac{30,0}{5,00}
\]

\[
I_{\max} = 6,0\,A
\]

\[
\boxed{I_{\max} = 6,0\,A}
\]

\textbf{(c) Frequência angular para metade da corrente máxima (\(\omega_{d1} < \omega_0\))}

Quando a corrente vale metade da máxima:

\[
I = \frac{I_{\max}}{2}
\]

Na impedância:

\[
I = \frac{\varepsilon_m}{Z}
\]

Logo:

\[
Z = 2R
\]

Sabendo que:

\[
Z = \sqrt{R^2 + \left(\omega L - \frac{1}{\omega C}\right)^2}
\]

Então:

\[
(2R)^2 = R^2 + \left(\omega L - \frac{1}{\omega C}\right)^2
\]

\[
4R^2 - R^2 = \left(\omega L - \frac{1}{\omega C}\right)^2
\]

\[
3R^2 = \left(\omega L - \frac{1}{\omega C}\right)^2
\]

Para \(\omega_{d1}\):

\[
\omega L - \frac{1}{\omega C} = -\sqrt{3}R
\]

Substituindo os valores:

\[
\omega - \frac{50000}{\omega} = -8,66
\]

Resolvendo:

\[
\omega_{d1} \approx 219\,rad/s
\]

\[
\boxed{\omega_{d1} \approx 219\,rad/s}
\]

\textbf{(d) Frequência angular para metade da corrente máxima (\(\omega_{d2} > \omega_0\))}

Agora:

\[
\omega L - \frac{1}{\omega C} = +\sqrt{3}R
\]

\[
\omega - \frac{50000}{\omega} = +8,66
\]

Resolvendo:

\[
\omega_{d2} \approx 228\,rad/s
\]

\[
\boxed{\omega_{d2} \approx 228\,rad/s}
\]

\textbf{(e) Largura relativa da curva de ressonância}

\[
\frac{\omega_{d2} - \omega_{d1}}{\omega_0}
\]

Substituindo:

\[
\frac{228 - 219}{224}
\]

\[
\approx 0,040
\]

\[
\boxed{0,040}
\]
\subsection{Questão 47 — Mateus Valentim}

\begin{enumerate}
    \item[(a)] Sim. Em um circuito \(RLC\) série, a tensão no indutor pode ser
        maior que a força eletromotriz do gerador. Isso ocorre porque, perto da
            ressonância, a corrente pode ser grande e a tensão no indutor é

                \[
                    V_L=I X_L=I\omega L.
                        \]

                            As tensões no indutor e no capacitor ficam defasadas de \(180^\circ\)
                                entre si e podem se cancelar parcialmente na soma fasorial, mesmo que cada
                                    uma delas seja grande.

                                        \item[(b)] Na ressonância,

                                            \[
                                                \omega_0=\frac{1}{\sqrt{LC}}.
                                                    \]

                                                        Assim,

                                                            \[
                                                                \omega_0=
                                                                    \frac{1}{\sqrt{(1{,}0)(1{,}0\times 10^{-6})}}
                                                                        =1{,}0\times 10^3\,\text{rad/s}.
                                                                            \]

                                                                                Também na ressonância a impedância do circuito é apenas \(R\), de modo que

                                                                                    \[
                                                                                        I_m=\frac{\mathcal{E}_m}{R}
                                                                                            =\frac{10}{10}=1{,}0\,\text{A}.
                                                                                                \]

                                                                                                    Portanto,

                                                                                                        \[
                                                                                                            V_{L,m}=I_m\omega_0 L
                                                                                                                =(1{,}0)(1{,}0\times 10^3)(1{,}0)
                                                                                                                    =1{,}0\times 10^3\,\text{V}.
                                                                                                                        \]
                                                                                                                        \end{enumerate}

\subsection{Questão 48 — Cicero Rodrigues}

\normalfont A Fig. 31.14 mostra um circuito RLC, alimentado por um gerador, que possui dois capacitores iguais e duas chaves. A amplitude da força eletromotriz é 12,0 V e a frequência do gerador é 60,0 Hz. Com as duas chaves abertas, a corrente está adiantada 30,9$^\circ$ em relação à tensão. Com a chave S1 fechada e a chave S2 aberta, a corrente está adiantada 15,0$^\circ$ em relação à tensão. Com as duas chaves fechadas, a amplitude da corrente é 447 mA. Determine o valor \textbf{(a)} de R, \textbf{(b)} de C e \textbf{(c)} de L.

\includegraphics[width=0.5\textwidth]{image_10.png}

\textbf{Resolução:}
\vspace{0.3cm}

\textbf{Dados:}
\[\mathcal{E}_m = 12,0\;V\]
\[f = 60,0\;Hz \Rightarrow \omega = 2\pi f \approx 377\;rad/s\]
\[I_m = 447\;mA \Rightarrow 0,447\;A\]
\[\phi_1 = -30,9^\circ \quad ; \quad \phi_2 = 15,0^\circ \quad\]

\vspace{0.2cm}
\textbf{Análise das Impedâncias:}
Com as duas chaves fechadas, S2 curto-circuita o resistor \(R\) e S1 curto-circuita um dos capacitores. A impedância restante é:
\[|X_L - X_C| = \frac{\mathcal{E}_m}{I_m} = \frac{12,0}{0,447} \approx 26,85\;\Omega\]
Como assumimos a tensão adiantada para S1 fechada, \(X_L > X_C\), logo \(X_L - X_C = 26,85\;\Omega\).

Com S1 fechada e S2 aberta, o resistor \(R\) volta à malha (junto com \(L\) e um \(C\)):
\[\tan(15,0^\circ) = \frac{X_L - X_C}{R} \Rightarrow X_L - X_C = R \tan(15,0^\circ)\]

Com ambas abertas, os dois capacitores ficam em série (\(C_{eq} = C/2 \Rightarrow X_{eq} = 2X_C\)):
\[\tan(-30,9^\circ) = \frac{X_L - 2X_C}{R} \Rightarrow X_L - 2X_C = -0,5985 \cdot R\]

\vspace{0.2cm}
\textbf{(a) Valor de \(R\)}
Substituindo o valor da impedância na equação da fase:
\[R = \frac{X_L - X_C}{\tan(15,0^\circ)}\]
\[R = \frac{26,85}{0,2679}\]
\[R \approx 100,2\;\Omega\]
\[\underline{R \approx 100\;\Omega}\]

\vspace{0.2cm}
\textbf{(b) Valor de \(C\)}
Temos agora um sistema de duas equações:
\[X_L - X_C = 26,85\]
\[X_L - 2X_C = -0,5985(100) = -59,85\]
Subtraindo a segunda equação da primeira:
\[X_C = 26,85 - (-59,85)\]
\[X_C = 86,70\;\Omega\]
Calculando a capacitância:
\[C = \frac{1}{\omega X_C}\]
\[C = \frac{1}{(377)(86,70)}\]
\[C \approx 30,6 \times 10^{-6}\;F\]
\[\underline{C = 30,6\;\mu F}\]

\vspace{0.2cm}
\textbf{(c) Valor de \(L\)}
Usando a primeira equação do sistema para achar \(X_L\):
\[X_L = 26,85 + X_C\]
\[X_L = 26,85 + 86,70\]
\[X_L = 113,55\;\Omega\]
Calculando a indutância:
\[L = \frac{X_L}{\omega}\]
\[L = \frac{113,55}{377}\]
\[L \approx 0,301\;H\]
\[\underline{L = 301\;mH}\]

\subsection{Questão 49 — Guilherme Araújo Floriano}

Usando as expressões do Problema 31-45,
\[
\omega_1 = \frac{-3CR + \sqrt{3C^2R^2+4LC}}{2LC},
\qquad
\omega_2 = \frac{3CR + \sqrt{3C^2R^2+4LC}}{2LC}.
\]
Com a Eq.~31-4,
\[
\frac{\Delta\omega}{\omega} = 
\frac{\omega_2-\omega_1}{\omega} =
\frac{2\sqrt{3}CR\sqrt{LC}}{2LC} =
R\sqrt{\frac{3C}{L}}
\]

\subsection{Questão 50 — Arthur Roberto da Silva}
\subsection*{(a) Cálculo da Reatância Capacitiva ($X_C$)}

A reatância capacitiva representa a oposição que um capacitor oferece à passagem da corrente alternada e é inversamente proporcional à frequência e à capacitância:
\begin{equation*}
X_C = \frac{1}{2\pi f C}
\end{equation*}

Substituindo os valores fornecidos no Sistema Internacional ($f = 400\text{ Hz}$ e $C = 24,0 \times 10^{-6}\text{ F}$):
\begin{equation*}
X_C = \frac{1}{2\pi (400\text{ Hz})(24,0 \times 10^{-6}\text{ F})} \approx 16,6\ \Omega
\end{equation*}

\subsection*{(b) Cálculo da Impedância do Circuito ($Z$)}

A impedância $Z$ é a oposição total do circuito RLC à corrente alternada, combinando a resistência ($R$), a reatância indutiva ($X_L = 2\pi f L$) e a reatância capacitiva ($X_C$):
\begin{equation*}
Z = \sqrt{R^2 + (X_L - X_C)^2} = \sqrt{R^2 + (2\pi f L - X_C)^2}
\end{equation*}

Substituindo os valores ($L = 150 \times 10^{-3}\text{ H}$):
\begin{align*}
Z &= \sqrt{(220\ \Omega)^2 + [2\pi(400\text{ Hz})(150 \times 10^{-3}\text{ H}) - 16,6\ \Omega]^2} \\
Z &= \sqrt{(220)^2 + [376,99 - 16,6]^2} \\
Z &= \sqrt{48400 + (360,39)^2} \approx 422\ \Omega
\end{align*}

\subsection*{(c) Cálculo da Amplitude da Corrente ($I$)}

Utilizando a variação da Lei de Ohm para circuitos de corrente alternada, a amplitude da corrente ($I$) é a razão entre a amplitude da fem ($\mathcal{E}_m$) e a impedância total ($Z$):
\begin{equation*}
I = \frac{\mathcal{E}_m}{Z} = \frac{220\text{ V}}{422\ \Omega} \approx 0,521\text{ A}
\end{equation*}

\subsection*{(d) Análise da Reatância Capacitiva com a Capacitância Equivalente ($C_{\text{eq}}$)}

A relação entre a reatância capacitiva e a capacitância é dada por $X_C \propto C_{\text{eq}}^{-1}$. Isto significa que a reatância e a capacitância são inversamente proporcionais. Portanto, se a capacitância equivalente ($C_{\text{eq}}$) do circuito diminuir, a reatância capacitiva ($X_C$) irá \textbf{aumentar}.

\subsection*{(e) Nova Impedância com Capacitores em Série}

Ao adicionar um segundo capacitor idêntico em série com o primeiro, a nova capacitância equivalente ($C_{\text{eq}}$) cai para metade do valor original:
\begin{equation*}
C_{\text{eq}} = \frac{C}{2}
\end{equation*}

Como $X_C$ é inversamente proporcional a $C$, reduzir a capacitância para a metade duplica o valor da reatância capacitiva original ($2 \times 16,6\ \Omega$). Substituindo este novo termo na equação da impedância:
\begin{align*}
Z &= \sqrt{(220\ \Omega)^2 + [2\pi(400\text{ Hz})(150 \times 10^{-3}\text{ H}) - 2(16,6\ \Omega)]^2} \\
Z &= \sqrt{(220)^2 + [376,99 - 33,2]^2} \\
Z &= \sqrt{48400 + (343,79)^2} \approx 408\ \Omega
\end{align*}

Comparando os resultados, a nova impedância é menor que a anterior ($408\ \Omega < 422\ \Omega$). Portanto, a impedância total do circuito \textbf{diminui}.

\subsection*{(f) Impacto na Amplitude da Corrente}

A relação entre a corrente e a impedância é inversamente proporcional, isto é, $I \propto Z^{-1}$. Como demonstrado no passo anterior que a impedância $Z$ diminuiu, a amplitude da corrente consequentemente \textbf{aumenta}.

\subsection{Questão 51 — Arthur Roberto da Silva}
Este problema analisa a frequência de ressonância de um circuito contendo múltiplos indutores e capacitores. Note que, embora a solução original utilize o símbolo $\omega$, a unidade final expressa em Hertz ($\text{Hz}$) indica que a grandeza calculada é, na verdade, a frequência linear $f$. A equação clássica para a frequência de ressonância em circuitos lineares é:
\begin{equation*}
f = \frac{1}{2\pi\sqrt{L_{\text{eq}}C_{\text{eq}}}}
\end{equation*}

\subsection*{(a) Cálculo da Frequência de Ressonância ($f$)}

O circuito possui indutores associados em série, o que significa que a indutância equivalente ($L_{\text{eq}}$) é a soma direta das indutâncias individuais:
\begin{equation*}
L_{\text{eq}} = L_1 + L_2
\end{equation*}

Os capacitores estão associados em paralelo, logo, a capacitância equivalente ($C_{\text{eq}}$) é a soma direta das capacitâncias individuais:
\begin{equation*}
C_{\text{eq}} = C_1 + C_2 + C_3
\end{equation*}

Substituindo as expressões equivalentes na fórmula da frequência de ressonância, obtemos:
\begin{equation*}
f = \frac{1}{2\pi\sqrt{(L_1 + L_2)(C_1 + C_2 + C_3)}}
\end{equation*}

Substituindo os valores numéricos fornecidos no Sistema Internacional:
\begin{itemize}
    \item $L_1 = 1,70\text{ mH} = 1,70 \times 10^{-3}\text{ H}$
    \item $L_2 = 2,30\text{ mH} = 2,30 \times 10^{-3}\text{ H}$
    \item $C_1 = 4,00\ \mu\text{F} = 4,00 \times 10^{-6}\text{ F}$
    \item $C_2 = 2,50\ \mu\text{F} = 2,50 \times 10^{-6}\text{ F}$
    \item $C_3 = 3,50\ \mu\text{F} = 3,50 \times 10^{-6}\text{ F}$
\end{itemize}

\begin{equation*}
f = \frac{1}{2\pi\sqrt{(1,70 \times 10^{-3}\text{ H} + 2,30 \times 10^{-3}\text{ H})(4,00 \times 10^{-6}\text{ F} + 2,50 \times 10^{-6}\text{ F} + 3,50 \times 10^{-6}\text{ F})}}
\end{equation*}

Simplificando os termos dentro da raiz:
\begin{align*}
L_{\text{eq}} &= 4,00 \times 10^{-3}\text{ H} \\
C_{\text{eq}} &= 10,00 \times 10^{-6}\text{ F}
\end{align*}

\begin{align*}
f &= \frac{1}{2\pi\sqrt{(4,00 \times 10^{-3}) \cdot (10,00 \times 10^{-6})}} \\
f &= \frac{1}{2\pi\sqrt{4,00 \times 10^{-8}}} \\
f &= \frac{1}{2\pi \cdot (2,00 \times 10^{-4})} \\
f &= \frac{1}{4\pi \times 10^{-4}} \approx 796\text{ Hz}
\end{align*}

\subsection*{(b) Efeito da Resistência ($R$) na Frequência de Ressonância}

A frequência de ressonância de um circuito depende estritamente dos componentes que armazenam energia, ou seja, dos indutores ($L$) e dos capacitores ($C$). Como a resistência elétrica $R$ do circuito não aparece na equação fundamental da frequência de ressonância, concluímos que \textbf{a frequência de ressonância não se alterará com o aumento de $R$}.

\subsection*{(c) Efeito do Aumento de $L_1$ na Frequência de Ressonância}

A relação matemática entre a frequência e a indutância total é dada por $f \propto (L_1 + L_2)^{-1/2}$. Como a frequência é inversamente proporcional à raiz quadrada da indutância, um aumento no valor do indutor $L_1$ provocará um aumento na indutância equivalente ($L_{\text{eq}}$). Consequentemente, a frequência de ressonância \textbf{irá diminuir}.

\subsection*{(d) Efeito da Remoção de $C_3$ na Frequência de Ressonância}

A relação entre a frequência e a capacitância total obedece à proporcionalidade $f \propto C_{\text{eq}}^{-1/2}$. Como os capacitores estão em paralelo, a remoção do capacitor $C_3$ fará com que a capacitância equivalente do sistema diminua ($C_{\text{eq, novo}} = C_1 + C_2$). Sendo a frequência inversamente proporcional à capacitância, essa redução em $C_{\text{eq}}$ fará com que a frequência de ressonância \textbf{aumente}.

\subsection{Questão 56 — Guilherme Araújo Floriano}

Como o circuito não contém reatâncias,
\[
\varepsilon_{\mathrm{rms}} = I_{\mathrm{rms}}(r+R).
\]
Pela Eq.~31-71, a potência média dissipada em $R$ é
\[
P_R = I_{\mathrm{rms}}^2R = \frac{\varepsilon_{\mathrm{rms}}^2R}{(r+R)^2}.
\]
Maximizando em relação a $R$,
\[
\frac{dP_R}{dR}=0 \quad\Rightarrow\quad R=r.
\]

\subsection{Questão 57 — Cicero Rodrigues}
 
\textbf{57. }\normalfont Em um circuito RLC como o da Fig. 31.3.2, suponha que \(R = 5,00\,\Omega\), \(L = 60,0\;mH\), \(f_d = 60,0\;Hz\) e \(\mathcal{E}_m = 30,0\;V\). \textbf{(a)} Para qual valor de capacitância a potência dissipada na resistência é máxima? \textbf{(b)} Para qual valor de capacitância a potência dissipada na resistência é mínima? Determine \textbf{(c)} a dissipação máxima, \textbf{(d)} o ângulo de fase correspondente e \textbf{(e)} o fator de potência correspondente. Determine também \textbf{(f)} a dissipação mínima, \textbf{(g)} o ângulo de fase correspondente e \textbf{(h)} o fator de potência correspondente.

\textbf{Resolução:}
\vspace{0.3cm}

\textbf{Dados:}
\[R = 5,00\;\Omega\]
\[L = 60,0\;mH = 60,0 \times 10^{-3}\;H\]
\[f_d = 60,0\;Hz \Rightarrow \omega_d = 2\pi(60,0) \approx 377\;rad/s\]
\[\mathcal{E}_m = 30,0\;V\]

\vspace{0.2cm}
\textbf{(a) Capacitância para potência máxima}
A potência dissipada é máxima na ressonância (\(X_L = X_C\)):
\[\omega_d L = \frac{1}{\omega_d C} \Rightarrow C = \frac{1}{\omega_d^2 L}\]
\[C = \frac{1}{(377)^2 (60,0 \times 10^{-3})}\]
\[C \approx 1,17 \times 10^{-4}\;F\]
\[\underline{C = 117\;\mu F}\]

\vspace{0.2cm}
\textbf{(b) Capacitância para potência mínima}
A potência é mínima quando a impedância do circuito tende ao infinito. Como \(R\), \(L\) e \(\omega_d\) são fixos, isso só ocorre se a reatância capacitiva for infinita (\(X_C \to \infty\)):
\[X_C = \frac{1}{\omega_d C} \to \infty \Rightarrow C \to 0\]
\[\underline{C \to 0}\]

\vspace{0.2cm}
\textbf{(c), (d) e (e) Parâmetros na dissipação máxima}
Na ressonância, a reatância total é nula e a impedância é puramente resistiva (\(Z = R\)).
\textbf{(c)} Dissipação máxima (\(P_{\max}\)):
\[P_{\max} = \frac{I_m^2 R}{2} = \frac{1}{2} \left(\frac{\mathcal{E}_m}{R}\right)^2 R = \frac{\mathcal{E}_m^2}{2R}\]
\[P_{\max} = \frac{(30,0)^2}{2(5,00)} = \frac{900}{10,0}\]
\[\underline{P_{\max} = 90,0\;W}\]
\textbf{(d)} Ângulo de fase (\(\phi\)):
\[\tan\phi = \frac{X_L - X_C}{R} = 0 \Rightarrow \underline{\phi = 0^\circ}\]
\textbf{(e)} Fator de potência (\(\cos\phi\)):
\[\underline{\cos\phi = 1}\]

\vspace{0.2cm}
\textbf{(f), (g) e (h) Parâmetros na dissipação mínima}
Para \(C \to 0\), a impedância \(Z \to \infty\) e a corrente \(I \to 0\).
\textbf{(f)} Dissipação mínima (\(P_{\min}\)):
\[\underline{P_{\min} = 0\;W}\]
\textbf{(g)} Ângulo de fase (\(\phi\)):
Como \(X_C \to \infty\), o termo capacitivo domina completamente a equação, tornando o circuito puramente capacitivo:
\[\tan\phi \to -\infty \Rightarrow \underline{\phi = -90,0^\circ}\]
\textbf{(h)} Fator de potência (\(\cos\phi\)):
\[\cos(-90,0^\circ) = 0 \Rightarrow \underline{\cos\phi = 0}\]

\subsection{Questão 58 — Mateus Valentim}
\subsection*{Resposta}

Dados:

\[
R = 16,0\,\Omega,
\qquad
L = 9,20\times10^{-3}\,\text{H},
\qquad
C = 31,2\times10^{-6}\,\text{F}
\]

\[
\mathcal{E}(t)=\mathcal{E}_m\sin(\omega t),
\qquad
\mathcal{E}_m=45,0\,\text{V},
\qquad
\omega=3000\,\text{rad/s}
\]

\[
t=0,442\,\text{ms}
=0,442\times10^{-3}\,\text{s}
\]

As reatâncias são

\[
X_L=\omega L
=(3000)(9,20\times10^{-3})
=27,6\,\Omega
\]

\[
X_C=\frac{1}{\omega C}
=
\frac{1}{(3000)(31,2\times10^{-6})}
=10,68\,\Omega
\]

A impedância do circuito é

\[
Z=\sqrt{R^2+(X_L-X_C)^2}
\]

\[
Z=\sqrt{16^2+(27,6-10,68)^2}
\]

\[
Z=23,29\,\Omega
\]

O ângulo de fase vale

\[
\tan\phi=
\frac{X_L-X_C}{R}
=
\frac{16,92}{16}
\]

\[
\phi=0,813\,\text{rad}
\]

A corrente máxima é

\[
I_m=\frac{\mathcal{E}_m}{Z}
=
\frac{45}{23,29}
=
1,93\,\text{A}
\]

Calculando

\[
\omega t
=
(3000)(0,442\times10^{-3})
=
1,326\,\text{rad}
\]

a corrente instantânea é

\[
i(t)
=
I_m\sin(\omega t-\phi)
\]

\[
i
=
1,93\sin(1,326-0,813)
=
0,947\,\text{A}
\]

e a tensão do gerador no instante considerado é

\[
\mathcal{E}
=
45\sin(1,326)
=
43,67\,\text{V}
\]

\subsubsection*{(a) Potência fornecida pelo gerador}

\[
P_G
=
\mathcal{E}\,i
\]

\[
P_G
=
(43,67)(0,947)
\]

\[
\boxed{P_G=41,4\,\text{W}}
\]

\subsubsection*{(b) Taxa de variação da energia no capacitor}

A amplitude da tensão no capacitor é

\[
V_{Cm}=I_mX_C
=(1,93)(10,68)
=20,64\,\text{V}
\]

Logo,

\[
v_C
=
20,64
\sin\!\left(
1,326-0,813-\frac{\pi}{2}
\right)
=
-17,98\,\text{V}
\]

Portanto,

\[
P_C=v_C\,i
\]

\[
P_C
=
(-17,98)(0,947)
\]

\[
\boxed{P_C=-17,0\,\text{W}}
\]

\subsubsection*{(c) Taxa de variação da energia no indutor}

A amplitude da tensão no indutor é

\[
V_{Lm}
=
I_mX_L
=
(1,93)(27,6)
=
53,3\,\text{V}
\]

Assim,

\[
v_L
=
53,3
\sin\!\left(
1,326-0,813+\frac{\pi}{2}
\right)
=
46,4\,\text{V}
\]

e

\[
P_L=v_L\,i
\]

\[
P_L
=
(46,4)(0,947)
\]

\[
\boxed{P_L=43,9\,\text{W}}
\]

\subsubsection*{(d) Potência dissipada no resistor}

\[
P_R=i^2R
\]

\[
P_R=(0,947)^2(16)
\]

\[
\boxed{P_R=14,4\,\text{W}}
\]

\subsubsection*{(e) Comparação entre as potências}

\[
P_C+P_L+P_R
=
-17,0+43,9+14,4
=
41,3\,\text{W}
\]

enquanto

\[
P_G=41,4\,\text{W}
\]

A pequena diferença decorre dos arredondamentos efetuados. Portanto,

\[
\boxed{
P_G=P_C+P_L+P_R
}
\]

confirmando a conservação de energia no circuito.

\subsection{Questão 59 — Fábio Rodrigues}
\textbf{Dados}

\[
R = 5,00\,\Omega
\]

\[
L = 60,0\,mH = 6,00 \times 10^{-2}\,H
\]

\[
f_d = 60,0\,Hz
\]

\[
\varepsilon_m = 30,0\,V
\]

A frequência angular do circuito é:

\[
\omega_d = 2\pi f_d = 2\pi(60,0) \approx 377\,rad/s
\]

\textbf{(a) Capacitância para potência máxima}

A potência dissipada é máxima na ressonância, onde a reatância indutiva é igual à reatância capacitiva:

\[
X_L = X_C
\]

\[
\omega_d L = \frac{1}{\omega_d C}
\]

\[
C = \frac{1}{\omega_d^2 L}
\]

Substituindo os valores:

\[
C = \frac{1}{(377)^2 (6,00 \times 10^{-2})}
\]

\[
C \approx 1,17 \times 10^{-4}\,F
\]

\[
\boxed{C = 117\,\mu F}
\]

\textbf{(b) Capacitância para potência mínima}

A potência é mínima quando a impedância \(Z\) é máxima. Sabendo que \(Z = \sqrt{R^2 + (X_L - X_C)^2}\), a impedância tende ao infinito quando a reatância capacitiva \(X_C \to \infty\). Isso ocorre quando a capacitância tende a zero:

\[
C \to 0
\]

\[
\boxed{C = 0}
\]

\textbf{(c) Dissipação máxima}

Na ressonância, a impedância é igual à resistência (\(Z = R\)). A potência máxima dissipada é:

\[
P_{\max} = \frac{\varepsilon_{\text{rms}}^2}{R} = \frac{(\varepsilon_m / \sqrt{2})^2}{R} = \frac{\varepsilon_m^2}{2R}
\]

Substituindo:

\[
P_{\max} = \frac{(30,0)^2}{2(5,00)}
\]

\[
P_{\max} = \frac{900}{10,0}
\]

\[
\boxed{P_{\max} = 90,0\,W}
\]

\textbf{(d) Ângulo de fase correspondente}

Na ressonância (\(X_L = X_C\)), a constante de fase \(\phi\) é:

\[
\tan \phi = \frac{X_L - X_C}{R} = 0
\]

\[
\boxed{\phi = 0^\circ}
\]

\textbf{(e) Fator de potência correspondente}

O fator de potência é o cosseno do ângulo de fase:

\[
\cos \phi = \cos(0^\circ)
\]

\[
\boxed{\cos \phi = 1,00}
\]

\textbf{(f) Dissipação mínima}

Quando \(C \to 0\), a impedância \(Z \to \infty\), fazendo com que a corrente no circuito tenda a zero. Logo:

\[
P_{\min} = 0
\]

\[
\boxed{P_{\min} = 0\,W}
\]

\textbf{(g) Ângulo de fase correspondente}

Quando \(C \to 0\), a reatância capacitiva \(X_C \to \infty\). O ângulo de fase tende para:

\[
\tan \phi = \frac{X_L - \infty}{R} \to -\infty
\]

\[
\boxed{\phi = -90^\circ}
\]

\textbf{(h) Fator de potência correspondente}

Com \(\phi = -90^\circ\), o fator de potência será:

\[
\cos \phi = \cos(-90^\circ)
\]

\[
\boxed{\cos \phi = 0}
\]
\subsection{Questão 60 — Rodrigo Albuquerque}


Um dimmer típico é composto por um indutor variável \(L\) ligado em série com uma lâmpada de \(120\ V\), \(1000\ W\). O circuito é alimentado com uma tensão de \(120\ V_{rms}\) e frequência de \(60\ Hz\).

\includegraphics[width=0.5\textwidth]{image_2.png}

\vspace{0.5cm}

\textbf{(a) Valor de \(L_{\max}\)}

Primeiro calculamos a resistência da lâmpada:

\[
P=\frac{V^2}{R}
\]

Logo:

\[
R=\frac{V^2}{P}
\]

Substituindo:

\[
R=\frac{(120)^2}{1000}
\]

\[
R=14,4\ \Omega
\]

Quando a potência mínima for \(200\ W\):

\[
P=I^2R
\]

Então:

\[
200=I^2(14,4)
\]

\[
I=\sqrt{\frac{200}{14,4}}
\]

\[
I\approx3,73\ A
\]

A impedância total do circuito é:

\[
Z=\frac{V}{I}
\]

\[
Z=\frac{120}{3,73}
\]

\[
Z\approx32,2\ \Omega
\]

Num circuito RL:

\[
Z=\sqrt{R^2+X_L^2}
\]

Logo:

\[
32,2^2=14,4^2+X_L^2
\]

\[
X_L^2=1036,8-207,4
\]

\[
X_L^2=829,4
\]

\[
X_L\approx28,8\ \Omega
\]

Sabendo que:

\[
X_L=\omega L
\]

e:

\[
\omega=2\pi f
\]

\[
\omega=2\pi(60)
\]

\[
\omega\approx377\ rad/s
\]

Então:

\[
L=\frac{X_L}{\omega}
\]

\[
L=\frac{28,8}{377}
\]

\[
L\approx0,076\ H
\]

\[
L_{\max}\approx7,6\times10^{-2}\ H
\]

\vspace{0.5cm}

\textbf{(b) É possível usar um resistor variável?}

Sim. Um resistor variável também reduziria a corrente do circuito.

\vspace{0.5cm}

\textbf{(c) Valor de \(R_{\max}\)}

Para potência mínima de \(200\ W\), a corrente já foi calculada:

\[
I\approx3,73\ A
\]

A resistência total necessária é:

\[
R_T=\frac{V}{I}
\]

\[
R_T=\frac{120}{3,73}
\]

\[
R_T\approx32,2\ \Omega
\]

Como a lâmpada possui:

\[
R=14,4\ \Omega
\]

o resistor adicional deve ser:

\[
R_{\max}=32,2-14,4
\]

\[
R_{\max}\approx17,8\ \Omega
\]

\vspace{0.5cm}

\textbf{(d) Por que não se usa esse método?}

Porque o resistor variável dissiparia muita potência em forma de calor, causando desperdício de energia.

Já o indutor praticamente não dissipa potência média, apenas limita a corrente do circuito.

\subsection{Questão 61 — Dávyla Maria}
Inserir resolução aqui.

\subsection{Questão 66 — Kaylane Castro}
\subsection*{Enunciado}

Um motor elétrico ligado a uma tomada de $120\,\text{V}$ e
$60\,\text{Hz}$ desenvolve uma potência mecânica de
$0,100\,\text{hp}$ $(1\,\text{hp}=746\,\text{W})$.

\begin{itemize}
    \item[(a)] Se o motor consome uma corrente rms de
    $0,650\,\text{A}$, qual é sua resistência efetiva em relação
    à transferência de energia?
    
    \item[(b)] Essa resistência efetiva é igual à resistência dos
    enrolamentos do motor, medida com um ohmímetro com o motor
    desligado da tomada?
\end{itemize}

\subsection*{Dados}

\[
V_{\text{rms}} = 120\,\text{V}
\]

\[
f = 60\,\text{Hz}
\]

\[
P = 0,100\,\text{hp}
\]

\[
1\,\text{hp} = 746\,\text{W}
\]

\[
I_{\text{rms}} = 0,650\,\text{A}
\]

Primeiro convertemos a potência para watts:

\[
P = 0,100 \times 746
\]

\[
P = 74,6\,\text{W}
\]

% ---------------------------------------------------------

\subsection*{(a) Resistência efetiva}

A potência elétrica associada à resistência efetiva é dada por:

\[
P = I_{\text{rms}}^2R
\]

Isolando $R$:

\[
R = \frac{P}{I_{\text{rms}}^2}
\]

Substituindo os valores:

\[
R = \frac{74,6}{(0,650)^2}
\]

\[
R = \frac{74,6}{0,4225}
\]

\[
R \approx 176,6\,\Omega
\]

Arredondando:

\[
\boxed{R \approx 177\,\Omega}
\]

% ---------------------------------------------------------

\subsection*{(b) Comparação com a resistência dos enrolamentos}

Não.

A resistência efetiva calculada acima não é igual à resistência
ôhmica dos enrolamentos medida com um ohmímetro.

O ohmímetro mede apenas a resistência elétrica do fio de cobre
das bobinas do motor.

Quando o motor está funcionando, aparece uma força
contraeletromotriz (fcem) devido à rotação do motor no campo
magnético. Essa fcem reduz a corrente do circuito e faz o motor
se comportar como se tivesse uma resistência muito maior.

Por isso, a resistência efetiva relacionada à transferência de
energia é muito maior que a resistência medida diretamente nos
enrolamentos.

\subsection{Questão 67 — Ana Alicy}
Inserir resolução aqui.

\subsection{Questão 68 — Ana Alicy}

Um capacitor de \SI{1,50}{\micro\farad} possui uma reatância capacitiva de \SI{12,0}{\ohm}. (a) Qual é a frequência de operação do circuito? (b) Qual será a reatância capacitiva do capacitor se a frequência for multiplicada por dois?

\vspace{0.8cm}


\textbf{Solução:}
\begin{quote}
\noindent A fórmula principal que relaciona a reatância capacitiva, a frequência de operação do circuito ($f$) e a capacitância ($C$) é:
\[
X_C = \frac{1}{\omega C} = \frac{1}{2\pi f C}
\]

\vspace{0.2cm}
\noindent Onde os dados fornecidos são:
\begin{itemize}
    \item $C = \SI{1,50}{\micro\farad} = 1,50 \times 10^{-6}\ \text{F}$
    \item $X_C = \SI{12,0}{\ohm}$
\end{itemize}

\vspace{0.4cm}
\noindent \textbf{(a) Qual é a frequência de operação do circuito?}

\vspace{0.2cm}
\noindent Para encontrar a frequência, precisamos isolar o $f$ na equação:
\[
f = \frac{1}{2\pi X_C C}
\]

\noindent Substituindo os valores no Sistema Internacional:
\[
f = \frac{1}{2\pi \cdot (12,0) \cdot (1,50 \times 10^{-6})}
\]
\[
f = \frac{1}{2\pi \cdot (18,0 \times 10^{-6})}
\]
\[
f = \frac{1}{113,097 \times 10^{-6}}
\]
\[
f \approx \SI{8841,9}{\hertz}
\]

\noindent Arredondando para três algarismos significativos (acompanhando a precisão dos dados do problema):

\resposta{(a) A frequência de operação é de aproximadamente \SI{8,84}{\kilo\hertz}.}

\vspace{0.4cm}
\noindent \textbf{(b) Qual será a reatância capacitiva se a frequência for multiplicada por dois?}

\vspace{0.2cm}
\noindent Aqui não precisamos refazer toda a conta do zero! Podemos resolver isso apenas analisando a proporção na fórmula.

\noindent Observe a equação da reatância:
\[
X_C = \frac{1}{2\pi f C}
\]

\noindent A reatância capacitiva ($X_C$) é inversamente proporcional à frequência ($f$). Isso significa que, se aumentarmos a frequência do sinal, o capacitor oferecerá menos oposição à passagem da corrente.

\noindent Portanto, se a frequência é multiplicada por dois ($2f$), a nova reatância ($X_C'$) será dividida por dois:
\[
X_C' = \frac{X_C}{2}
\]
\[
X_C' = \frac{12,0}{2}
\]
\[
X_C' = \SI{6,00}{\ohm}
\]

\resposta{(b) A nova reatância capacitiva será de \SI{6,00}{\ohm}.}
\end{quote}

\vspace{0.2cm}
\subsection{Questão 69 — Kaylane Castro}

\subsection*{Enunciado}

Em um certo circuito $RLC$ série, a força eletromotriz máxima
do gerador é $125\,\text{V}$ e a corrente máxima é
$3,20\,\text{A}$.

Se a corrente está adiantada de $0,982\,\text{rad}$ em relação
à força eletromotriz do gerador, determine:

\begin{itemize}
    \item[(a)] a impedância do circuito;
    \item[(b)] a resistência do circuito;
    \item[(c)] se o circuito é principalmente capacitivo ou indutivo.
\end{itemize}

\subsection*{Dados}

\[
\mathcal{E}_m = 125\,\text{V}
\]

\[
I = 3,20\,\text{A}
\]

\[
\phi = -0,982\,\text{rad}
\]

O sinal negativo é utilizado porque a corrente está adiantada em
relação à tensão.

% ---------------------------------------------------------

\subsection*{(a) Impedância do circuito}

A impedância é dada pela relação:

\[
Z=\frac{\mathcal{E}_m}{I}
\]

Substituindo os valores:

\[
Z=\frac{125}{3,20}
\]

\[
Z=39,0625\,\Omega
\]

Arredondando:

\[
\boxed{Z\approx 39,1\,\Omega}
\]

% ---------------------------------------------------------

\subsection*{(b) Resistência do circuito}

A resistência do circuito pode ser calculada por:

\[
R=Z\cos\phi
\]

Substituindo os valores:

\[
R=(39,0625)\cos(-0,982)
\]

Como:

\[
\cos(-0,982)\approx 0,555
\]

temos:

\[
R=(39,0625)(0,555)
\]

\[
R\approx 21,7\,\Omega
\]

Portanto:

\[
\boxed{R\approx 21,7\,\Omega}
\]

% ---------------------------------------------------------

\subsection*{(c) Tipo do circuito}

A corrente está adiantada em relação à tensão.

Em circuitos de corrente alternada:

\begin{itemize}
    \item Em circuitos indutivos, a corrente atrasa a tensão;
    \item Em circuitos capacitivos, a corrente adianta a tensão.
\end{itemize}

Como a corrente está adiantada, o circuito é predominantemente
capacitivo.

\[
\boxed{\text{Circuito principalmente capacitivo}}
\]

\subsection{Questão 70 — Dávyla Maria}
Inserir resolução aqui.

\subsection{Questão 71 — Rodrigo Albuquerque}

Em um circuito RLC série excitado com frequência de \(60,0\ Hz\), a tensão máxima no indutor é \(2,00\) vezes a tensão máxima no resistor e \(2,00\) vezes a tensão máxima no capacitor.

\vspace{0.5cm}

\textbf{(a) Ângulo de fase}

Sabemos que:

\[
V_L=2V_R
\]

e

\[
V_L=2V_C
\]

Logo:

\[
V_C=\frac{V_L}{2}
\]

A tensão resultante reativa é:

\[
V_X=V_L-V_C
\]

Substituindo:

\[
V_X=2V_R-\frac{2V_R}{2}
\]

\[
V_X=2V_R-V_R
\]

\[
V_X=V_R
\]

No circuito RLC:

\[
\tan\phi=\frac{V_L-V_C}{V_R}
\]

Então:

\[
\tan\phi=\frac{V_R}{V_R}
\]

\[
\tan\phi=1
\]

Logo:

\[
\phi=45^\circ
\]

Como \(V_L>V_C\), o circuito é indutivo, então a corrente atrasa em relação à força eletromotriz.

\[
\phi=45^\circ
\]

\vspace{0.5cm}

\textbf{(b) Resistência do circuito}

A força eletromotriz máxima é:

\[
\mathcal{E}_{max}=30,0\ V
\]

e a corrente máxima:

\[
I_{max}=300\ mA=0,300\ A
\]

A impedância é:

\[
Z=\frac{\mathcal{E}_{max}}{I_{max}}
\]

\[
Z=\frac{30,0}{0,300}
\]

\[
Z=100\ \Omega
\]

Num circuito RLC:

\[
\cos\phi=\frac{R}{Z}
\]

Como:

\[
\phi=45^\circ
\]

\[
\cos45^\circ=\frac{R}{100}
\]

\[
\frac{\sqrt2}{2}=\frac{R}{100}
\]

\[
R=100\cdot\frac{\sqrt2}{2}
\]

\[
R\approx70,7\ \Omega
\]


\subsection{Questão 72 — Fábio Rodrigues}
\textbf{Dados}

\[
L = 1,30\,mH = 1,30 \times 10^{-3}\,H
\]

\[
f = 3,50\,kHz = 3,50 \times 10^3\,Hz
\]

A frequência de ressonância de um circuito LC é dada por:

\[
f = \frac{1}{2\pi\sqrt{LC}}
\]

Elevando ambos os lados ao quadrado para isolar a capacitância \(C\):

\[
f^2 = \frac{1}{4\pi^2 L C}
\]

Multiplicando cruzado para isolar \(C\):

\[
C = \frac{1}{4\pi^2 f^2 L}
\]

Substituindo os valores fornecidos:

\[
C = \frac{1}{4\pi^2 (3,50 \times 10^3)^2 (1,30 \times 10^{-3})}
\]

\[
C = \frac{1}{4\pi^2 (12,25 \times 10^6) (1,30 \times 10^{-3})}
\]

\[
C = \frac{1}{4\pi^2 (15925)}
\]

\[
C \approx \frac{1}{628693,5}
\]

\[
C \approx 1,59 \times 10^{-6}\,F
\]

\[
\boxed{C = 1,59\,\mu F}
\]

\subsection{Questão 73 — Mateus Valentim}
\subsubsection{Resolução}

A frequência angular do circuito é

\[
\omega=2\pi f
=2\pi(10{,}4\times 10^3)
=6{,}53\times 10^4\,\text{rad/s}.
\]

\begin{enumerate}
    \item[(a)] Para um circuito \(LC\),

    \[
    \omega=\frac{1}{\sqrt{LC}}.
    \]

    Portanto,

    \[
    L=\frac{1}{\omega^2 C}.
    \]

    Substituindo os valores,

    \[
    L=
    \frac{1}{(6{,}53\times 10^4)^2(340\times 10^{-6})}
    =
    6{,}89\times 10^{-7}\,\text{H}.
    \]

    \item[(b)] A energia total é igual à energia máxima no indutor:

    \[
    U=\frac{1}{2}LI_{\max}^2.
    \]

    Logo,

    \[
    U=
    \frac{1}{2}(6{,}89\times 10^{-7})(7{,}20\times 10^{-3})^2
    =
    1{,}79\times 10^{-11}\,\text{J}.
    \]

    \item[(c)] Como

    \[
    I_{\max}=\omega Q_{\max},
    \]

    temos

    \[
    Q_{\max}=\frac{I_{\max}}{\omega}
    =
    \frac{7{,}20\times 10^{-3}}{6{,}53\times 10^4}
    =
    1{,}10\times 10^{-7}\,\text{C}.
    \]
\end{enumerate}

\subsection{Questão 74 — Cicero Rodrigues}

\textbf{74. }\normalfont Um circuito LC oscilante tem uma indutância de 3,00 mH e uma capacitância de 10,0 $\mu$F. Determine \textbf{(a)} a frequência angular e \textbf{(b)} o período de oscilação. \textbf{(c)} No instante $t = 0$, o capacitor é carregado com 200 $\mu$C e a corrente é zero. Faça um esboço da carga do capacitor em função do tempo.

\textbf{Resolução:}
\vspace{0.3cm}

\textbf{Dados:}
\[L = 3,00\;mH \Rightarrow 3,00 \times 10^{-3}\;H\]
\[C = 10,0\,\mu F \Rightarrow 10,0 \times 10^{-6}\;F\]
\[Q_{\max} = 200\,\mu C \Rightarrow 200 \times 10^{-6}\;C\]

\vspace{0.2cm}
\textbf{(a) Frequência angular (\(\omega\))}
A frequência angular de um circuito LC é dada por:
\[\omega = \frac{1}{\sqrt{LC}}\]
Substituindo os valores:
\[\omega = \frac{1}{\sqrt{(3,00 \times 10^{-3})(10,0 \times 10^{-6})}}\]
\[\omega = \frac{1}{\sqrt{30,0 \times 10^{-9}}}\]
\[\omega = \frac{1}{1,732 \times 10^{-4}}\]
\[\omega \approx 5773,5\;rad/s\]
\[\underline{\omega \approx 5,77 \times 10^3\;rad/s}\]

\vspace{0.2cm}
\textbf{(b) Período de oscilação (\(T\))}
O período é a relação entre \(2\pi\) e a frequência angular:
\[T = \frac{2\pi}{\omega}\]
\[T = \frac{2\pi}{5773,5}\]
\[T \approx 1,088 \times 10^{-3}\;s\]
\[\underline{T \approx 1,09\;ms}\]

\vspace{0.2cm}
\textbf{(c) Esboço da carga em função do tempo}
Como em \(t = 0\) a carga é máxima e a corrente é nula, a função obedece a um cosseno:
\[q(t) = Q_{\max} \cos(\omega t)\]
\[q(t) = 200 \cos(5774\,t)\;\;(\mu C)\]
O esboço gráfico consiste em uma curva cossenoide que se inicia no pico positivo de \(200\,\mu C\) no instante \(t = 0\), cruza o zero, atinge o pico negativo de \(-200\,\mu C\) em \(t = T/2 \approx 0,545\;ms\) e retorna ao pico positivo em \(t = T \approx 1,09\;ms\), repetindo o ciclo.


\vspace{0.5cm}
\textbf{Esboço do gráfico:}

\begin{center}
\begin{tikzpicture}
\begin{axis}[
    width=13cm,
    height=7cm,
    xlabel={$t\;(ms)$},
    ylabel={$q(t)\;(\mu C)$},
    xmin=0,
    xmax=4,
    ymin=-220,
    ymax=220,
    samples=300,
    domain=0:4,
    grid=both,
]

\addplot[
    thick
]
{200*cos(deg(5770*x/1000))};

\end{axis}
\end{tikzpicture}
\end{center}

\subsection{Questão 75 — Guilherme Araújo Floriano}

\textbf{(a)} Pela Eq.~31-4,
\[
L = (\omega^2C)^{-1} = \bigl[(2\pi f)^2C\bigr]^{-1}=2.41\,\mu\mathrm{H}.
\]

\textbf{(b)} A energia total é a energia máxima em qualquer elemento:
\[
U_{\max} = \frac{1}{2}LI^2 = 21.4\,\mathrm{pJ}.
\]

\textbf{(c)} Usando a relação do capacitor,
\[
U_{\max} = \frac{Q^2}{2C}
\quad\Rightarrow\quad
Q = \sqrt{2CU_{\max}} = 82.2\,\mathrm{nC}.
\]

\subsection{Questão 80 — Arthur Roberto da Silva}

\subsection*{Dados do enunciado:}
\begin{itemize}
    \item Indutância: $L = 1,50\text{ mH} = 1,50 \times 10^{-3}\text{ H}$
    \item Energia máxima no indutor: $U_B = 10,0\text{ }\mu\text{J} = 10,0 \times 10^{-6}\text{ J}$
\end{itemize}

\subsection*{Cálculo da corrente máxima}

A energia magnética máxima armazenada em um indutor pertencente a um circuito $LC$ oscilante é relacionada com a corrente máxima ($I$) pela seguinte expressão:
\begin{equation*}
U_B = \frac{1}{2} L I^2
\end{equation*}

Isolando a corrente máxima ($I$):
\begin{equation*}
I^2 = \frac{2 U_B}{L} \implies I = \sqrt{\frac{2 U_B}{L}}
\end{equation*}

Substituindo os valores numéricos fornecidos:
\begin{align*}
I &= \sqrt{\frac{2 \cdot (10,0 \times 10^{-6}\text{ J})}{1,50 \times 10^{-3}\text{ H}}} \\
I &= \sqrt{\frac{20,0 \times 10^{-6}}{1,50 \times 10^{-3}}} \\
I &= \sqrt{13,33 \times 10^{-3}} \\
I &= \sqrt{1,333 \times 10^{-2}} \\
I &\approx 0,1155\text{ A} = 115,5\text{ mA}
\end{align*}

\subsection*{Resposta Final}
A corrente máxima no circuito é de aproximadamente \textbf{$0,116\text{ A}$} (ou \textbf{$116\text{ mA}$}).



% ================= CAPÍTULO 32 =================

\section{Capítulo 32}

\subsection{Questão 3 — Arthur Roberto da Silva}
\subsection*{(a)}
Utilizamos a lei de Gauss para o magnetismo: $\oint \vec{B} \cdot d\vec{A} = 0$. Portanto,

\begin{equation*}
\oint \vec{B} \cdot d\vec{A} = \Phi_1 + \Phi_2 + \Phi_C \,,
\end{equation*}

\noindent onde $\Phi_1$ é o fluxo magnético através da primeira extremidade mencionada, $\Phi_2$ é o fluxo magnético através da segunda extremidade mencionada e $\Phi_C$ é o fluxo magnético através da superfície curva lateral. 

Sobre a primeira extremidade, o campo magnético aponta para dentro, de modo que o fluxo é $\Phi_1 = -25,0\ \mu\text{Wb}$. Sobre a segunda extremidade, o campo magnético é uniforme, normal à superfície e aponta para fora, de forma que o fluxo é $\Phi_2 = AB = \pi r^2 B$, onde $A$ é a área da extremidade e $r$ é o raio do cilindro. Seu valor é:

\begin{equation*}
\Phi_2 = \pi (0,120\text{ m})^2 (1,60 \times 10^{-3}\text{ T}) = +7,24 \times 10^{-5}\text{ Wb} = +72,4\ \mu\text{Wb} \,.
\end{equation*}

Como a soma dos três fluxos deve ser igual a zero, temos:

\begin{equation*}
\Phi_C = -\Phi_1 - \Phi_2 = 25,0\ \mu\text{Wb} - 72,4\ \mu\text{Wb} = -47,4\ \mu\text{Wb} \,.
\end{equation*}

Portanto, o módulo é $|\Phi_C| = 47,4\ \mu\text{Wb}$.

\subsection*{(b)}
O sinal de menos em $\Phi_C$ indica que o fluxo magnético aponta para dentro através da superfície curva lateral.

\subsection{Questão 7 — Guilherme Araújo Floriano}

	\textbf{(a)} Com $E=V/d$ e $dV/dt = -V_{\max}\omega\sin(\omega t)$, a lei de Ampère-Maxwell dá, para $r\le R$,
\[
B(r,t)=\frac{\mu_0\varepsilon_0\,r}{2d}\,\frac{dV}{dt}.
\]
No máximo,
\[
B_{\max}(r)=\frac{\mu_0\varepsilon_0\,r\,V_{\max}\omega}{2d}.
\]

Em $r=R=30 mm$, com $V_{\max}= 150 V$ e $\omega\approx {377}{s^{-1}}$,

\[
B_{\max}=1.9\times10^{-12}\,\mathrm{T}.
\]

	\textbf{(b)} Para $r\ge R$,
\[
B_{\max}(r)=\frac{\mu_0\varepsilon_0\,R^2V_{\max}\omega}{2rd},
\]
mostrando a dependência $B\propto r^{-1}$ fora da região central.

A Figura abaixo ilustra o gráfico do campo magnético com relação ao raio, onde até a distância igual ao raio da esfera, o campo cresce linearmente e decresce pela recíproca após esse raio.

\begin{figure}[ht]
    \centering
    \begin{tikzpicture}
        \begin{axis}[
            axis lines = left,
            xlabel = {$r$ (mm)},
            ylabel = {$B_{\max}$ ($\times 10^{-12}$ T)},
            xmin = 0, xmax = 60,
            ymin = 0, ymax = 2.2,
            xtick = {0, 15, 30, 45, 60},
            ytick = {0, 0.5, 1.0, 1.5, 1.9, 2.2},
            grid = both,
            grid style = {dashed, gray!30},
            width = 11cm,
            height = 7.5cm,
            title = {Gráfico de $B_{\max}$ em função de $r$ ($B_{\text{pico}} = 1,9 \times 10^{-12}$ T)},
            legend pos = north east
        ]
            \addplot [
                domain=0:30, 
                samples=100, 
                color=blue,
                ultra thick
            ] {(1.9 / 30) * x};
            \addlegendentry{$r < R$ (Linear)}
            
            \addplot [
                domain=30:60, 
                samples=100, 
                color=red,
                ultra thick
            ] {57 / x};
            \addlegendentry{$r \geq R$ ($1/r$)}
            
            \draw[dashed, thick, black!60] (axis cs:30,0) -- (axis cs:30,1.9);
            \draw[dashed, thick, black!60] (axis cs:0,1.9) -- (axis cs:30,1.9);
            
            \node[circle, fill=black, inner sep=1.5pt] at (axis cs:30,1.9) {};
        \end{axis}
    \end{tikzpicture}
\end{figure}
\newpage
\subsection{Questão 8 — Cicero Rodrigues}

\textbf{8. }\normalfont Fluxo elétrico não uniforme. A Fig. 32.12 mostra uma região circular de raio $R = 3,00$ cm na qual um fluxo elétrico aponta para fora do papel. O fluxo elétrico envolvido por uma circunferência concêntrica de raio $r$ é dado por $\Phi_{E,env} = (0,600\;V\cdot m/s)(r/R)t$, em que $r \le R$ e $t$ está em segundos. Determine o módulo do campo magnético induzido a uma distância radial \textbf{(a)} de 2,00 cm e \textbf{(b)} de 5,00 cm.

\includegraphics[width=0.2\textwidth]{image_11.png}

\textbf{Resolução:}
\vspace{0.3cm}

\textbf{Dados:}
\[R = 3,00\;cm = 0,03\;m\]
\[\Phi_{E,env} = 0,600 \left(\frac{r}{R}\right) t\]
\[\frac{d\Phi_{E,env}}{dt} = 0,600 \left(\frac{r}{R}\right)\]
\[\mu_0 = 4\pi \times 10^{-7}\;T\cdot m/A\]
\[\epsilon_0 = 8,85 \times 10^{-12}\;F/m\]

\vspace{0.2cm}
\textbf{(a) Campo magnético a \(r = 2,00\;cm\)}
O ponto de medição está dentro da região de fluxo (\(r \le R\)). Aplicando a Lei de Ampère-Maxwell:
\[\oint \vec{B} \cdot d\vec{s} = \mu_0 \epsilon_0 \frac{d\Phi_{E,env}}{dt}\]
\[B(2\pi r) = \mu_0 \epsilon_0 \left[ 0,600 \left(\frac{r}{R}\right) \right]\]
Observe que o termo \(r\) se cancela de ambos os lados da equação:
\[B(2\pi) = \mu_0 \epsilon_0 \left( \frac{0,600}{R} \right) \Rightarrow B = \frac{\mu_0 \epsilon_0 \cdot 0,600}{2\pi R}\]
Substituindo os valores numéricos:
\[B = \frac{(4\pi \times 10^{-7})(8,85 \times 10^{-12})(0,600)}{2\pi (0,03)}\]
\[B = \frac{2 \times 10^{-7} \cdot 8,85 \times 10^{-12} \cdot 0,600}{0,03}\]
\[B = \frac{10,62 \times 10^{-19}}{0,03}\]
\[B = 354 \times 10^{-19}\;T\]
\[\underline{B = 3,54 \times 10^{-17}\;T}\]

\vspace{0.2cm}
\textbf{(b) Campo magnético a \(r = 5,00\;cm\)}
O ponto de medição está fora da região de fluxo (\(r > R\)). O fluxo elétrico envolvido pela trajetória é o fluxo total disponível, que avaliamos no limite \(r = R\):
\[\Phi_{E,total} = 0,600 \left(\frac{R}{R}\right) t = 0,600\,t\]
\[\frac{d\Phi_{E,total}}{dt} = 0,600\;V\cdot m/s\]
Aplicando novamente a Lei de Ampère-Maxwell:
\[B(2\pi r) = \mu_0 \epsilon_0 \frac{d\Phi_{E,total}}{dt}\]
\[B = \frac{\mu_0 \epsilon_0 (0,600)}{2\pi r}\]
Substituindo os valores (lembrando que agora \(r = 0,05\;m\)):
\[B = \frac{(4\pi \times 10^{-7})(8,85 \times 10^{-12})(0,600)}{2\pi (0,05)}\]
\[B = \frac{2 \times 10^{-7} \cdot 8,85 \times 10^{-12} \cdot 0,600}{0,05}\]
\[B = \frac{10,62 \times 10^{-19}}{0,05}\]
\[B = 212,4 \times 10^{-19}\;T\]
\[\underline{B \approx 2,12 \times 10^{-17}\;T}\]

\subsection{Questão 10 — Mateus Valentim}
\subsubsection{Resolução}

***
Usamos a lei de Ampere-Maxwell:

\[
\oint \vec{B}\cdot d\vec{s}
=
\mu_0\epsilon_0\frac{d\Phi_E}{dt}.
\]

Por simetria, o campo magnético induzido é tangente a uma circunferência de
raio \(r\), logo

\[
B(2\pi r)=\mu_0\epsilon_0\frac{d\Phi_{E,\text{env}}}{dt}.
\]

\begin{enumerate}
    \item[(a)] Para \(r=2{,}00\,\text{cm}<R\),

    \[
    \frac{d\Phi_{E,\text{env}}}{dt}
    =
    (0{,}600)\frac{r}{R}.
    \]

    Assim,

    \[
    B=
    \frac{\mu_0\epsilon_0(0{,}600)(r/R)}{2\pi r}
    =
    \frac{\mu_0\epsilon_0(0{,}600)}{2\pi R}.
    \]

    Com \(R=3{,}00\,\text{cm}=3{,}00\times 10^{-2}\,\text{m}\),

    \[
    B=3{,}54\times 10^{-17}\,\text{T}.
    \]

    \item[(b)] Para \(r=5{,}00\,\text{cm}>R\), o fluxo envolvido é o fluxo
    total da região circular:

    \[
    \frac{d\Phi_E}{dt}=0{,}600.
    \]

    Portanto,

    \[
    B=
    \frac{\mu_0\epsilon_0(0{,}600)}{2\pi r}
    =
    \frac{(4\pi\times 10^{-7})(8{,}85\times 10^{-12})(0{,}600)}
    {2\pi(5{,}00\times 10^{-2})}.
    \]

    Logo,

    \[
    B=2{,}13\times 10^{-17}\,\text{T}.
    \]
\end{enumerate}

\subsection{Questão 11 — Fábio Rodrigues}
\textbf{Dados}

\[
R = 3,00\,\text{cm} = 0,0300\,\text{m}
\]

\[
\frac{dE}{dt} = 4,50 \times 10^{-3}\,\text{V/(m}\cdot\text{s)}
\]

\[
\mu_0 = 4\pi \times 10^{-7}\,\text{T}\cdot\text{m/A}
\]

\[
\varepsilon_0 = 8,85 \times 10^{-12}\,\text{C}^2/\text{(N}\cdot\text{m}^2)
\]

\[
\mu_0\varepsilon_0 \approx 1,11 \times 10^{-17}\,\text{s}^2/\text{m}^2
\]

Pela Lei de Ampère-Maxwell, o campo magnético induzido \(B\) ao longo de uma trajetória circular de raio \(r\) é dado por:

\[
\oint \vec{B} \cdot d\vec{s} = \mu_0 \varepsilon_0 \frac{d\Phi_E}{dt}
\]

\[
B(2\pi r) = \mu_0 \varepsilon_0 A_{\text{enc}} \frac{dE}{dt}
\]

\textbf{(a) Módulo do campo magnético a \(r = 2,00\,\text{cm}\)}

Como \(r < R\), a área envolvida pelo campo elétrico é \(A_{\text{enc}} = \pi r^2\).

\[
B(2\pi r) = \mu_0 \varepsilon_0 (\pi r^2) \frac{dE}{dt}
\]

\[
B = \frac{\mu_0 \varepsilon_0 r}{2} \frac{dE}{dt}
\]

Substituindo os valores:

\[
B = \frac{(1,11 \times 10^{-17})(0,0200)}{2} (4,50 \times 10^{-3})
\]

\[
B = (1,11 \times 10^{-17})(0,0100)(4,50 \times 10^{-3})
\]

\[
B \approx 5,00 \times 10^{-22}\,\text{T}
\]

\[
\boxed{B = 5,00 \times 10^{-22}\,\text{T}}
\]

\textbf{(b) Módulo do campo magnético a \(r = 5,00\,\text{cm}\)}

Como \(r > R\), o campo elétrico só existe até o raio \(R\). A área envolvida pelo fluxo elétrico é fixa em \(A_{\text{enc}} = \pi R^2\).

\[
B(2\pi r) = \mu_0 \varepsilon_0 (\pi R^2) \frac{dE}{dt}
\]

\[
B = \frac{\mu_0 \varepsilon_0 R^2}{2r} \frac{dE}{dt}
\]

Substituindo os valores:

\[
B = \frac{(1,11 \times 10^{-17})(0,0300)^2}{2(0,0500)} (4,50 \times 10^{-3})
\]

\[
B = \frac{(1,11 \times 10^{-17})(0,0009)}{0,100} (4,50 \times 10^{-3})
\]

\[
B = (1,11 \times 10^{-17})(0,009)(4,50 \times 10^{-3})
\]

\[
B \approx 4,50 \times 10^{-22}\,\text{T}
\]

\[
\boxed{B = 4,50 \times 10^{-22}\,\text{T}}
\]

\subsection{Questão 12 — Rodrigo Albuquerque}

Um campo elétrico aponta para fora do papel em uma região circular de raio

\[
R=3,00\ cm=0,0300\ m
\]

O módulo do campo elétrico é:

\[
E=(0,500\ \text{V/m}\cdot\text{s})\left(1-\frac{r}{R}\right)t
\]

Deseja-se determinar o campo magnético induzido.

\vspace{0.5cm}

Pela lei de Ampère-Maxwell:

\[
\oint \vec B\cdot d\vec s
=
\mu_0\varepsilon_0
\frac{d\Phi_E}{dt}
\]

Como o problema possui simetria circular:

\[
B(2\pi r)
=
\mu_0\varepsilon_0
\frac{d\Phi_E}{dt}
\]

\vspace{0.5cm}

\textbf{(a) Para \(r=2,00\ cm\)}

\[
r=0,0200\ m
\]

Como \(r<R\), o fluxo elétrico é:

\[
\Phi_E
=
\int E\,dA
\]

\[
\Phi_E
=
\int_0^r
(0,500)\left(1-\frac{r'}{R}\right)t
(2\pi r' dr')
\]

Derivando em relação ao tempo:

\[
\frac{d\Phi_E}{dt}
=
0,500\cdot2\pi
\int_0^r
\left(1-\frac{r'}{R}\right)r' dr'
\]

Integrando:

\[
\frac{d\Phi_E}{dt}
=
\pi
\left[
\frac{r'^2}{2}
-
\frac{r'^3}{3R}
\right]_0^r
\]

\[
\frac{d\Phi_E}{dt}
=
\pi
\left(
\frac{r^2}{2}
-
\frac{r^3}{3R}
\right)
\]

Substituindo \(r=0,0200\ m\):

\[
\frac{d\Phi_E}{dt}
=
\pi
\left(
\frac{(0,0200)^2}{2}
-
\frac{(0,0200)^3}{3(0,0300)}
\right)
\]

\[
\frac{d\Phi_E}{dt}
=
4,89\times10^{-4}
\]

Então:

\[
B=
\frac{\mu_0\varepsilon_0}{2\pi r}
\frac{d\Phi_E}{dt}
\]

Substituindo:

\[
\mu_0\varepsilon_0
=
1,11\times10^{-17}
\]

\[
B=
\frac{(1,11\times10^{-17})(4,89\times10^{-4})}
{2\pi(0,0200)}
\]

\[
B\approx4,3\times10^{-20}\ T
\]

\vspace{0.5cm}

\textbf{(b) Para \(r=5,00\ cm\)}

\[
r=0,0500\ m
\]

Agora \(r>R\), então o fluxo deve ser calculado usando toda a região de raio \(R\).

\[
\frac{d\Phi_E}{dt}
=
\pi
\left(
\frac{R^2}{2}
-
\frac{R^3}{3R}
\right)
\]

\[
\frac{d\Phi_E}{dt}
=
\pi R^2
\left(
\frac12-\frac13
\right)
\]

\[
\frac{d\Phi_E}{dt}
=
\frac{\pi R^2}{6}
\]

Substituindo:

\[
\frac{d\Phi_E}{dt}
=
\frac{\pi(0,0300)^2}{6}
\]

\[
\frac{d\Phi_E}{dt}
=
4,71\times10^{-4}
\]

Então:

\[
B=
\frac{(1,11\times10^{-17})(4,71\times10^{-4})}
{2\pi(0,0500)}
\]

\[
B\approx1,7\times10^{-20}\ T
\]
```


\subsection{Questão 16 — Dávyla Maria}
Inserir resolução aqui.

\subsection{Questão 18 — Kaylane Castro}
\subsection*{Enunciado}

O módulo do campo elétrico entre duas placas paralelas circulares é dado por

\[
E(t) = (4,0\times10^5) - (6,0\times10^4)t
\]

com $E$ em volts por metro e $t$ em segundos.

Em $t=0$, o vetor campo elétrico aponta para cima.

A área das placas é:

\[
A = 4,0\times10^{-2}\,\text{m}^2
\]

Para $t \ge 0$, determine:

\begin{itemize}
    \item[(a)] o módulo da corrente de deslocamento;
    \item[(b)] o sentido da corrente de deslocamento;
    \item[(c)] o sentido do campo magnético induzido.
\end{itemize}

\subsection*{Dados}

\[
E(t)=4,0\times10^5-6,0\times10^4 t
\]

\[
A=4,0\times10^{-2}\,\text{m}^2
\]

\[
\varepsilon_0 = 8,85\times10^{-12}\,\text{F/m}
\]

% ---------------------------------------------------------

\subsection*{(a) Corrente de deslocamento}

A corrente de deslocamento é dada por:

\[
i_d=\varepsilon_0\frac{d\Phi_E}{dt}
\]

Como o campo elétrico é uniforme:

\[
\Phi_E = EA
\]

Então:

\[
i_d=\varepsilon_0 A\frac{dE}{dt}
\]

Derivando o campo elétrico:

\[
\frac{dE}{dt}=
\frac{d}{dt}
\left(
4,0\times10^5-6,0\times10^4 t
\right)
\]

\[
\frac{dE}{dt}=-6,0\times10^4
\]

Usando o módulo:

\[
|i_d|=
\varepsilon_0 A
\left|
\frac{dE}{dt}
\right|
\]

Substituindo os valores:

\[
|i_d|=
(8,85\times10^{-12})
(4,0\times10^{-2})
(6,0\times10^4)
\]

\[
|i_d|\approx 2,12\times10^{-8}\,\text{A}
\]

Portanto:

\[
\boxed{|i_d|\approx 2,1\times10^{-8}\,\text{A}}
\]

% ---------------------------------------------------------

\subsection*{(b) Sentido da corrente de deslocamento}

Como:

\[
\frac{dE}{dt}<0
\]

o campo elétrico está diminuindo com o tempo.

A corrente de deslocamento possui sentido oposto ao campo elétrico.

Como o campo elétrico aponta para cima, a corrente de deslocamento aponta para baixo.

\[
\boxed{\text{Sentido para baixo}}
\]

% ---------------------------------------------------------

\subsection*{(c) Sentido do campo magnético induzido}

A corrente de deslocamento aponta para baixo.

Aplicando a regra da mão direita:

\begin{itemize}
    \item polegar apontando para baixo;
    \item os dedos indicam o sentido do campo magnético.
\end{itemize}

Assim, o campo magnético circula no sentido horário.

\[
\boxed{\text{Sentido horário}}
\]
\subsection{Questão 19 — Ana Alicy}
Inserir resolução aqui.

\subsection{Questão 20 — Ana Alicy}
Inserir resolução aqui.

\subsection{Questão 21 — Kaylane Castro}
\subsection*{Enunciado}

Um capacitor de placas paralelas possui placas quadradas de lado

\[
L = 1,0\,\text{m}
\]

Uma corrente de

\[
i = 2,0\,\text{A}
\]

carrega o capacitor, produzindo um campo elétrico uniforme
$\vec{E}$ entre as placas.

Determine:

\begin{itemize}
    \item[(a)] a corrente de deslocamento na região entre as placas;
    \item[(b)] o valor de $\dfrac{dE}{dt}$;
    \item[(c)] a corrente de deslocamento envolvida por uma trajetória quadrada de lado
    $d = 0,50\,\text{m}$;
    \item[(d)] o valor de
    $\displaystyle \oint \vec{B}\cdot d\vec{s}$.
\end{itemize}

\subsection*{Dados}

\[
L = 1,0\,\text{m}
\]

\[
A = L^2 = 1,0\,\text{m}^2
\]

\[
i = 2,0\,\text{A}
\]

\[
d = 0,50\,\text{m}
\]

\[
A_{\text{int}} = d^2 = 0,25\,\text{m}^2
\]

\[
\varepsilon_0 = 8,85\times10^{-12}\,\text{F/m}
\]

\[
\mu_0 = 4\pi\times10^{-7}\,\text{T}\cdot\text{m/A}
\]

% ---------------------------------------------------------

\subsection*{(a) Corrente de deslocamento}

A corrente de deslocamento entre as placas é igual à corrente de
condução do circuito:

\[
i_d = i
\]

Portanto:

\[
\boxed{i_d = 2,0\,\text{A}}
\]

% ---------------------------------------------------------

\subsection*{(b) Cálculo de \texorpdfstring{$dE/dt$}{dE/dt}}

A corrente de deslocamento é dada por:

\[
i_d = \varepsilon_0 \frac{d\Phi_E}{dt}
\]

Como o campo elétrico é uniforme:

\[
\Phi_E = EA
\]

Então:

\[
i_d = \varepsilon_0 A \frac{dE}{dt}
\]

Isolando $\dfrac{dE}{dt}$:

\[
\frac{dE}{dt}
=
\frac{i_d}{\varepsilon_0 A}
\]

Substituindo os valores:

\[
\frac{dE}{dt}
=
\frac{2,0}
{(8,85\times10^{-12})(1,0)}
\]

\[
\frac{dE}{dt}
\approx 2,26\times10^{11}\,\text{V/(m}\cdot\text{s)}
\]

Portanto:

\[
\boxed{
\frac{dE}{dt}
\approx 2,3\times10^{11}\,\text{V/(m}\cdot\text{s)}
}
\]

% ---------------------------------------------------------

\subsection*{(c) Corrente de deslocamento envolvida}

Como o campo elétrico é uniforme, a corrente de deslocamento é
proporcional à área envolvida.

Assim:

\[
i_{d,\text{env}}
=
i_d
\left(
\frac{A_{\text{int}}}{A}
\right)
\]

Substituindo os valores:

\[
i_{d,\text{env}}
=
2,0
\left(
\frac{0,25}{1,0}
\right)
\]

\[
i_{d,\text{env}} = 0,50\,\text{A}
\]

Portanto:

\[
\boxed{i_{d,\text{env}} = 0,50\,\text{A}}
\]

% ---------------------------------------------------------

\subsection*{(d) Integral de linha}

Pela Lei de Ampère-Maxwell:

\[
\oint \vec{B}\cdot d\vec{s}
=
\mu_0 i_{d,\text{env}}
\]

Substituindo os valores:

\[
\oint \vec{B}\cdot d\vec{s}
=
(4\pi\times10^{-7})(0,50)
\]

\[
\oint \vec{B}\cdot d\vec{s}
=
2\pi\times10^{-7}
\]

\[
\oint \vec{B}\cdot d\vec{s}
\approx 6,28\times10^{-7}\,\text{T}\cdot\text{m}
\]

Portanto:

\[
\boxed{
\oint \vec{B}\cdot d\vec{s}
\approx 6,3\times10^{-7}\,\text{T}\cdot\text{m}
}
\]

\subsection{Questão 22 — Dávyla Maria}
Inserir resolução aqui.

\subsection{Questão 23 — Rodrigo Albuquerque}

Um fio de prata tem uma resistividade $\rho = 1,62 \times 10^{-8}\ \Omega\cdot\text{m}$ e uma seção reta de $5,00\ \text{mm}^2$. A corrente no fio é uniforme e varia à taxa de $2000\ \text{A/s}$ quando a corrente é $100\ \text{A}$. (a) Determine o módulo do campo elétrico (uniforme) no fio quando a corrente é $100\ \text{A}$. (b) Determine a corrente de deslocamento no fio nesse instante. (c) Determine a razão entre o módulo do campo magnético produzido pela corrente de deslocamento e o módulo do campo magnético produzido pela corrente a uma distância $r$ do fio.


Um fio de prata possui:

\[
\rho=1,62\times10^{-8}\ \Omega\cdot m
\]

\[
A=5,00\ mm^2=5,00\times10^{-6}\ m^2
\]

A corrente varia à taxa:

\[
\frac{dI}{dt}=2000\ A/s
\]

quando:

\[
I=100\ A
\]

\vspace{0.5cm}

\textbf{(a) Campo elétrico no fio}

Sabemos que:

\[
J=\frac{I}{A}
\]

Substituindo:

\[
J=\frac{100}{5,00\times10^{-6}}
\]

\[
J=2,00\times10^7\ A/m^2
\]

Pela forma microscópica da lei de Ohm:

\[
E=\rho J
\]

Então:

\[
E=(1,62\times10^{-8})(2,00\times10^7)
\]

\[
E=3,24\times10^{-1}\ V/m
\]

\[
E=0,324\ V/m
\]

\vspace{0.5cm}

\textbf{(b) Corrente de deslocamento}

A corrente de deslocamento é:

\[
I_d=\varepsilon_0\frac{d\Phi_E}{dt}
\]

Como o campo elétrico é uniforme:

\[
\Phi_E=EA
\]

Logo:

\[
I_d=\varepsilon_0A\frac{dE}{dt}
\]

Mas:

\[
E=\rho\frac{I}{A}
\]

Derivando:

\[
\frac{dE}{dt}
=
\frac{\rho}{A}\frac{dI}{dt}
\]

Substituindo:

\[
I_d
=
\varepsilon_0A
\left(
\frac{\rho}{A}\frac{dI}{dt}
\right)
\]

\[
I_d
=
\varepsilon_0\rho\frac{dI}{dt}
\]

Substituindo os valores:

\[
I_d
=
(8,85\times10^{-12})
(1,62\times10^{-8})
(2000)
\]

\[
I_d\approx2,87\times10^{-16}\ A
\]

\vspace{0.5cm}

\textbf{(c) Razão entre os campos magnéticos}

O campo magnético devido à corrente normal é:

\[
B=\frac{\mu_0 I}{2\pi r}
\]

O campo magnético devido à corrente de deslocamento é:

\[
B_d=\frac{\mu_0 I_d}{2\pi r}
\]

A razão entre eles é:

\[
\frac{B_d}{B}=\frac{I_d}{I}
\]

Substituindo:

\[
\frac{B_d}{B}
=
\frac{2,87\times10^{-16}}{100}
\]

\[
\frac{B_d}{B}
=
2,87\times10^{-18}
\]
```


\subsection{Questão 24 — Fábio Rodrigues}
\textbf{Dados}

\[
\rho = 1,62 \times 10^{-8}\,\Omega\cdot\text{m}
\]

\[
i_s = 10,0\,\text{A}
\]

\[
t_s = 50,0\,\text{ms} = 0,050\,\text{s}
\]

\[
r = 9,00\,\text{mm} = 9,00 \times 10^{-3}\,\text{m}
\]

\[
\mu_0 = 4\pi \times 10^{-7}\,\text{T}\cdot\text{m/A}
\]

\[
\varepsilon_0 = 8,85 \times 10^{-12}\,\text{F/m}
\]

A função da corrente real no tempo \(t\) (para \(t \le t_s\)) é linear:

\[
i(t) = \frac{i_s}{t_s} t = \frac{10,0}{0,050} t = 200 t
\]

A taxa de variação da corrente real para \(0 \le t \le 50,0\,\text{ms}\) é constante:

\[
\frac{di}{dt} = 200\,\text{A/s}
\]

Para \(t > 50,0\,\text{ms}\), a corrente é constante (\(10,0\,\text{A}\)) e \(\frac{di}{dt} = 0\).

O campo magnético devido à corrente real é dado por:

\[
B_i = \frac{\mu_0 i}{2\pi r}
\]

\textbf{(a) Módulo do campo \(B_i\) em \(t = 20\,\text{ms}\)}

A corrente em \(t = 0,020\,\text{s}\) é:

\[
i = 200(0,020) = 4,00\,\text{A}
\]

Substituindo na fórmula do campo:

\[
B_i = \frac{(4\pi \times 10^{-7})(4,00)}{2\pi(9,00 \times 10^{-3})} = \frac{2 \times 10^{-7} \times 4,00}{9,00 \times 10^{-3}}
\]

\[
B_i \approx 8,89 \times 10^{-5}\,\text{T}
\]

\[
\boxed{B_i = 8,89 \times 10^{-5}\,\text{T}}
\]

\textbf{(b) Módulo do campo \(B_i\) em \(t = 40\,\text{ms}\)}

A corrente em \(t = 0,040\,\text{s}\) é:

\[
i = 200(0,040) = 8,00\,\text{A}
\]

\[
B_i = \frac{2 \times 10^{-7} \times 8,00}{9,00 \times 10^{-3}}
\]

\[
B_i \approx 1,78 \times 10^{-4}\,\text{T}
\]

\[
\boxed{B_i = 1,78 \times 10^{-4}\,\text{T}}
\]

\textbf{(c) Módulo do campo \(B_i\) em \(t = 60\,\text{ms}\)}

Nesse instante, a corrente atingiu o valor máximo constante \(i = 10,0\,\text{A}\).

\[
B_i = \frac{2 \times 10^{-7} \times 10,0}{9,00 \times 10^{-3}}
\]

\[
B_i \approx 2,22 \times 10^{-4}\,\text{T}
\]

\[
\boxed{B_i = 2,22 \times 10^{-4}\,\text{T}}
\]

A corrente de deslocamento \(i_d\) é dada por:

\[
i_d = \varepsilon_0 \frac{d\Phi_E}{dt}
\]

O fluxo elétrico no fio é \(\Phi_E = E A\). Sabendo que o campo elétrico é \(E = \rho J = \rho \frac{i}{A}\), temos:

\[
\Phi_E = \left(\rho \frac{i}{A}\right) A = \rho i
\]

Portanto, a corrente de deslocamento total no fio é:

\[
i_d = \varepsilon_0 \frac{d(\rho i)}{dt} = \varepsilon_0 \rho \frac{di}{dt}
\]

O campo magnético induzido pela corrente de deslocamento é:

\[
B_{id} = \frac{\mu_0 i_d}{2\pi r} = \frac{\mu_0 \varepsilon_0 \rho}{2\pi r} \frac{di}{dt}
\]

\textbf{(d) Módulo do campo \(B_{id}\) em \(t = 20\,\text{ms}\)}

Como a taxa \(\frac{di}{dt} = 200\,\text{A/s}\), calculamos:

\[
B_{id} = \frac{(4\pi \times 10^{-7})(8,85 \times 10^{-12})(1,62 \times 10^{-8})}{2\pi(9,00 \times 10^{-3})} (200)
\]

\[
B_{id} = \frac{(2 \times 10^{-7})(1,4337 \times 10^{-19})}{9,00 \times 10^{-3}} (200)
\]

\[
B_{id} = \frac{5,7348 \times 10^{-24}}{9,00 \times 10^{-3}}
\]

\[
B_{id} \approx 6,37 \times 10^{-22}\,\text{T}
\]

\[
\boxed{B_{id} = 6,37 \times 10^{-22}\,\text{T}}
\]

\textbf{(e) Módulo do campo \(B_{id}\) em \(t = 40\,\text{ms}\)}

A taxa de variação \(\frac{di}{dt}\) continua sendo \(200\,\text{A/s}\). Logo, o campo é o mesmo:

\[
\boxed{B_{id} = 6,37 \times 10^{-22}\,\text{T}}
\]

\textbf{(f) Módulo do campo \(B_{id}\) em \(t = 60\,\text{ms}\)}

A corrente é constante, logo \(\frac{di}{dt} = 0\). A corrente de deslocamento é nula.

\[
\boxed{B_{id} = 0}
\]

\textbf{(g) Sentido de \(\vec{B}_i\) no ponto P}

Pela regra da mão direita, apontando o polegar no sentido da corrente (para a direita), os dedos indicam que o campo acima do fio aponta para fora do plano do papel.

\[
\boxed{\text{Para fora do papel}}
\]

\textbf{(h) Sentido de \(\vec{B}_{id}\) no ponto P}

Como a corrente aumenta para a direita, o campo elétrico também aponta e cresce para a direita. A corrente de deslocamento aponta no mesmo sentido. Novamente, pela regra da mão direita, o campo magnético induzido acima do fio também aponta para fora.

\[
\boxed{\text{Para fora do papel}}
\]

\subsection{Questão 25 — Mateus Valentim}


\subsubsection{Resolução}

Pela lei de Ampere-Maxwell,

\[
B(2\pi r)=\mu_0 i_{d,\text{env}}.
\]

A densidade de corrente de deslocamento é uniforme, então a corrente de
deslocamento envolvida depende da área envolvida.

\begin{enumerate}
    \item[(a)] Para \(r=2{,}00\,\text{cm}<R\),

    \[
    i_{d,\text{env}}=J_d\pi r^2.
    \]

    Portanto,

    \[
    B(2\pi r)=\mu_0J_d\pi r^2,
    \]

    ou

    \[
    B=\frac{\mu_0J_dr}{2}.
    \]

    Assim,

    \[
    B=
    \frac{(4\pi\times 10^{-7})(6{,}00)(2{,}00\times 10^{-2})}{2}
    =
    7{,}54\times 10^{-8}\,\text{T}.
    \]

    \item[(b)] Para \(r=5{,}00\,\text{cm}>R\), a corrente envolvida é a
    corrente total da região:

    \[
    i_{d,\text{env}}=J_d\pi R^2.
    \]

    Logo,

    \[
    B=
    \frac{\mu_0J_dR^2}{2r}.
    \]

    Substituindo \(R=3{,}00\,\text{cm}\) e \(r=5{,}00\,\text{cm}\),

    \[
    B=
    \frac{(4\pi\times 10^{-7})(6{,}00)(3{,}00\times 10^{-2})^2}
    {2(5{,}00\times 10^{-2})}
    =
    6{,}79\times 10^{-8}\,\text{T}.
    \]
\end{enumerate}

\subsection{Questão 26 — Cicero Rodrigues}

\textbf{26. }\normalfont Um capacitor com placas paralelas circulares, de raio R = 1,20 cm, está sendo descarregado por uma corrente de 12,0 A. Considere um anel de raio R/3, concêntrico com o capacitor, situado entre as placas. \textbf{(a)} Qual é a corrente de deslocamento envolvida pelo anel? O campo magnético máximo induzido tem módulo de 12,0 mT. A que distância radial \textbf{(b)} do lado de dentro e \textbf{(c)} do lado de fora do espaço entre as placas o módulo do campo magnético induzido é 3,00 mT?

\textbf{Resolução:}
\vspace{0.3cm}

\textbf{Dados:}
\[R = 1,20\;cm\]
\[i = i_d = 12,0\;A \quad \text{(Corrente total de deslocamento)}\]
\[B_{\max} = 12,0\;mT \quad \text{(Ocorre na borda da placa, onde } r = R \text{)}\]
\[B = 3,00\;mT \quad \text{(Campo alvo)}\]

\vspace{0.2cm}
\textbf{(a) Corrente de deslocamento envolvida pelo anel}
A corrente de deslocamento se distribui uniformemente pela área da placa. A fração envolvida em um anel de raio \(r\) é proporcional à razão das áreas:
\[i_{d,env} = i_d \left( \frac{\pi r^2}{\pi R^2} \right) = i_d \left( \frac{r}{R} \right)^2\]
Substituindo para o anel de \(r = R/3\):
\[i_{d,env} = 12,0 \left( \frac{R/3}{R} \right)^2 = 12,0 \left( \frac{1}{3} \right)^2\]
\[i_{d,env} = 12,0 \left( \frac{1}{9} \right) \approx 1,333\;A\]
\[\underline{i_{d,env} \approx 1,33\;A}\]

\vspace{0.2cm}
\textbf{(b) Distância radial do lado de dentro (\(r < R\))}
Dentro do capacitor, o campo magnético cresce linearmente com o raio em relação ao seu valor máximo na borda (\(B_{\max}\)):
\[B_{in} = B_{\max} \left( \frac{r}{R} \right)\]
Isolando \(r\) e substituindo os valores desejados:
\[3,00 = 12,0 \left( \frac{r}{1,20} \right)\]
\[\frac{r}{1,20} = \frac{3,00}{12,0} = 0,25\]
\[r = 0,25 \cdot 1,20\]
\[r = 0,300\;cm\]
\[\underline{r = 3,00\;mm}\]

\vspace{0.2cm}
\textbf{(c) Distância radial do lado de fora (\(r > R\))}
Fora do capacitor, o campo magnético decai de forma inversamente proporcional à distância radial em relação à borda:
\[B_{out} = B_{\max} \left( \frac{R}{r} \right)\]
Isolando \(r\) e substituindo os valores:
\[3,00 = 12,0 \left( \frac{1,20}{r} \right)\]
\[r = 1,20 \left( \frac{12,0}{3,00} \right)\]
\[r = 1,20 \cdot 4,00\]
\[r = 4,80\;cm\]
\[\underline{r = 48,0\;mm}\]

\subsection{Questão 27 — Guilherme Araújo Floriano}

	\textbf{(a)} O valor da corrente de deslocamento enclausurada é
\[
i_{d,\mathrm{enc}} = \int_0^r J_d\,2\pi\rho\,d\rho
= (4.00\,\mathrm{A/m^2})(2\pi)\int_0^r\left(1-\frac{\rho}{R}\right)\rho\,d\rho
= 8\pi\left(\frac{r^2}{2}-\frac{r^3}{3R}\right).
\]
Pela Eq.~32-17,
\[
B=\frac{\mu_0 i_{d,\mathrm{enc}}}{2\pi r}=27.9\,\mathrm{nT}.
\]

\textbf{(b)} Agora o limite superior da integral é $R$, então
\[
i_{d,\mathrm{enc}}=8\pi\left(\frac{R^2}{2}-\frac{R^3}{3R}\right)=\frac{4}{3}\pi R^2.
\]
Logo,
\[
B=\frac{\mu_0 i_d}{2\pi r}=15.1\,\mathrm{nT}.
\]

\subsection{Questão 28 — Arthur Roberto da Silva}
\subsection*{Dados do enunciado:}
\begin{itemize}
    \item Raio da região circular: $R = 3,00\text{ cm} = 0,03\text{ m}$
    \item Corrente de deslocamento delimitada por um raio $r \le R$: $i_d = (3,00\text{ A})\left(\frac{r}{R}\right)$
\end{itemize}

Pela Lei de Ampère-Maxwell, o campo magnético $B$ a uma distância $r$ do centro é dado por:
\begin{equation*}
B = \frac{\mu_0 i_{d,\text{env}}}{2\pi r}
\end{equation*}

\subsection*{(a) Módulo do campo magnético em $r = 2,00\text{ cm}$ (Região Interna)}

Como $r = 2,00\text{ cm} \le R$, a corrente de deslocamento envolvida pela curva amperiana é calculada diretamente pela expressão fornecida:
\begin{equation*}
i_{d,\text{env}} = 3,00 \cdot \left(\frac{2,00}{3,00}\right) = 2,00\text{ A}
\end{equation*}

Substituindo na equação do campo magnético ($r = 0,02\text{ m}$ e $\mu_0 = 4\pi \times 10^{-7}\text{ T}\cdot\text{m/A}$):
\begin{align*}
B &= \frac{(4\pi \times 10^{-7}) \cdot 2,00}{2\pi \cdot 0,02} \\
B &= \frac{2 \times 10^{-7} \cdot 2,00}{0,02} \\
B &= 2,00 \times 10^{-5}\text{ T} = 20,0\text{ }\mu\text{T}
\end{align*}

\subsection*{(b) Módulo do campo magnético em $r = 5,00\text{ cm}$ (Região Externa)}

Como $r = 5,00\text{ cm} > R$, a curva amperiana engloba \textbf{toda} a corrente de deslocamento da região circular, cujo valor máximo ocorre na borda ($r = R$):
\begin{equation*}
i_{d,\text{total}} = 3,00 \cdot \left(\frac{R}{R}\right) = 3,00\text{ A}
\end{equation*}

Substituindo na equação com a nova distância externa ($r = 0,05\text{ m}$):
\begin{align*}
B &= \frac{(4\pi \times 10^{-7}) \cdot 3,00}{2\pi \cdot 0,05} \\
B &= \frac{2 \times 10^{-7} \cdot 3,00}{0,05} \\
B &= 1,20 \times 10^{-5}\text{ T} = 12,0\text{ }\mu\text{T}
\end{align*}

\subsection*{Respostas Finais}
\begin{itemize}
    \item \textbf{(a)} $2,00 \times 10^{-5}\text{ T}$ \quad (ou $20,0\text{ }\mu\text{T}$)
    \item \textbf{(b)} $1,20 \times 10^{-5}\text{ T}$ \quad (ou $12,0\text{ }\mu\text{T}$)
\end{itemize}

\subsection{Questão 29 — Ana Alicy}
Inserir resolução aqui.

\subsection{Questão 38 — Kaylane Castro}
\subsection*{Enunciado}

Um elétron move-se em uma órbita circular de raio $r$ com
velocidade orbital $v_0$.

Um campo magnético uniforme $\vec{B}$ é aplicado perpendicularmente
ao plano da órbita.

Determine a variação do momento dipolar magnético orbital do elétron.


\subsection*{Momento dipolar magnético inicial}

O momento dipolar magnético orbital é dado por:

\[
\mu = IA
\]

A corrente equivalente produzida pelo elétron é:

\[
I=\frac{e}{T}
\]

Como o período orbital vale:

\[
T=\frac{2\pi r}{v_0}
\]

temos:

\[
I=\frac{ev_0}{2\pi r}
\]

A área da órbita circular é:

\[
A=\pi r^2
\]

Logo:

\[
\mu = \frac{ev_0}{2\pi r}\cdot\pi r^2
\]

\[
\mu = \frac{ev_0r}{2}
\]

% ---------------------------------------------------------

\subsection*{Variação do momento magnético}

Como $e$ e $r$ permanecem constantes, a variação do momento
dipolar depende apenas da variação da velocidade orbital:

\[
\Delta\mu = \frac{er}{2}\Delta v
\]

% ---------------------------------------------------------

\subsection*{Equação dinâmica}

Aplicando a segunda lei de Newton na direção radial:

\[
\frac{mv^2}{r}
=
\frac{mv_0^2}{r}+evB
\]

Reorganizando:

\[
m(v^2-v_0^2)=evBr
\]

Usando diferença de quadrados:

\[
m(v-v_0)(v+v_0)=evBr
\]

Definindo:

\[
\Delta v = v-v_0
\]

obtemos:

\[
m\Delta v (v+v_0)=evBr
\]

Como a variação da velocidade é pequena:

\[
v\approx v_0
\]

Assim:

\[
v+v_0 \approx 2v_0
\]

e também podemos aproximar:

\[
v\approx v_0
\]

no lado direito da equação.

Portanto:

\[
2m\Delta v \approx eBr
\]

Isolando $\Delta v$:

\[
\Delta v = \frac{eBr}{2m}
\]

% ---------------------------------------------------------

\subsection*{Expressão final para $\Delta\mu$}

Substituindo $\Delta v$ em

\[
\Delta\mu = \frac{er}{2}\Delta v
\]

temos:

\[
\Delta\mu
=
\frac{er}{2}
\left(
\frac{eBr}{2m}
\right)
\]

\[
\Delta\mu
=
\frac{e^2Br^2}{4m}
\]

Portanto:

\[
\boxed{
\Delta\mu
=
\frac{e^2Br^2}{4m}
}
\]

\subsection{Questão 43 — Dávyla Maria}
Inserir resolução aqui.

\subsection{Questão 44 — Rodrigo Albuquerque}

A Fig. 32-39 mostra a curva de magnetização de um material paramagnético. A escala do eixo vertical é definida por \(a = 0,15\) e a escala do eixo horizontal é definida por \(b = 0,2\ \text{T/K}\). Seja \(\mu_{\text{exp}}\) o valor experimental do momento magnético de uma amostra e \(\mu_{\text{máx}}\) o valor máximo possível do momento magnético da mesma amostra. De acordo com a lei de Curie, qual é o valor da razão \(\mu_{\text{exp}}/\mu_{\text{máx}}\) quando a amostra é submetida a um campo magnético de \(0,800\ \text{T}\) a uma temperatura de \(2,00\ \text{K}\)?}

\includegraphics[width=0.5\textwidth]{image_3.png}

\vspace{0.5cm}

Pela lei de Curie:

\[
\frac{\mu_{\text{exp}}}{\mu_{\text{máx}}}
=
\frac{B/T}{b}
\]

Dados:

\[
B=0,800\ T
\]

\[
T=2,00\ K
\]

\[
b=0,2\ \text{T/K}
\]

Calculando:

\[
\frac{B}{T}
=
\frac{0,800}{2,00}
\]

\[
\frac{B}{T}=0,400\ \text{T/K}
\]

Então:

\[
\frac{\mu_{\text{exp}}}{\mu_{\text{máx}}}
=
\frac{0,400}{0,2}
\]

\[
\frac{\mu_{\text{exp}}}{\mu_{\text{máx}}}=2,0
\]

Como o valor máximo físico possível da razão é \(1\), o material encontra-se saturado magneticamente.

Logo:

\[
\frac{\mu_{\text{exp}}}{\mu_{\text{máx}}}=1
\]

\subsection{Questão 50 — Fábio Rodrigues}
\textbf{Dados}

\[
f = 0,312\,\text{Hz}
\]

\[
B = 18,0\,\mu\text{T} = 18,0 \times 10^{-6}\,\text{T}
\]

\[
\mu = 0,680\,\text{mJ/T} = 0,680 \times 10^{-3}\,\text{J/T}
\]

A frequência de oscilação de um dipolo magnético (como a agulha de uma bússola) em um campo magnético uniforme é dada pela equação:

\[
f = \frac{1}{2\pi} \sqrt{\frac{\mu B}{I}}
\]

Elevando ambos os lados da equação ao quadrado para isolar o momento de inércia \(I\):

\[
f^2 = \frac{1}{4\pi^2} \frac{\mu B}{I}
\]

Multiplicando cruzado e isolando \(I\):

\[
I = \frac{\mu B}{4\pi^2 f^2}
\]

Substituindo os valores fornecidos na questão:

\[
I = \frac{(0,680 \times 10^{-3})(18,0 \times 10^{-6})}{4\pi^2 (0,312)^2}
\]

\[
I = \frac{12,24 \times 10^{-9}}{4\pi^2 (0,097344)}
\]

\[
I = \frac{12,24 \times 10^{-9}}{3,843}
\]

\[
I \approx 3,1849 \times 10^{-9}\,\text{kg}\cdot\text{m}^2
\]

\[
\boxed{I = 3,18 \times 10^{-9}\,\text{kg}\cdot\text{m}^2}
\]

\subsection{Questão 51 — Mateus Valentim}
\subsubsection{Resolução}

O raio médio do anel é

\[
r_m=\frac{5{,}0\,\text{cm}+6{,}0\,\text{cm}}{2}
=5{,}5\,\text{cm}
=5{,}5\times 10^{-2}\,\text{m}.
\]

\begin{enumerate}
    \item[(a)] Para um toroide,

    \[
    B_0=\frac{\mu_0Ni}{2\pi r_m}.
    \]

    Assim,

    \[
    i=\frac{B_0(2\pi r_m)}{\mu_0N}.
    \]

    Substituindo os valores,

    \[
    i=
    \frac{(0{,}20\times 10^{-3})(2\pi)(5{,}5\times 10^{-2})}
    {(4\pi\times 10^{-7})(400)}
    =
    1{,}38\times 10^{-1}\,\text{A}.
    \]

    Portanto,

    \[
    i=0{,}14\,\text{A}.
    \]

    \item[(b)] A seção reta circular do anel tem raio

    \[
    a=\frac{6{,}0\,\text{cm}-5{,}0\,\text{cm}}{2}
    =0{,}50\,\text{cm}
    =5{,}0\times 10^{-3}\,\text{m}.
    \]

    A área da seção reta é

    \[
    A=\pi a^2
    =\pi(5{,}0\times 10^{-3})^2
    =7{,}85\times 10^{-5}\,\text{m}^2.
    \]

    O campo total no material é

    \[
    B=B_0+B_M=B_0+800B_0=801B_0.
    \]

    Quando a corrente começa a circular, a variação do fluxo na bobina
    secundária é

    \[
    \Delta \Phi=B A.
    \]

    Como

    \[
    \mathcal{E}=-N_s\frac{d\Phi}{dt}
    \]

    e

    \[
    i_s=\frac{\mathcal{E}}{R},
    \]

    a carga total que atravessa a bobina secundária é

    \[
    q=\int i_s\,dt
    =
    \frac{N_s\Delta\Phi}{R}
    =
    \frac{N_sBA}{R}.
    \]

    Logo,

    \[
    q=
    \frac{(50)(801)(0{,}20\times 10^{-3})(7{,}85\times 10^{-5})}{8{,}0}
    =
    7{,}86\times 10^{-5}\,\text{C}.
    \]
\end{enumerate}

\subsection{Questão 52 — Cicero Rodrigues}

\textbf{52. }\normalfont Medidas realizadas em minas e poços revelam que a temperatura no interior da Terra aumenta com a profundidade à taxa média de 30$^\circ$C/km. Supondo que a temperatura na superfície seja 10$^\circ$C, a que profundidade o ferro deixa de ser ferromagnético? (A temperatura de Curie do ferro varia muito pouco com a pressão.)

\textbf{Resolução:}
\vspace{0.3cm}

\textbf{Dados:}
\[\frac{dT}{dy} = 30^\circ\text{C/km} \quad \text{(Taxa de aquecimento)}\]
\[T_0 = 10^\circ\text{C} \quad \text{(Temperatura na superfície)}\]
\[T_C = 770^\circ\text{C} \quad \text{(Temperatura de Curie do ferro)}\]

\vspace{0.2cm}
\textbf{Cálculo da profundidade (\(d\))}
A temperatura no interior da Terra em função da profundidade \(d\) (em km) pode ser modelada por uma função de primeiro grau:
\[T(d) = T_0 + \left( \frac{dT}{dy} \right) d\]
Queremos descobrir a profundidade \(d\) em que a temperatura atinge \(T_C\). Substituindo os valores:
\[770 = 10 + 30d\]
Isolando o termo com a profundidade:
\[30d = 770 - 10\]
\[30d = 760\]
\[d = \frac{760}{30}\]
\[d \approx 25,333\;km\]
\[\underline{d \approx 25,3\;km}\]

\subsection{Questão 53 — Guilherme Araújo Floriano}

	\textbf{(a)} Se a magnetização da esfera está saturada, o momento dipolar total é
\[
\mu_{\mathrm{total}}=N\mu,
\]
onde $N$ é o número de átomos de ferro e $\mu$ é o momento dipolar de um átomo. Como a massa da esfera é $Nm$ e também $\frac{4\pi\rho R^3}{3}$,
\[
N=\frac{4\pi\rho R^3}{3m}.
\]
Substituindo,
\[
\mu_{\mathrm{total}}=\frac{4\pi\rho R^3\mu}{3m}
\quad\Rightarrow\quad
R=\left(\frac{3m\mu_{\mathrm{total}}}{4\pi\rho\mu}\right)^{1/3}.
\]
Com $m=56u=9.30\times10^{-26}\,\mathrm{kg}$, obtém-se
\[
R=1.8\times10^5\,\mathrm{m}.
\]

	\textbf{(b)} O volume da esfera é
\[
V_s=\frac{4\pi R^3}{3}=2.53\times10^{16}\,\mathrm{m^3}.
\]
Comparando com o volume da Terra,
\[
V_T=1.08\times10^{21}\,\mathrm{m^3},
\]
temos a fração
\[
\frac{V_s}{V_T}=2.3\times10^{-5}.
\]

\subsection{Questão 54 — Arthur Roberto da Silva}
\subsection*{Dados do enunciado:}
\begin{itemize}
    \item Raio das placas: $R = 4,00\text{ cm} = 0,04\text{ m}$
    \item Ponto interno ($r_1 < R$): $r_1 = 2,00\text{ cm} = 0,02\text{ m}$
    \item Campo no ponto interno: $B_1 = 12,5\text{ nT} = 12,5 \times 10^{-9}\text{ T}$
    \item Ponto externo ($r_2 > R$): $r_2 = 6,00\text{ cm} = 0,06\text{ m}$
\end{itemize}

\subsection*{(a) Módulo do campo magnético a $6,00\text{ cm}$ (ponto externo)}

Pela Lei de Ampère-Maxwell, as expressões para os campos magnéticos interno e externo são:
\begin{equation*}
B_{\text{int}} = \frac{\mu_0 i}{2\pi} \frac{r}{R^2} \quad \text{e} \quad B_{\text{ext}} = \frac{\mu_0 i}{2\pi r}
\end{equation*}

Fazendo a razão entre o campo externo desejado ($B_2$) e o campo interno conhecido ($B_1$):
\begin{equation*}
\frac{B_2}{B_1} = \frac{\frac{\mu_0 i}{2\pi r_2}}{\frac{\mu_0 i \cdot r_1}{2\pi R^2}} = \frac{R^2}{r_1 \cdot r_2}
\end{equation*}

Isolando $B_2$ e substituindo os valores:
\begin{align*}
B_2 &= B_1 \cdot \frac{R^2}{r_1 \cdot r_2} \\
B_2 &= 12,5\text{ nT} \cdot \frac{(4,00)^2}{2,00 \cdot 6,00} \\
B_2 &= 12,5\text{ nT} \cdot \frac{16,00}{12,00} \\
B_2 &\approx 16,7\text{ nT}
\end{align*}

\subsection*{(b) Corrente nos fios ligados às placas}

Utilizando a fórmula do campo magnético na região interna ($B_1$):
\begin{equation*}
B_1 = \frac{\mu_0 i}{2\pi} \frac{r_1}{R^2}
\end{equation*}

Isolando a corrente $i$:
\begin{equation*}
i = \frac{2\pi \cdot B_1 \cdot R^2}{\mu_0 \cdot r_1}
\end{equation*}

Substituindo os valores no Sistema Internacional ($\mu_0 = 4\pi \times 10^{-7}\text{ T}\cdot\text{m/A}$):
\begin{align*}
i &= \frac{2\pi \cdot (12,5 \times 10^{-9}) \cdot (0,04)^2}{(4\pi \times 10^{-7}) \cdot 0,02} \\
i &= \frac{(12,5 \times 10^{-9}) \cdot 0,0016}{2 \times 10^{-7} \cdot 0,02} \\
i &= \frac{2,0 \times 10^{-11}}{4,0 \times 10^{-9}} \\
i &= 0,005\text{ A} = 5,00\text{ mA}
\end{align*}

\subsection*{Respostas Finais}
\begin{itemize}
    \item \textbf{(a)} $16,7\text{ nT}$
    \item \textbf{(b)} $5,00\text{ mA}$
\end{itemize}

\subsection{Questão 56 — Guilherme Araújo Floriano}

\textbf{(a)} No equador magnético ($\lambda=0$),
\[
B=\frac{\mu_0\mu}{4\pi r^3}=3.10\times10^{-5}\,\mathrm{T}.
\]

	\textbf{(b)}
\[
\phi_i=\tan^{-1}(2\tan\lambda)=\tan^{-1}(0)=0^\circ.
\]

	\textbf{(c)} Para $\lambda=60^\circ$,
\[
B=\frac{\mu_0\mu}{4\pi r^3}(1+3\sin^2\lambda)=5.59\times10^{-5}\,\mathrm{T}.
\]

	\textbf{(d)}
\[
\phi_i=\tan^{-1}(2\tan 60^\circ)=73.9^\circ.
\]

	\textbf{(e)} No polo magnético norte ($\lambda=90^\circ$),
\[
B=\frac{\mu_0\mu}{4\pi r^3}(1+3\sin^2 90^\circ)=6.20\times10^{-5}\,\mathrm{T}.
\]

	\textbf{(f)}
\[
\phi_i=\tan^{-1}(2\tan 90^\circ)=90^\circ.
\]

\subsection{Questão 57 — Cicero Rodrigues}

\textbf{57. }\normalfont A agulha de uma bússola, com 0,050 kg de massa e 4,0 cm de comprimento, está alinhada com a componente horizontal do campo magnético da Terra em um local em que a componente tem o valor $B_h = 16\;\mu$T. Depois que a agulha recebe um leve empurrão, ela começa a oscilar com uma frequência angular $\omega = 45$ rad/s. Supondo que a agulha seja uma barra fina e uniforme, livre para girar em torno do centro, determine o módulo do momento dipolar magnético da agulha.

\textbf{Resolução:}
\vspace{0.3cm}

\textbf{Dados:}
\[m = 0,050\;kg\]
\[L = 4,0\;cm = 0,040\;m\]
\[B_h = 16\,\mu T = 16 \times 10^{-6}\;T\]
\[\omega = 45\;rad/s\]

\vspace{0.2cm}
\textbf{(1) Momento de inércia da agulha (\(I\))}
Considerando a agulha como uma barra fina e uniforme livre para girar em torno de seu centro de massa, o momento de inércia é dado por:
\[I = \frac{1}{12}mL^2\]
Substituindo os valores numéricos:
\[I = \frac{1}{12}(0,050)(0,040)^2\]
\[I = \frac{1}{12}(0,050)(0,0016)\]
\[I = \frac{0,000080}{12}\]
\[I \approx 6,667 \times 10^{-6}\;kg\cdot m^2\]

\vspace{0.2cm}
\textbf{(2) Momento dipolar magnético (\(\mu\))}
A agulha se comporta como um pêndulo físico sujeito ao torque magnético restaurador. A frequência angular de pequenas oscilações é:
\[\omega = \sqrt{\frac{\mu B_h}{I}}\]
Elevando ambos os lados ao quadrado para isolar \(\mu\):
\[\omega^2 = \frac{\mu B_h}{I} \Rightarrow \mu = \frac{I \omega^2}{B_h}\]
Substituindo os valores:
\[\mu = \frac{(6,667 \times 10^{-6})(45)^2}{16 \times 10^{-6}}\]
Podemos cancelar a potência de \(10^{-6}\) no numerador e denominador:
\[\mu = \frac{(6,667)(2025)}{16}\]
\[\mu = \frac{13500}{16}\]
\[\mu = 843,75\;A\cdot m^2\]
Arredondando para os dois algarismos significativos dados no problema (\(16\,\mu T\), \(45\;rad/s\)):
\[\underline{\mu \approx 8,4 \times 10^2\;A\cdot m^2}\]

\subsection{Questão 58 — Mateus Valentim}
Inserir resolução aqui.

\subsection{Questão 59 — Fábio Rodrigues}
\textbf{Dados}

\[
R = 55,0\,\text{mm} = 0,0550\,\text{m}
\]

O campo magnético induzido \(B\) no interior e no exterior de um capacitor de placas circulares paralelas em processo de carga é dado pelas seguintes expressões (onde \(i_d\) é a corrente de deslocamento):

\begin{itemize}
    \item Para \(r \le R\) (dentro do capacitor): \(B_{\text{in}} = \frac{\mu_0 i_d r}{2\pi R^2}\)
    \item Para \(r \ge R\) (fora do capacitor): \(B_{\text{out}} = \frac{\mu_0 i_d}{2\pi r}\)
\end{itemize}

O valor máximo do campo magnético ocorre exatamente na borda das placas, onde \(r = R\). Substituindo \(r = R\) em qualquer uma das equações acima, obtemos:

\[
B_{\text{max}} = \frac{\mu_0 i_d}{2\pi R}
\]

A questão pede a distância \(r\) onde o campo é igual a \(50,0\%\) do valor máximo, ou seja, \(B = 0,500 B_{\text{max}}\).

\textbf{(a) Distância dentro do espaço entre as placas (\(r < R\))}

Igualamos a expressão do campo interno a \(0,500 B_{\text{max}}\):

\[
B_{\text{in}} = 0,500 B_{\text{max}}
\]

\[
\frac{\mu_0 i_d r}{2\pi R^2} = 0,500 \left( \frac{\mu_0 i_d}{2\pi R} \right)
\]

Cancelando os termos comuns (\(\frac{\mu_0 i_d}{2\pi}\)):

\[
\frac{r}{R^2} = \frac{0,500}{R}
\]

\[
r = 0,500 R
\]

Substituindo o valor do raio \(R = 55,0\,\text{mm}\):

\[
r = 0,500(55,0)
\]

\[
\boxed{r = 27,5\,\text{mm}}
\]

\textbf{(b) Distância fora do espaço entre as placas (\(r > R\))}

Igualamos a expressão do campo externo a \(0,500 B_{\text{max}}\):

\[
B_{\text{out}} = 0,500 B_{\text{max}}
\]

\[
\frac{\mu_0 i_d}{2\pi r} = 0,500 \left( \frac{\mu_0 i_d}{2\pi R} \right)
\]

Cancelando os termos comuns (\(\frac{\mu_0 i_d}{2\pi}\)):

\[
\frac{1}{r} = \frac{0,500}{R}
\]

\[
r = \frac{R}{0,500}
\]

\[
r = 2R
\]

Substituindo o valor do raio \(R = 55,0\,\text{mm}\):

\[
r = 2(55,0)
\]

\[
\boxed{r = 110\,\text{mm}}
\]

\subsection{Questão 60 — Rodrigo Albuquerque}

Um anel delgado de raio \(r\) possui uma carga \(q\) uniformemente distribuída. O anel gira com velocidade angular constante \(\omega\) em torno do seu eixo central.}

\textbf{(a) Mostre que o momento magnético do anel é dado por}

\[
\mu=\frac{q\omega r^2}{2}
\]

\textbf{(b) Qual é a orientação desse momento magnético se a carga é positiva?}

\vspace{0.5cm}

A carga em movimento equivale a uma corrente elétrica.

O período de rotação do anel é:

\[
T=\frac{2\pi}{\omega}
\]

A corrente elétrica associada é:

\[
i=\frac{q}{T}
\]

Substituindo \(T\):

\[
i=\frac{q}{2\pi/\omega}
\]

\[
i=\frac{q\omega}{2\pi}
\]

O momento magnético de uma espira circular é:

\[
\mu=iA
\]

A área do anel é:

\[
A=\pi r^2
\]

Então:

\[
\mu=\left(\frac{q\omega}{2\pi}\right)(\pi r^2)
\]

Simplificando:

\[
\mu=\frac{q\omega r^2}{2}
\]

\vspace{0.5cm}

\textbf{(b)}

Se a carga é positiva, o momento magnético possui a direção dada pela regra da mão direita.

Assim, ao enrolar os dedos no sentido da rotação da carga, o polegar aponta a direção do momento magnético.
```


\subsection{Questão 61 — Dávyla Maria}
Inserir resolução aqui.

\subsection{Questão 62 — Kaylane Castro}
\subsection*{Enunciado}

O capacitor da Fig. 32-7 está sendo carregado com uma corrente de

\[
i = 2,50\,\text{A}
\]

O raio do fio é:

\[
R_{\text{fio}} = 1,50\,\text{mm}
\]

e o raio das placas do capacitor é:

\[
R_{\text{placa}} = 2,00\,\text{cm}
\]

Suponha que as distribuições da corrente de condução e da corrente de
deslocamento sejam uniformes.

Determine o módulo do campo magnético:

\begin{itemize}
    \item[(a)] a $1,00\,\text{mm}$ do eixo do fio;
    \item[(b)] a $3,00\,\text{mm}$ do eixo do fio;
    \item[(c)] a $2,20\,\text{cm}$ do eixo do fio;
    \item[(d)] a $1,00\,\text{mm}$ do eixo entre as placas;
    \item[(e)] a $3,00\,\text{mm}$ do eixo entre as placas;
    \item[(f)] a $2,20\,\text{cm}$ do eixo entre as placas.
\end{itemize}

% ---------------------------------------------------------

\subsection*{Dados}

\[
i=i_d=2,50\,\text{A}
\]

\[
R_{\text{fio}} = 1,50\times10^{-3}\,\text{m}
\]

\[
R_{\text{placa}} = 2,00\times10^{-2}\,\text{m}
\]

\[
\mu_0 = 4\pi\times10^{-7}\,\text{T}\cdot\text{m/A}
\]

% =========================================================
\subsection*{Campo magnético produzido pelo fio}
% =========================================================

\subsubsection*{(a) Para $r = 1,00\,\text{mm}$}

Como:

\[
r < R_{\text{fio}}
\]

usa-se a expressão do campo no interior do fio:

\[
B =
\frac{\mu_0 i r}
{2\pi R_{\text{fio}}^2}
\]

Substituindo os valores:

\[
B =
\frac{
(4\pi\times10^{-7})(2,50)(1,00\times10^{-3})
}
{
2\pi(1,50\times10^{-3})^2
}
\]

\[
B \approx 2,22\times10^{-4}\,\text{T}
\]

Portanto:

\[
\boxed{
B \approx 2,2\times10^{-4}\,\text{T}
}
\]

% ---------------------------------------------------------

\subsubsection*{(b) Para $r = 3,00\,\text{mm}$}

Agora:

\[
r > R_{\text{fio}}
\]

Logo, toda a corrente é envolvida:

\[
B =
\frac{\mu_0 i}{2\pi r}
\]

Substituindo:

\[
B =
\frac{
(4\pi\times10^{-7})(2,50)
}
{
2\pi(3,00\times10^{-3})
}
\]

\[
B \approx 1,67\times10^{-4}\,\text{T}
\]

Portanto:

\[
\boxed{
B \approx 1,7\times10^{-4}\,\text{T}
}
\]

% ---------------------------------------------------------

\subsubsection*{(c) Para $r = 2,20\,\text{cm}$}

\[
r = 2,20\times10^{-2}\,\text{m}
\]

Usando novamente:

\[
B =
\frac{\mu_0 i}{2\pi r}
\]

\[
B =
\frac{
(4\pi\times10^{-7})(2,50)
}
{
2\pi(2,20\times10^{-2})
}
\]

\[
B \approx 2,27\times10^{-5}\,\text{T}
\]

Portanto:

\[
\boxed{
B \approx 2,3\times10^{-5}\,\text{T}
}
\]

% =========================================================
\subsection*{Campo magnético produzido pela corrente de deslocamento}
% =========================================================

\subsubsection*{(d) Para $r = 1,00\,\text{mm}$}

Como:

\[
r < R_{\text{placa}}
\]

a corrente envolvida é proporcional à área:

\[
i_{d,\text{env}}
=
i_d
\frac{r^2}{R_{\text{placa}}^2}
\]

Assim:

\[
B =
\frac{\mu_0 i_d r}
{2\pi R_{\text{placa}}^2}
\]

Substituindo os valores:

\[
B =
\frac{
(4\pi\times10^{-7})(2,50)(1,00\times10^{-3})
}
{
2\pi(2,00\times10^{-2})^2
}
\]

\[
B = 1,25\times10^{-6}\,\text{T}
\]

Portanto:

\[
\boxed{
B \approx 1,3\times10^{-6}\,\text{T}
}
\]

% ---------------------------------------------------------

\subsubsection*{(e) Para $r = 3,00\,\text{mm}$}

\[
B =
\frac{
(4\pi\times10^{-7})(2,50)(3,00\times10^{-3})
}
{
2\pi(2,00\times10^{-2})^2
}
\]

\[
B = 3,75\times10^{-6}\,\text{T}
\]

Portanto:

\[
\boxed{
B \approx 3,8\times10^{-6}\,\text{T}
}
\]

% ---------------------------------------------------------

\subsubsection*{(f) Para $r = 2,20\,\text{cm}$}

Como:

\[
r > R_{\text{placa}}
\]

toda a corrente de deslocamento é envolvida:

\[
B =
\frac{\mu_0 i_d}{2\pi r}
\]

Substituindo:

\[
B =
\frac{
(4\pi\times10^{-7})(2,50)
}
{
2\pi(2,20\times10^{-2})
}
\]

\[
B \approx 2,27\times10^{-5}\,\text{T}
\]

Portanto:

\[
\boxed{
B \approx 2,3\times10^{-5}\,\text{T}
}
\]

% =========================================================
\subsection*{Explicação final}
% =========================================================

Nas menores distâncias, os campos são diferentes porque a corrente no
fio está concentrada em uma região muito pequena, enquanto a corrente
de deslocamento está distribuída por toda a área das placas.

Assim, perto do eixo, a fração da corrente envolvida entre as placas é
muito menor.

Já para a distância maior:

\[
r > R_{\text{fio}}
\qquad \text{e} \qquad
r > R_{\text{placa}}
\]

toda a corrente é envolvida nos dois casos.

Por isso, os campos magnéticos possuem praticamente o mesmo valor.


\subsection{Questão 63 — Ana Alicy}
Inserir resolução aqui.

\subsection{Questão 64 — Ana Alicy}
Inserir resolução aqui.

\subsection{Questão 65 — Kaylane Castro}

\subsection*{Enunciado}

A agulha de uma bússola, com massa

\[
m = 0,050\,\text{kg}
\]

e comprimento

\[
L = 4,0\,\text{cm}
\]

está alinhada com a componente horizontal do campo magnético da Terra,
cujo módulo é

\[
B_h = 16\,\mu\text{T}
\]

Após receber um pequeno impulso, a agulha passa a oscilar com frequência angular

\[
\omega = 45\,\text{rad/s}
\]

Supondo que a agulha seja uma barra fina e uniforme livre para girar
em torno do centro, determine o módulo do momento dipolar magnético
da agulha.


\subsection*{Dados}

\[
m = 0,050\,\text{kg}
\]

\[
L = 4,0\,\text{cm} = 0,040\,\text{m}
\]

\[
B_h = 16\,\mu\text{T} = 1,6\times10^{-5}\,\text{T}
\]

\[
\omega = 45\,\text{rad/s}
\]


\subsection*{Momento de inércia da agulha}

A agulha é modelada como uma barra fina uniforme girando em torno do centro.

Assim, o momento de inércia é:

\[
I = \frac{mL^2}{12}
\]

Substituindo os valores:

\[
I =
\frac{
(0,050)(0,040)^2
}{12}
\]

\[
I =
\frac{
(0,050)(1,6\times10^{-3})
}{12}
\]

\[
I =
\frac{
8,0\times10^{-5}
}{12}
\]

\[
I \approx 6,67\times10^{-6}\,\text{kg}\cdot\text{m}^2
\]


\subsection*{Equação do movimento}

O torque magnético sobre a agulha é:

\[
\tau = -\mu B_h \sin\theta
\]

Para pequenas oscilações:

\[
\sin\theta \approx \theta
\]

Logo:

\[
\tau \approx -\mu B_h \theta
\]

Usando a relação dinâmica rotacional:

\[
\tau = I\alpha = I\ddot{\theta}
\]

temos:

\[
I\ddot{\theta} + \mu B_h \theta = 0
\]

Essa é a equação do movimento harmônico simples, cuja frequência angular é:

\[
\omega^2 = \frac{\mu B_h}{I}
\]

Isolando o momento dipolar magnético:

\[
\mu = \frac{\omega^2 I}{B_h}
\]

% ---------------------------------------------------------

\subsection*{Cálculo do momento dipolar magnético}

Substituindo os valores:

\[
\mu =
\frac{
(45)^2(6,67\times10^{-6})
}{
1,6\times10^{-5}
}
\]

\[
\mu =
\frac{
2025(6,67\times10^{-6})
}{
1,6\times10^{-5}
}
\]

\[
\mu =
\frac{
1,35\times10^{-2}
}{
1,6\times10^{-5}
}
\]

\[
\mu \approx 8,44\times10^{2}\,\text{A}\cdot\text{m}^2
\]

% ---------------------------------------------------------

\subsection*{Resposta}

\[
\boxed{
\mu \approx 8,4\times10^{2}\,\text{A}\cdot\text{m}^2
}
\]


\subsection{Questão 66 — Dávyla Maria}
Inserir resolução aqui.

\subsection{Questão 67 — Rodrigo Albuquerque}

\textbf{Na Fig. 32-40 um capacitor de placas paralelas está sendo descarregado por uma corrente}

\[
i=5,0\ A
\]

\textbf{As placas são quadrados de lado}

\[
L=8,0\ mm
\]

\textbf{(a) Qual é a taxa de variação do campo elétrico entre as placas?}

\textbf{(b) Qual é o valor de}

\[
\oint \vec{B}\cdot d\vec{s}
\]

\textbf{ao longo da linha tracejada, na qual}

\[
H=2,0\ mm
\]

\[
W=3,0\ mm
\]

\includegraphics[width=0.5\textwidth]{image_4.png}

\vspace{0.5cm}

\textbf{(a)}

Pela corrente de deslocamento:

\[
i_d=\varepsilon_0\frac{d\Phi_E}{dt}
\]

Como o campo elétrico é uniforme:

\[
\Phi_E=EA
\]

Então:

\[
i=\varepsilon_0A\frac{dE}{dt}
\]

A área das placas é:

\[
A=L^2
\]

\[
A=(8,0\times10^{-3})^2
\]

\[
A=6,4\times10^{-5}\ m^2
\]

Logo:

\[
\frac{dE}{dt}
=
\frac{i}{\varepsilon_0A}
\]

Substituindo:

\[
\frac{dE}{dt}
=
\frac{5,0}
{(8,85\times10^{-12})(6,4\times10^{-5})}
\]

\[
\frac{dE}{dt}
\approx8,83\times10^{15}\ V/(m\cdot s)
\]

\vspace{0.5cm}

\textbf{(b)}

Pela lei de Ampère-Maxwell:

\[
\oint \vec B\cdot d\vec s
=
\mu_0\varepsilon_0
\frac{d\Phi_E}{dt}
\]

Como:

\[
\frac{d\Phi_E}{dt}
=
\frac{dE}{dt}(HW)
\]

temos:

\[
\oint \vec B\cdot d\vec s
=
\mu_0\varepsilon_0
(HW)
\frac{dE}{dt}
\]

Mas:

\[
\varepsilon_0\frac{dE}{dt}
=
\frac{i}{L^2}
\]

Então:

\[
\oint \vec B\cdot d\vec s
=
\mu_0 i\frac{HW}{L^2}
\]

Substituindo:

\[
H=2,0\times10^{-3}\ m
\]

\[
W=3,0\times10^{-3}\ m
\]

\[
L=8,0\times10^{-3}\ m
\]

\[
\oint \vec B\cdot d\vec s
=
(4\pi\times10^{-7})(5,0)
\frac{(2,0\times10^{-3})(3,0\times10^{-3})}
{(8,0\times10^{-3})^2}
\]

\[
\oint \vec B\cdot d\vec s
\approx5,89\times10^{-7}\ T\cdot m
\]
```


\subsection{Questão 68 — Fábio Rodrigues}
\textbf{Dados}

\[
\Phi_{\text{inf}} = +7,0\,\text{mWb} = +7,0 \times 10^{-3}\,\text{Wb} \quad \text{(direcionado para fora, logo, fluxo positivo)}
\]

\[
r = 4,2\,\text{cm} = 0,042\,\text{m}
\]

\[
B_{\text{sup}} = 0,40\,\text{T} \quad \text{(perpendicular e para cima, ou seja, para fora da superfície fechada, logo, fluxo positivo)}
\]

De acordo com a Lei de Gauss para o magnetismo, o fluxo magnético total através de qualquer superfície fechada é sempre nulo:

\[
\Phi_{\text{total}} = \Phi_{\text{inf}} + \Phi_{\text{sup}} + \Phi_{\text{curva}} = 0
\]

Primeiro, calculamos o fluxo na face superior plana (\(\Phi_{\text{sup}}\)). A área dessa face é:

\[
A_{\text{sup}} = \pi r^2 = \pi (0,042)^2 \approx 5,54 \times 10^{-3}\,\text{m}^2
\]

O fluxo através dessa face superior é:

\[
\Phi_{\text{sup}} = B_{\text{sup}} A_{\text{sup}} \cos(0^\circ)
\]

\[
\Phi_{\text{sup}} = (0,40)(5,54 \times 10^{-3})(1)
\]

\[
\Phi_{\text{sup}} \approx 2,22 \times 10^{-3}\,\text{Wb} = +2,22\,\text{mWb}
\]

\textbf{(a) Valor absoluto do fluxo na parte curva}

Substituindo os valores conhecidos na equação da Lei de Gauss:

\[
7,0\,\text{mWb} + 2,22\,\text{mWb} + \Phi_{\text{curva}} = 0
\]

\[
9,22\,\text{mWb} + \Phi_{\text{curva}} = 0
\]

\[
\Phi_{\text{curva}} = -9,22\,\text{mWb}
\]

Considerando os algarismos significativos dos dados fornecidos, o valor absoluto do fluxo é:

\[
|\Phi_{\text{curva}}| = 9,2\,\text{mWb}
\]

\[
\boxed{|\Phi_{\text{curva}}| = 9,2\,\text{mWb}}
\]

\textbf{(b) Sentido do fluxo magnético}

Por convenção, um fluxo magnético com sinal negativo indica que as linhas de campo estão entrando na superfície fechada. Portanto, o sentido é para dentro.

\[
\boxed{\text{Para dentro}}
\]

\subsection{Questão 69 — Mateus Valentim}
\subsubsection{Resolução}

Entre as placas,

\[
E=\frac{V}{d}.
\]

Como

\[
V=(100\,\text{V})e^{-t/\tau},
\]

temos

\[
E=\frac{100}{d}e^{-t/\tau}.
\]

Pela lei de Ampere-Maxwell,

\[
B(2\pi r)=\mu_0\epsilon_0\frac{d\Phi_E}{dt}.
\]

Para \(r<R\), o fluxo envolvido é

\[
\Phi_E=E(\pi r^2).
\]

Logo,

\[
B(2\pi r)=\mu_0\epsilon_0\pi r^2
\left|\frac{dE}{dt}\right|.
\]

Assim,

\[
B=\frac{\mu_0\epsilon_0 r}{2}
\left|\frac{dE}{dt}\right|.
\]

Como

\[
\left|\frac{dE}{dt}\right|
=
\frac{100}{d\tau}e^{-t/\tau},
\]

obtemos

\[
B(t)=
\frac{\mu_0\epsilon_0 r}{2}
\frac{100}{d\tau}e^{-t/\tau}.
\]

\begin{enumerate}
    \item[(a)] Como \(r=0{,}80R=0{,}80(16\,\text{mm})=12{,}8\,\text{mm}\),

    \[
    B(t)=
    \frac{(4\pi\times 10^{-7})(8{,}85\times 10^{-12})
    (12{,}8\times 10^{-3})}{2}
    \frac{100}{(5{,}0\times 10^{-3})(12\times 10^{-3})}
    e^{-t/\tau}.
    \]

    Portanto,

    \[
    B(t)=
    (1{,}19\times 10^{-13})e^{-t/\tau}\,\text{T}.
    \]

    \item[(b)] Para \(t=3\tau\),

    \[
    B(3\tau)=
    (1{,}19\times 10^{-13})e^{-3}
    =
    5{,}91\times 10^{-15}\,\text{T}.
    \]
\end{enumerate}

\subsection{Questão 70 — Cicero Rodrigues}

\textbf{70. }\normalfont Paramagnetismo, diamagnetismo. A Fig. mostra um dispositivo usado em demonstrações em sala de aula de paramagnetismo e diamagnetismo. Uma amostra de um material é pendurada em uma corda de comprimento $L = 2$ m em uma região de largura $d = 2$ cm entre os polos de um eletroímã. Como mostra a figura, o polo P1 tem forma de cunha e o polo de P2 é côncavo. Os desvios da corda em relação à vertical são mostrados à plateia por meio de um sistema de projeção que não é mostrado na figura. \textbf{(a)} Primeiro, uma amostra de bismuto (um elemento diamagnético, mau condutor de eletricidade) é usada. Quando o eletroímã é ligado, a amostra se desloca ligeiramente (cerca de 1 mm) em direção a um dos polos. Qual é o sentido desse movimento? \textbf{(b)} Em seguida, uma amostra de alumínio (um elemento paramagnético, bom condutor de eletricidade) é usada. A amostra sofre um grande deslocamento (cerca de 1 cm) em direção a um dos polos durante cerca de um segundo e depois um deslocamento moderado (de alguns milímetros) em direção ao outro polo. Explique por que isso acontece e indique o sentido das deflexões. \textbf{(c)} O que aconteceria com uma amostra de um material ferromagnético?

\includegraphics[width=0.4\textwidth]{image_12.png}


\textbf{Resolução:}
\vspace{0.3cm}

\textbf{Análise Preliminar do Campo Magnético:}
A geometria dos polos cria um campo magnético \textit{não uniforme}. O polo P1 (forma de cunha) concentra as linhas de campo, tornando a intensidade do campo magnético muito maior perto dele. O polo P2 (côncavo) espalha as linhas, criando uma região de campo mais fraco. 

\vspace{0.2cm}
\textbf{(a) Amostra de bismuto (diamagnético)}
Materiais diamagnéticos desenvolvem um momento de dipolo induzido que se opõe ao campo magnético externo. Em um campo não uniforme, isso resulta em uma força resultante que empurra o material para a região de \textit{menor} intensidade de campo.
\[\underline{\text{Sentido: Afasta-se de P1, em direção ao polo P2.}}\]

\vspace{0.2cm}
\textbf{(b) Amostra de alumínio (paramagnético e condutor)}
Temos dois fenômenos ocorrendo em instantes diferentes:
\begin{itemize}
    \item \textbf{Fase 1 (Transiente - durante $\approx 1$ segundo):} Quando o eletroímã é ligado, o fluxo magnético aumenta rapidamente. Como o alumínio é um bom condutor, a Lei de Faraday-Lenz nos diz que correntes de Foucault serão induzidas na amostra. O campo magnético gerado por essas correntes repele fortemente a amostra da região de maior variação de fluxo (perto de P1). Isso causa o grande deslocamento inicial para longe da cunha.
    \item \textbf{Fase 2 (Regime permanente):} Quando o campo se estabiliza, não há mais variação de fluxo e as correntes induzidas cessam. Resta apenas o efeito paramagnético natural do alumínio. Materiais paramagnéticos são fracamente atraídos para regiões de \textit{maior} intensidade de campo magnético.
\end{itemize}
\[\underline{\text{Sentido: Primeiro sofre forte repulsão em direção a P2,}}\]
\[\underline{\text{depois é moderadamente atraído em direção a P1.}}\]

\vspace{0.2cm}
\textbf{(c) Amostra de material ferromagnético}
Materiais ferromagnéticos têm seus domínios magnéticos fortemente alinhados com o campo externo, gerando uma atração magnética extremamente intensa (ordens de grandeza maior que o paramagnetismo) para regiões de campo magnético mais forte.
\[\underline{\text{A amostra seria puxada violentamente em direção ao polo P1 (cunha),}}\]
\[\underline{\text{muito provavelmente colidindo e ficando grudada nele.}}\]
 

\subsection{Questão 71 — Guilherme Araújo Floriano}

Se a área de cada placa circular é $A$ e a seção circular central é $a$, então
\[
\frac{A}{a}=4.
\]
Pelas Eqs.~32-14 e 32-15, a corrente total de deslocamento é
\[
i_d=4(2.0\,\mathrm{A})=8.0\,\mathrm{A}.
\]


\subsection{Questão 72 — Arthur Roberto da Silva}
A corrente de deslocamento é dada pela Lei de Ampère-Maxwell:
\begin{equation*}
i_d = \varepsilon_0 \frac{d\Phi_E}{dt}
\end{equation*}

Como o campo elétrico $E$ é perpendicular à área constante $A$, o fluxo elétrico é $\Phi_E = E \cdot A$. Substituindo na equação, obtemos:
\begin{equation*}
i_d = \varepsilon_0 A \frac{dE}{dt}
\end{equation*}

Para encontrar a \textbf{maior corrente de deslocamento}, precisamos do intervalo onde a taxa de variação do campo elétrico em módulo, $\left| \frac{\Delta E}{\Delta t} \right|$, seja máxima. 

Analisando os intervalos do gráfico (convertendo $\mu\text{s}$ para segundos: $1\,\mu\text{s} = 10^{-6}\,\text{s}$):
\begin{itemize}
    \item De $0$ a $2\,\mu\text{s}$: $\frac{\Delta E}{\Delta t} = 0$
    \item De $2$ a $6\,\mu\text{s}$: $\frac{\Delta E}{\Delta t} = \frac{2 - 6}{(6 - 2) \times 10^{-6}} = -1,0 \times 10^6\,\text{V/(m}\cdot\text{s)}$
    \item De $6$ a $7\,\mu\text{s}$: $\frac{\Delta E}{\Delta t} = \frac{4 - 2}{(7 - 6) \times 10^{-6}} = 2,0 \times 10^6\,\text{V/(m}\cdot\text{s)}$ \quad $\leftarrow$ \textbf{Máximo em módulo}
    \item De $7$ a $12\,\mu\text{s}$: $\frac{\Delta E}{\Delta t} = \frac{0 - 4}{(12 - 7) \times 10^{-6}} = -0,8 \times 10^6\,\text{V/(m}\cdot\text{s)}$
\end{itemize}

Substituindo a taxa máxima de $\frac{dE}{dt} = 2,0 \times 10^6\,\text{V/(m}\cdot\text{s)}$, a área $A = 2,0\,\text{m}^2$ e a permissividade do vácuo $\varepsilon_0 \approx 8,85 \times 10^{-12}\,\text{F/m}$:

\begin{align*}
i_d &= (8,85 \times 10^{-12}\,\text{F/m}) \cdot (2,0\,\text{m}^2) \cdot (2,0 \times 10^6\,\text{V/(m}\cdot\text{s)}) \\
i_d &= 35,4 \times 10^{-6}\,\text{A} = 35,4\,\mu\text{A}
\end{align*}

\textbf{Resposta:} A maior corrente de deslocamento é \textbf{$35,4\,\mu\text{A}$} (ou $3,54 \times 10^{-5}\,\text{A}$).


\subsection{Questão 74 - Arthur Roberto}
\subsection*{(a) Módulo do campo elétrico produzido pelo próton}

Tratando o próton como uma carga puntiforme, o módulo do campo elétrico $E$ a uma distância $r$ é calculado por:
\begin{equation*}
E = \frac{1}{4\pi\varepsilon_0} \frac{e}{r^2}
\end{equation*}

Substituindo os valores numéricos:
\begin{align*}
E &= (8,99 \times 10^9) \cdot \frac{1,60 \times 10^{-19}}{(5,2 \times 10^{-11})^2} \\
E &= \frac{1,4384 \times 10^{-9}}{2,704 \times 10^{-21}} \\
E &\approx 5,32 \times 10^{11} \text{ N/C}
\end{align*}

\subsection*{(b) Módulo do campo magnético do próton sobre o eixo $z$}

O campo magnético gerado por um dipolo magnético ao longo do seu eixo (eixo $z$) a uma distância $r$ é dado por:
\begin{equation*}
B = \frac{\mu_0}{4\pi} \frac{2\mu_{s,z}}{r^3}
\end{equation*}

Substituindo os valores numéricos:
\begin{align*}
B &= (10^{-7}) \cdot \frac{2 \cdot (1,4 \times 10^{-26})}{(5,2 \times 10^{-11})^3} \\
B &= 10^{-7} \cdot \frac{2,8 \times 10^{-26}}{1,406 \times 10^{-31}} \\
B &\approx 1,99 \times 10^{-2} \text{ T}
\end{align*}

\subsection*{(c) Razão entre o momento dipolar magnético do elétron e do próton}

Utilizando o valor do momento magnético do elétron ($\mu_{e,z} \approx \mu_B$) e dividindo pelo momento do próton fornecido:
\begin{equation*}
\text{Razão} = \frac{\mu_{e,z}}{\mu_{s,z}}
\end{equation*}

\begin{align*}
\text{Razão} &= \frac{9,27 \times 10^{-24} \text{ J/T}}{1,4 \times 10^{-26} \text{ J/T}} \\
\text{Razão} &\approx 662
\end{align*}

\subsection*{Respostas Finais}
\begin{itemize}
    \item \textbf{(a)} $5,32 \times 10^{11}$ N/C
    \item \textbf{(b)} $1,99 \times 10^{-2}$ T
    \item \textbf{(c)} $662$
\end{itemize}
% ================= CAPÍTULO 33 =================

\section{Capítulo 33}

\subsection{Questão 13 — Ana Alicy}
Inserir resolução aqui.

\subsection{Questão 14 — Kaylane Castro}
\textbf{Enunciado:}

Frank D. Drake, investigador do programa SETI, afirmou que o radiotelescópio de Arecibo é capaz de detectar um sinal que deposite em toda a superfície da Terra uma potência de apenas $1,0\text{ pW}$.

\begin{itemize}
    \item[(a)] Determinar a potência recebida pela antena de Arecibo, sabendo que seu diâmetro é de $300\text{ m}$.
    
    \item[(b)] Determinar a potência que uma fonte isotrópica situada no centro da galáxia deveria emitir para que esse sinal chegasse à Terra.
\end{itemize}

\vspace{0.5cm}

\textbf{Dados:}
$P_\text{Terra} = 1{,}0\times10^{-12}\ \text{W}$,\quad
$d_\text{antena} = 300\ \text{m}$,\quad
$R_\oplus = 6{,}37\times10^6\ \text{m}$,\quad
$R_\text{gal} = 2{,}2\times10^4\ \text{al}$,\quad
$1\ \text{al} = 9{,}46\times10^{15}\ \text{m}$.

\noindent\rule{\linewidth}{0.4pt}

\subsection*{(a) Potência recebida pela antena de Arecibo}

O enunciado afirma que $P_\text{Terra}$ é a potência depositada sobre
\emph{toda} a superfície terrestre. A antena capta uma fração
proporcional à sua área de abertura em relação ao disco projetado
da Terra (seção transversal circular que intercepta a onda plana
incidente):

\begin{equation}
    \frac{P_\text{antena}}{A_\text{antena}}
    =
    \frac{P_\text{Terra}}{A_\text{disco}}
    =
    \frac{P_\text{Terra}}{\pi R_\oplus^2}
    \label{eq:proporcao}
\end{equation}

com $A_\text{antena} = \pi r^2$ e $r = 150\ \text{m}$. Isolando:

\begin{equation}
    P_\text{antena}
    =
    P_\text{Terra}
    \left(\frac{r}{R_\oplus}\right)^2
    =
    (1{,}0\times10^{-12})
    \left(\frac{150}{6{,}37\times10^6}\right)^2
\end{equation}

\[
    \left(\frac{150}{6{,}37\times10^6}\right)^2
    = (2{,}355\times10^{-5})^2
    = 5{,}55\times10^{-10}
\]

\[
    \boxed{P_\text{antena} \approx 5{,}5\times10^{-22}\ \text{W}}
\]


\subsection*{(b) Potência da fonte isotrópica}

Para uma fonte isotrópica a distância $R$, a intensidade a essa
distância é:
\begin{equation}
    I = \frac{P_\text{fonte}}{4\pi R^2}
    \label{eq:isotropica}
\end{equation}

Essa mesma intensidade, ao incidir sobre a Terra, deposita uma potência
igual à intensidade multiplicada pela seção transversal (disco projetado)
da Terra:
\begin{equation}
    P_\text{Terra} = I \cdot \pi R_\oplus^2
    \implies
    I = \frac{P_\text{Terra}}{\pi R_\oplus^2}
    \label{eq:intensidade_terra}
\end{equation}

Igualando \eqref{eq:isotropica} e \eqref{eq:intensidade_terra}:

\begin{equation}
    \frac{P_\text{fonte}}{4\pi R^2}
    =
    \frac{P_\text{Terra}}{\pi R_\oplus^2}
    \implies
    P_\text{fonte}
    =
    4P_\text{Terra}
    \left(\frac{R}{R_\oplus}\right)^2
\end{equation}

Convertendo a distância galáctica:

\[
    R = (2{,}2\times10^4)(9{,}46\times10^{15})
    = 2{,}081\times10^{20}\ \text{m}
\]

\[
    \frac{R}{R_\oplus}
    = \frac{2{,}081\times10^{20}}{6{,}37\times10^6}
    = 3{,}267\times10^{13}
\]

\[
    P_\text{fonte}
    = 4(1{,}0\times10^{-12})(3{,}267\times10^{13})^2
    = (4{,}0\times10^{-12})(1{,}067\times10^{27})
\]

\[
    \boxed{P_\text{fonte} \approx 4{,}3\times10^{15}\ \text{W}}
\]

\noindent
\textit{Nota:} $4{,}3\times10^{15}\ \text{W}$ corresponde a
aproximadamente $4300$ vezes a potência total do Sol emitida
em rádio — o que ilustra o quão sensível o radiotelescópio
de Arecibo realmente é.

\subsection{Questão 15 — Dávyla Maria}
Inserir resolução aqui.

\subsection{Questão 16 — Rodrigo Albuquerque}

\tEspelhos esféricos. Um objeto \(O\) está sobre o eixo central de um espelho esférico. Para o problema 16, a Tabela 34-3 mostra a distância do objeto}

\[
p=+22\ cm
\]

\textbf{o tipo de espelho (convexo) e a distância focal}

\[
f=-35\ cm
\]

\textbf{Determine:}

\begin{itemize}
\item[(a)] o raio de curvatura \(r\) do espelho;
\item[(b)] a distância da imagem \(i\);
\item[(c)] a ampliação lateral \(m\);
\item[(d)] se a imagem é real (R) ou virtual (V);
\item[(e)] se é invertida (I) ou não-invertida (NI);
\item[(f)] se está do mesmo lado (M) do espelho que o objeto ou do lado oposto (O).
\end{itemize}
\includegraphics[width=0.5\textwidth]{image_5.png}
\vspace{0.5cm}

\textbf{(a) Raio de curvatura}

Para espelhos esféricos:

\[
r=2f
\]

Logo:

\[
r=2(-35)
\]

\[
r=-70\ cm
\]

\vspace{0.5cm}

\textbf{(b) Distância da imagem}

Equação dos espelhos:

\[
\frac{1}{f}=\frac{1}{p}+\frac{1}{i}
\]

Substituindo:

\[
\frac{1}{-35}=\frac{1}{22}+\frac{1}{i}
\]

\[
\frac{1}{i}
=
\frac{1}{-35}-\frac{1}{22}
\]

\[
\frac{1}{i}
=
\frac{-22-35}{770}
\]

\[
\frac{1}{i}
=
\frac{-57}{770}
\]

\[
i\approx-13,5\ cm
\]

\vspace{0.5cm}

\textbf{(c) Ampliação lateral}

\[
m=-\frac{i}{p}
\]

\[
m=-\frac{-13,5}{22}
\]

\[
m\approx0,614
\]

\vspace{0.5cm}

\textbf{(d) Real ou virtual}

Como:

\[
i<0
\]

a imagem é virtual (V).

\vspace{0.5cm}

\textbf{(e) Invertida ou não-invertida}

Como:

\[
m>0
\]

a imagem é não-invertida (NI).

\vspace{0.5cm}

\textbf{(f) Mesmo lado ou lado oposto}

Imagens virtuais em espelhos ficam atrás do espelho.

Logo, a imagem está do lado oposto (O).

\subsection{Questão 17 — Fábio Rodrigues}
\textbf{Dados}

\[
r = 10\,\text{km} = 10 \times 10^3\,\text{m}
\]

\[
I = 10\,\mu\text{W/m}^2 = 10 \times 10^{-6}\,\text{W/m}^2
\]

\[
c = 3,00 \times 10^8\,\text{m/s}
\]

\[
\mu_0 = 4\pi \times 10^{-7}\,\text{T}\cdot\text{m/A}
\]

\textbf{(a) Amplitude do campo elétrico (\(E_m\))}

A intensidade da onda eletromagnética relaciona-se com a amplitude do campo elétrico pela equação:

\[
I = \frac{E_m^2}{2 \mu_0 c}
\]

Isolando \(E_m\):

\[
E_m = \sqrt{2 \mu_0 c I}
\]

Substituindo os valores:

\[
E_m = \sqrt{2 (4\pi \times 10^{-7}) (3,00 \times 10^8) (10 \times 10^{-6})}
\]

\[
E_m = \sqrt{240\pi \times 10^{-5}}
\]

\[
E_m \approx \sqrt{0,00754}
\]

\[
E_m \approx 0,0868\,\text{V/m}
\]

\[
\boxed{E_m \approx 86,8\,\text{mV/m}}
\]

\textbf{(b) Amplitude do campo magnético (\(B_m\))}

A relação entre as amplitudes dos campos elétrico e magnético no vácuo é:

\[
B_m = \frac{E_m}{c}
\]

Substituindo o valor encontrado:

\[
B_m = \frac{0,0868}{3,00 \times 10^8}
\]

\[
B_m \approx 2,89 \times 10^{-10}\,\text{T}
\]

\[
\boxed{B_m = 2,89 \times 10^{-10}\,\text{T}}
\]

\textbf{(c) Potência da transmissão (\(P\))}

A intensidade é a potência transmitida por unidade de área. Como o transmissor irradia uniformemente sobre um hemisfério, a área considerada é a de uma semi-esfera (\(A = 2\pi r^2\)).

\[
I = \frac{P}{A} = \frac{P}{2\pi r^2}
\]

Isolando a potência \(P\):

\[
P = I (2\pi r^2)
\]

Substituindo os valores:

\[
P = (10 \times 10^{-6}) [2\pi (10 \times 10^3)^2]
\]

\[
P = (10^{-5}) (2\pi \times 10^8)
\]

\[
P = 2\pi \times 10^3
\]

\[
P \approx 6283\,\text{W}
\]

\[
\boxed{P \approx 6,28\,\text{kW}}
\]

\subsection{Questão 18 — Mateus Valentim}
\subsubsection{Resolução}
\textbf{Resolução}

Pelo gráfico, observa-se que:

\[
m = 1 \quad \text{quando} \quad p = 0,
\]

e

\[
m \approx 0{,}45 \quad \text{quando} \quad p = p_s = 10\,\text{cm}.
\]

Para um espelho esférico convexo, a ampliação lateral é dada por

\[
m=\frac{f}{p+f}.
\]

Utilizando o ponto \(p=10\,\text{cm}\) e \(m\approx 0{,}45\), temos:

\[
0{,}45=\frac{f}{10+f}.
\]

Multiplicando ambos os lados por \((10+f)\):

\[
0{,}45(10+f)=f.
\]

Desenvolvendo:

\[
4{,}5+0{,}45f=f.
\]

Logo,

\[
4{,}5=0{,}55f,
\]

e

\[
f\approx 8{,}18\,\text{cm}.
\]

Para \(p=21\,\text{cm}\), a ampliação será

\[
m=\frac{8{,}18}{21+8{,}18}.
\]

Assim,

\[
m\approx \frac{8{,}18}{29{,}18}
\approx 0{,}28.
\]

Portanto, a ampliação do objeto quando ele se encontra a \(21\,\text{cm}\) do espelho é

\[
\boxed{m \approx 0{,}28}.
\]

\subsection{Questão 23 — Cicero Rodrigues}

\textbf{23. }\normalfont Pretende-se levitar uma pequena esfera, totalmente absorvente, 0,500 m acima de uma fonte luminosa pontual, fazendo com que a força para cima exercida pela radiação seja igual ao peso da esfera. A esfera tem 2,00 mm de raio e massa específica de 19,0 g/cm$^3$. \textbf{(a)} Qual deve ser a potência da fonte luminosa? \textbf{(b)} Mesmo que fosse possível construir uma fonte com essa potência, por que o equilíbrio da esfera seria instável?

\textbf{Resolução:}
\vspace{0.3cm}

\textbf{Dados:}
\[y = 0,500\;m \quad \text{(Distância da fonte)}\]
\[r = 2,00\;mm = 2,00 \times 10^{-3}\;m \quad \text{(Raio da esfera)}\]
\[\rho = 19,0\;g/cm^3 = 19,0 \times 10^3\;kg/m^3 \quad \text{(Massa específica)}\]
\[c = 3,00 \times 10^8\;m/s\]
\[g = 9,8\;m/s^2\]

\vspace{0.2cm}
\textbf{(a) Potência da fonte luminosa (\(P\))}
Para a esfera levitar, a força de radiação (\(F_r\)) apontando para cima deve equilibrar a força peso (\(F_g\)) apontando para baixo.
A força peso da esfera é dada pelo seu volume e massa específica:
\[F_g = mg = (\rho V)g = \rho \left( \frac{4}{3}\pi r^3 \right) g\]

A força de radiação para uma superfície totalmente absorvente é a pressão de radiação multiplicada pela área da seção transversal interceptada pela luz (\(A = \pi r^2\)):
\[F_r = \left( \frac{I}{c} \right) A = \left( \frac{P}{4\pi y^2 c} \right) (\pi r^2) = \frac{P r^2}{4 y^2 c}\]

Igualando as duas forças (\(F_r = F_g\)) e isolando \(P\):
\[\frac{P r^2}{4 y^2 c} = \rho \frac{4}{3}\pi r^3 g\]
\[P = \frac{16\pi}{3} \rho r g y^2 c\]

Substituindo os valores no SI:
\[P = \frac{16\pi}{3} (19,0 \times 10^3)(2,00 \times 10^{-3})(9,8)(0,500)^2(3,00 \times 10^8)\]
\[P = \frac{16\pi}{3} (38,0)(9,8)(0,25)(3,00 \times 10^8)\]
\[P = 16\pi (38,0)(9,8)(0,25)(10^8)\]
\[P = 4\pi (372,4) \times 10^8\]
\[P = 1489,6\pi \times 10^8 \approx 4679,7 \times 10^8\;W\]
Ajustando a notação científica:
\[\underline{P \approx 4,68 \times 10^{11}\;W}\]

\vspace{0.2cm}
\textbf{(b) Por que o equilíbrio seria instável?}
A fonte luminosa é pontual, o que significa que seus raios de luz divergem radialmente em todas as direções. 
\[\underline{\text{Se a esfera sofrer qualquer pequeno deslocamento lateral (horizontal),}}\]
\[\underline{\text{os raios de luz incidentes não apontarão mais perfeitamente para cima.}}\]
Em vez disso, a força de radiação passará a ter uma componente horizontal que empurrará a esfera ainda mais para o lado, afastando-a do eixo central irreversivelmente, até que ela caia (não há força restauradora horizontal).

\subsection{Questão 24 — Guilherme Araújo Floriano}

\subsection*{(a)}

Como a área da seção reta do feixe é $\frac{\pi d^2}{4}$, na qual d é o diâmetro da esfera, a intensidade
do feixe é:

\[
I = \frac{P}{A} = \frac{P}{\frac{\pi d^2}{4}} = \frac{5 * 10^{-3} W}{0.786 (1266*10^{-9}m)^2} = 3,97 GW/m^2
\]

\subsection*{(b)}

A pressão de radiação é

\[
p_r = \frac{I}{c} = \frac{3,97 * 10^9 W/m^2}{3 * 10^8} = 13,2 Pa    
\]

\subsection*{(c)}

A força associada à radiação é igual à pressão multiplicada pela área da seção reta do feixe, que, por sua vez, é igual a P/I

\[
F_r = p_r\frac{\pi d^2}{4} = p_r\frac{P}{I} = 13,2 Pa * \frac{5*10^{-3}W}{3,97*10^9 W/m^2} = 1,67 * 10^{-11} N
\]

\subsection*{(d)}

A aceleração da esfera é

\[
a = \frac{F_r}{m} = \frac{F_r}{p(\pi*d^3/6)} = \frac{6*1,67*10^{-11}}{\pi * 5*10^3 * (1266*10^{-9})^3} = 3,14 * 10^3 m/s^2
\]

\subsection{Questão 25 — Arthur Roberto da Silva}
\subsection*{(a) Cálculo da Frequência ($f$)}

A relação entre a velocidade da luz no vácuo ($c \approx 2,998 \times 10^8\text{ m/s}$), o comprimento de onda ($\lambda$) e a frequência ($f$) é dada por $c = \lambda f$. Isolando a frequência:
\begin{equation*}
f = \frac{c}{\lambda} = \frac{2,998 \times 10^8\text{ m/s}}{3,0\text{ m}} \approx 1,0 \times 10^8\text{ Hz}
\end{equation*}

\subsection*{(b) Cálculo da Frequência Angular ($\omega$)}

A frequência angular representa a taxa de variação da fase da onda em radianos por segundo:
\begin{equation*}
\omega = 2\pi f = 2\pi (1,0 \times 10^8\text{ Hz}) \approx 6,3 \times 10^8\text{ rad/s}
\end{equation*}

\subsection*{(c) Cálculo do Número de Onda ($k$)}

O número de onda representa a magnitude do vetor de onda espacial, indicando quantos ciclos a onda possui por unidade de espaço:
\begin{equation*}
k = \frac{2\pi}{\lambda} = \frac{2\pi}{3,0\text{ m}} \approx 2,1\text{ rad/m}
\end{equation*}

\subsection*{(d) Cálculo da Amplitude do Campo Magnético ($B_m$)}

Nas ondas eletromagnéticas que se propagam no vácuo, as amplitudes dos campos elétrico e magnético estão diretamente acopladas através da velocidade da luz pela relação $E_m = c B_m$:
\begin{equation*}
B_m = \frac{E_m}{c} = \frac{300\text{ V/m}}{2,998 \times 10^8\text{ m/s}} \approx 1,0 \times 10^{-6}\text{ T} = 1,0\text{ }\mu\text{T}
\end{equation*}

\subsection*{(e) Direção de Oscilação do Campo Magnético}

Os vetores do campo elétrico ($\vec{E}$), do campo magnético ($\vec{B}$) e da direção de propagação (vetor de Poynting $\vec{S}$) são mutuamente ortogonais. Como a onda se propaga paralelamente ao eixo $x$ e o campo elétrico oscila paralelamente ao eixo $y$, o campo magnético deve oscilar obrigatoriamente de forma paralela ao \textbf{eixo $z$}.

\subsection*{(f) Cálculo da Intensidade da Onda ($I$)}

A intensidade média de uma onda eletromagnética senoidal corresponde à magnitude média do vetor de Poynting, dependendo da amplitude do campo elétrico e da permeabilidade magnética do vácuo ($\mu_0 = 4\pi \times 10^{-7}\text{ H/m}$):
\begin{equation*}
I = \frac{E_m^2}{2\mu_0 c} = \frac{(300\text{ V/m})^2}{2 \cdot (4\pi \times 10^{-7}\text{ H/m}) \cdot (2,998 \times 10^8\text{ m/s})} \approx 119\text{ W/m}^2 \approx 1,2 \times 10^2\text{ W/m}^2
\end{equation*}

\subsection*{(g) Taxa de Transferência de Momento Linear ($\frac{dp}{dt}$)}

Visto que a placa absorve totalmente a radiação incidente, a taxa com que o momento linear é transferido para a superfície por unidade de área é dada pela razão entre a intensidade e a velocidade da luz ($I/c$). Para encontrar a força resultante exercida sobre a área total $A$, multiplicamos o termo pela área da placa:
\begin{equation*}
\frac{dp}{dt} = \frac{I A}{c} = \frac{(119\text{ W/m}^2) \cdot (2,0\text{ m}^2)}{2,998 \times 10^8\text{ m/s}} \approx 8,0 \times 10^{-7}\text{ N}
\end{equation*}

\subsection*{(h) Cálculo da Pressão da Radiação ($p_r$)}

A pressão de radiação é definida como a força exercida por unidade de área ($F/A$). Como a força é equivalente à taxa de variação do momento linear temporal ($\frac{dp}{dt}$), dividimos a força obtida no item anterior pela área da placa (o que retorna diretamente à expressão de absorção total $p_r = \frac{I}{c}$):
\begin{equation*}
p_r = \frac{1}{A}\left(\frac{dp}{dt}\right) = \frac{8,0 \times 10^{-7}\text{ N}}{2,0\text{ m}^2} = 4,0 \times 10^{-7}\text{ Pa}
\end{equation*}

\subsection{Questão 26 — Arthur Roberto da Silva}
No equilíbrio mecânico, a força de atração gravitacional solar ($F_g$) é igual à força da pressão de radiação ($F_r$):
\begin{equation*}
F_g = F_r \implies G \frac{m M_S}{d_{TS}^2} = \frac{2 I A}{c}
\end{equation*}

Isolando a área $A$ da vela solar:
\begin{equation*}
A = \frac{c G m M_S}{2 I d_{TS}^2}
\end{equation*}

Substituindo os dados no Sistema Internacional (S.I.):
\begin{align*}
A &= \frac{(2,998 \times 10^8) \cdot (6,67 \times 10^{-11}) \cdot 1500 \cdot (1,99 \times 10^{30})}{2 \cdot (1,40 \times 10^3) \cdot (1,50 \times 10^{11})^2} \\
A &= \frac{5,969 \times 10^{30}}{6,30 \times 10^{25}} \\
A &\approx 9,5 \times 10^4\text{ m}^2 = 0,095\text{ km}^2
\end{align*}

\noindent \textit{Nota:} O gabarito original do livro contém um erro de digitação no expoente ao indicar $9,5 \times 10^5\text{ m}^2$ ($0,95\text{ km}^2$). O cálculo correto resulta em $9,5 \times 10^4\text{ m}^2$.

\subsection*{Resposta Final}
\textbf{$A = 9,5 \times 10^4\text{ m}^2$} (ou \textbf{$0,095\text{ km}^2$}).

\subsection{Questão 27 — Guilherme Araújo Floriano}

Se a luz do laser remove uma energia $U$ da espaçonave, remove também um momento $p = U/c$. Como o momento da espaçonave e o momento da luz são conservados, este é o momento adquirido pela espaçonave. 

Se $P$ é a potência do laser, a energia removida em um intervalo de tempo $\Delta t$ é $U = P\Delta t$. 

Assim, $p = \frac{P\Delta t}{c}$ e, se $m$ é a massa da espaçonave, a velocidade que ela atinge é

\[
v = \frac{p}{m} = \frac{Pt}{mc} = \frac{10^4W*86400s}{1,5*10^3 kg*3*10^8 m/s^2} = 1,9*10^{-3} m/s = 1,9 mm/s
\]

\subsection{Questão 28 — Cicero Rodrigues}

\textbf{28. }\normalfont A intensidade média da radiação solar que incide perpendicularmente em uma superfície situada logo acima da atmosfera da Terra é 1,4 kW/m$^2$. \textbf{(a)} Qual é a pressão de radiação $p_r$ exercida pelo Sol sobre a superfície, supondo que toda a radiação é absorvida? \textbf{(b)} Calcule a razão entre essa pressão e a pressão atmosférica ao nível do mar, que é $1,0 \times 10^5$ Pa.

\textbf{Resolução:}
\vspace{0.3cm}

\textbf{Dados:}
\[I = 1,4\;kW/m^2 = 1,4 \times 10^3\;W/m^2\]
\[c = 3,00 \times 10^8\;m/s \quad \text{(Velocidade da luz)}\]
\[p_{atm} = 1,0 \times 10^5\;Pa\]

\vspace{0.2cm}
\textbf{(a) Pressão de radiação (\(p_r\))}
Quando a radiação eletromagnética incide perpendicularmente e é totalmente absorvida pela superfície, a pressão de radiação é a razão entre a intensidade da onda e a velocidade da luz:
\[p_r = \frac{I}{c}\]
Substituindo os valores:
\[p_r = \frac{1,4 \times 10^3}{3,00 \times 10^8}\]
\[p_r \approx 0,4667 \times 10^{-5}\;Pa\]
Ajustando para notação científica padrão:
\[\underline{p_r \approx 4,7 \times 10^{-6}\;Pa}\]

\vspace{0.2cm}
\textbf{(b) Razão com a pressão atmosférica (\(p_r / p_{atm}\))}
Basta dividir o valor encontrado na letra (a) pela pressão atmosférica ao nível do mar fornecida no enunciado:
\[\text{Razão} = \frac{p_r}{p_{atm}}\]
\[\text{Razão} = \frac{4,667 \times 10^{-6}}{1,0 \times 10^5}\]
Como estamos dividindo por uma potência de 10, basta subtrair os expoentes \((-6 - 5 = -11)\):
\[\text{Razão} \approx 4,667 \times 10^{-11}\]
\[\underline{\text{Razão} \approx 4,7 \times 10^{-11}}\]

\subsection{Questão 29 — Mateus Valentim}
\subsubsection{Resolução}
\textbf{Dados}

\[
r = 2,00\,\text{mm} = 2,00\times10^{-3}\,\text{m},
\]

\[
h = 0,500\,\text{m},
\]

\[
\rho = 19,0\,\text{g/cm}^3
      = 1,90\times10^{4}\,\text{kg/m}^3.
\]

A esfera é totalmente absorvente, logo a força de radiação é

\[
F_{\text{rad}}=\frac{IA}{c},
\]

onde \(I\) é a intensidade da radiação na posição da esfera,
\(A=\pi r^2\) é sua área de seção reta e \(c\) é a velocidade da luz.

Para que a esfera fique em equilíbrio,

\[
F_{\text{rad}}=mg.
\]

\bigskip

\textbf{(a) Potência necessária da fonte}

O volume da esfera é

\[
V=\frac{4}{3}\pi r^3,
\]

e sua massa é

\[
m=\rho V
  =\rho\frac{4}{3}\pi r^3.
\]

Substituindo os valores:

\[
m
=(1,90\times10^{4})
\frac{4}{3}\pi(2,00\times10^{-3})^3
\approx 6,37\times10^{-4}\,\text{kg}.
\]

Assim,

\[
mg
=(6,37\times10^{-4})(9,8)
\approx 6,24\times10^{-3}\,\text{N}.
\]

Da condição de equilíbrio,

\[
\frac{I\pi r^2}{c}=mg,
\]

logo

\[
I=\frac{mg\,c}{\pi r^2}.
\]

Substituindo os valores:

\[
I=
\frac{(6,24\times10^{-3})(3,00\times10^{8})}
{\pi(2,00\times10^{-3})^2}
\approx 1,49\times10^{11}\,\text{W/m}^2.
\]

Para uma fonte isotrópica,

\[
I=\frac{P}{4\pi h^2},
\]

portanto

\[
P=4\pi h^2 I.
\]

Com \(h=0,500\,\text{m}\),

\[
P
=4\pi(0,500)^2(1,49\times10^{11})
\approx 4,7\times10^{11}\,\text{W}.
\]

Logo,

\[
\boxed{P \approx 4,7\times10^{11}\ \text{W}}.
\]

\bigskip

\textbf{(b) Estabilidade do equilíbrio}

A intensidade da radiação varia com a distância segundo

\[
I=\frac{P}{4\pi r^2}.
\]

Assim, a força de radiação também varia como

\[
F_{\text{rad}}\propto \frac{1}{r^2}.
\]

Como a força de radiação varia com \(1/r^2\), um pequeno deslocamento altera essa força de modo a afastar ainda mais a esfera da posição de equilíbrio. Portanto, o equilíbrio é

\[
\boxed{\text{instável}}.
\]

\subsection{Questão 30 — Fábio Rodrigues}
\textbf{Dados}

\[
I_s = 200\,\text{W/m}^2
\]

\[
r_s^{-2} = 8,0\,\text{m}^{-2}
\]

A intensidade \(I\) da luz emitida por uma fonte pontual e isotrópica é dada por:

\[
I = \frac{P}{4\pi r^2}
\]

Onde \(P\) é a potência da fonte e \(r\) é a distância.

Podemos reescrever a equação isolando a relação entre \(I\) e \(r^{-2}\):

\[
I = \left(\frac{P}{4\pi}\right) r^{-2}
\]

Isso mostra que o gráfico de \(I\) em função de \(r^{-2}\) é uma reta que passa pela origem. A inclinação (coeficiente angular) dessa reta, \(m\), corresponde ao termo entre parênteses:

\[
m = \frac{I}{r^{-2}} = \frac{P}{4\pi}
\]

Utilizando os valores máximos das escalas fornecidas pelo gráfico, onde a reta atinge \(I = 200\,\text{W/m}^2\) para \(r^{-2} = 8,0\,\text{m}^{-2}\), calculamos a inclinação:

\[
m = \frac{200}{8,0}
\]

\[
m = 25\,\text{W}
\]

Igualando esse valor à expressão da potência:

\[
\frac{P}{4\pi} = 25
\]

\[
P = 25(4\pi)
\]

\[
P = 100\pi
\]

\[
P \approx 314,16\,\text{W}
\]

\[
\boxed{P \approx 314\,\text{W}}
\]

\subsection{Questão 36 — Rodrigo Albuquerque}

\textbf{ Superfícies refratoras esféricas. Um objeto \(O\) está sobre o eixo central de uma superfície refratora esférica. Para o problema 36, a Tabela 34-5 fornece:}

\[
n_1=1,0
\]

\[
n_2=1,5
\]

\[
p=+10\ cm
\]

\[
r=+30\ cm
\]

\textbf{Determine:}

\includegraphics[width=0.5\textwidth]{image_7.png}

\begin{itemize}
\item[(a)] o índice de refração \(n_2\);
\item[(b)] a distância do objeto \(p\);
\item[(c)] o raio de curvatura \(r\);
\item[(d)] a distância da imagem \(i\);
\item[(e)] se a imagem é real (R) ou virtual (V);
\item[(f)] se a imagem fica do mesmo lado da superfície que o objeto (M) ou do lado oposto (O).
\end{itemize}

\vspace{0.5cm}

\textbf{(a)}

\[
n_2=1,5
\]

\vspace{0.5cm}

\textbf{(b)}

\[
p=+10\ cm
\]

\vspace{0.5cm}

\textbf{(c)}

\[
r=+30\ cm
\]

\vspace{0.5cm}

\textbf{(d) Distância da imagem}

Para superfícies refratoras esféricas:

\[
\frac{n_1}{p}+\frac{n_2}{i}
=
\frac{n_2-n_1}{r}
\]

Substituindo:

\[
\frac{1,0}{10}
+
\frac{1,5}{i}
=
\frac{1,5-1,0}{30}
\]

\[
0,1+\frac{1,5}{i}
=
0,0167
\]

\[
\frac{1,5}{i}
=
0,0167-0,1
\]

\[
\frac{1,5}{i}
=
-0,0833
\]

\[
i=\frac{1,5}{-0,0833}
\]

\[
i\approx-18\ cm
\]

\vspace{0.5cm}

\textbf{(e) Real ou virtual}

Como:

\[
i<0
\]

a imagem é virtual (V).

\vspace{0.5cm}

\textbf{(f) Lado da imagem}

Imagens virtuais ficam do mesmo lado da superfície que o objeto.

Logo:

\[
\text{lado} = M
\]
```


\subsection{Questão 37 — Dávyla Maria}
Inserir resolução aqui.

\subsection{Questão 38 — Kaylane Castro}
\subsection*{Enunciado}

Um sinal de rádio proveniente de uma fonte isotrópica localizada no centro da nossa galáxia (a uma distância de $2,2 \times 10^4 \text{ al}$) atinge a Terra. A potência total do sinal depositada sobre toda a superfície terrestre é de $1,0 \times 10^{-12} \text{ W}$. Sabendo que o radiotelescópio de Arecibo possui uma antena circular com diâmetro de $300 \text{ m}$:

\begin{enumerate}
    \item[\textbf{(a)}] Que potência é recebida pela antena de Arecibo?
    \item[\textbf{(b)}] Qual é a potência total emitida por essa fonte isotrópica?
\end{enumerate}

% ---------------------------------------------------------

\subsection*{Análise Inicial e Dados}

Para resolver o problema, coletamos os dados fornecidos e os convertemos para o Sistema Internacional (SI):

* Potência total na Terra: $P_{\text{Terra}} = 1,0 \times 10^{-12} \text{ W}$
* Diâmetro da antena: $d_{\text{antena}} = 300 \text{ m} \implies \text{Raio } r = 150 \text{ m}$
* Raio da Terra: $R_{\oplus} = 6,37 \times 10^6 \text{ m}$
* Distância até a fonte: $R = 2,2 \times 10^4 \text{ al}$
* Conversão de ano-luz: $1 \text{ al} = 9,46 \times 10^{15} \text{ m}$

Convertendo a distância $R$ para metros:
\[
R = (2,2 \times 10^4) \times (9,46 \times 10^{15} \text{ m}) = 2,081 \times 10^{20} \text{ m}
\]

% ---------------------------------------------------------

\subsection*{Resolução: Item (a)}

A onda plana incidente é interceptada pelo disco projetado da Terra (seção transversal circular de raio $R_{\oplus}$). A antena de Arecibo capta uma fração dessa potência que é diretamente proporcional à sua própria área de abertura ($A_{\text{antena}} = \pi r^2$) em relação à área desse disco terrestre ($A_{\text{disco}} = \pi R_{\oplus}^2$).

Montando a relação de proporção:
\[
\frac{P_{\text{antena}}}{A_{\text{antena}}} = \frac{P_{\text{Terra}}}{A_{\text{disco}}} \implies \frac{P_{\text{antena}}}{\pi r^2} = \frac{P_{\text{Terra}}}{\pi R_{\oplus}^2}
\]

Isolando a potência da antena ($P_{\text{antena}}$):
\[
P_{\text{antena}} = P_{\text{Terra}} \left( \frac{r}{R_{\oplus}} \right)^2
\]

Substituindo os valores numéricos:
\[
P_{\text{antena}} = (1,0 \times 10^{-12}) \left( \frac{150}{6,37 \times 10^6} \right)^2
\]

\[
P_{\text{antena}} = (1,0 \times 10^{-12}) \times (2,355 \times 10^{-5})^2 = (1,0 \times 10^{-12}) \times (5,55 \times 10^{-10})
\]

\[
P_{\text{antena}} \approx 5,55 \times 10^{-22} \text{ W}
\]

% ---------------------------------------------------------

\subsection*{Resolução: Item (b)}

A intensidade $I$ da onda esférica emitida por uma fonte isotrópica a uma distância $R$ espalha-se pela área de uma esfera de raio $R$:
\[
I = \frac{P_{\text{fonte}}}{4\pi R^2}
\]

Essa mesma intensidade, ao atingir a Terra, distribui a potência $P_{\text{Terra}}$ sobre a área de sua seção transversal ($A_{\text{disco}} = \pi R_{\oplus}^2$):
\[
I = \frac{P_{\text{Terra}}}{\pi R_{\oplus}^2}
\]

Igualando as duas expressões para a intensidade $I$:
\[
\frac{P_{\text{fonte}}}{4\pi R^2} = \frac{P_{\text{Terra}}}{\pi R_{\oplus}^2}
\]

Isolando a potência da fonte ($P_{\text{fonte}}$):
\[
P_{\text{fonte}} = 4 P_{\text{Terra}} \left( \frac{R}{R_{\oplus}} \right)^2
\]

Substituindo os valores conhecidos:
\[
P_{\text{fonte}} = 4 \times (1,0 \times 10^{-12}) \times \left( \frac{2,081 \times 10^{20}}{6,37 \times 10^6} \right)^2
\]

\[
P_{\text{fonte}} = (4,0 \times 10^{-12}) \times (3,267 \times 10^{13})^2
\]

\[
P_{\text{fonte}} = (4,0 \times 10^{-12}) \times (1,067 \times 10^{27})
\]

\[
P_{\text{fonte}} \approx 4,27 \times 10^{15} \text{ W}
\]

% ---------------------------------------------------------

\subsection*{Respostas}

\[
\text{\textbf{(a)}} \quad \boxed{P_{\text{antena}} \approx 5,5 \times 10^{-22} \text{ W}}
\]

\[
\text{\textbf{(b)}} \quad \boxed{P_{\text{fonte}} \approx 4,3 \times 10^{15} \text{ W}}
\]

\vspace{0.3cm}
\noindent\textit{Nota: Esta potência astronômica equivale a milhares de vezes a emissão total do Sol na faixa de rádio, o que destaca a incrível sensibilidade de captação do radiotelescópio de Arecibo.}

\subsection{Questão 39 — Ana Alicy}
Inserir resolução aqui.

\subsection{Questão 40 — Ana Alicy}
Inserir resolução aqui.

\subsection{Questão 41 — Kaylane Castro}

\subsection*{Enunciado}

Um feixe de luz parcialmente polarizada pode ser considerado uma mistura de luz polarizada e não-polarizada. Suponha que um feixe desse tipo atravesse um filtro polarizador e que o filtro seja girado de $360^\circ$ enquanto se mantém perpendicular ao feixe.

Se a intensidade da luz transmitida varia por um fator de $5,0$ durante a rotação do filtro, determine que fração da intensidade da luz incidente está associada à luz polarizada do feixe.

% ---------------------------------------------------------

\subsection*{Modelagem Matemática}

A intensidade total incidente é dada pela soma da intensidade da componente não-polarizada com a intensidade da componente polarizada:

\[
I_0 = I_{np} + I_p
\]

Ao atravessar o polarizador:

\begin{itemize}
    \item A componente não-polarizada perde metade de sua intensidade:
    
    \[
    I'_{np} = \frac{1}{2}I_{np}
    \]

    \item A componente polarizada obedece à Lei de Malus:
    
    \[
    I'_p = I_p\cos^2\theta
    \]
\end{itemize}

Logo, a intensidade transmitida em função do ângulo é:

\[
I(\theta) = \frac{1}{2}I_{np} + I_p\cos^2\theta
\]

% ---------------------------------------------------------

\subsection*{Intensidade Máxima e Mínima}

A intensidade transmitida será máxima quando:

\[
\cos^2\theta = 1
\]

Assim:

\[
I_{\text{máx}} = \frac{1}{2}I_{np} + I_p
\]

A intensidade será mínima quando:

\[
\cos^2\theta = 0
\]

Portanto:

\[
I_{\text{mín}} = \frac{1}{2}I_{np}
\]

O problema informa que a intensidade varia por um fator de $5,0$:

\[
\frac{I_{\text{máx}}}{I_{\text{mín}}} = 5,0
\]

% ---------------------------------------------------------

\subsection*{Resolução}

Substituindo as expressões de intensidade máxima e mínima:

\[
\frac{
\frac{1}{2}I_{np} + I_p
}{
\frac{1}{2}I_{np}
}
= 5,0
\]

Multiplicando ambos os lados por $\frac{1}{2}I_{np}$:

\[
\frac{1}{2}I_{np} + I_p
=
5,0\left(\frac{1}{2}I_{np}\right)
\]

Multiplicando toda a equação por 2:

\[
I_{np} + 2I_p = 5I_{np}
\]

Isolando $I_p$:

\[
2I_p = 5I_{np} - I_{np}
\]

\[
2I_p = 4I_{np}
\]

\[
I_p = 2I_{np}
\]

A intensidade total incidente é:

\[
I_0 = I_{np} + I_p
\]

Substituindo $I_p = 2I_{np}$:

\[
I_0 = I_{np} + 2I_{np}
\]

\[
I_0 = 3I_{np}
\]

A fração associada à luz polarizada é:

\[
\frac{I_p}{I_0}
=
\frac{2I_{np}}{3I_{np}}
\]

\[
\frac{I_p}{I_0} = \frac{2}{3}
\]

Em porcentagem:

\[
\frac{2}{3} \approx 0,667
\]

\[
0,667 \times 100\% \approx 66,7\%
\]

% ---------------------------------------------------------

\subsection*{Resposta}

\[
\boxed{
\frac{I_p}{I_0} = \frac{2}{3} \approx 66,7\%
}
\]


```


\subsection{Questão 42 — Dávyla Maria}
Inserir resolução aqui.

\subsection{Questão 43 — Rodrigo Albuquerque}

\textbf{ Você produz uma imagem do Sol em uma tela usando uma lente delgada com uma distância focal de \(20,0\ cm\). Qual é o diâmetro da imagem?}

\vspace{0.5cm}

Como o Sol está muito distante da lente:

\[
p\rightarrow\infty
\]

Logo, a imagem forma-se aproximadamente no foco:

\[
i=f=20,0\ cm
\]

A ampliação lateral é:

\[
m=\frac{i}{p}
\]

Mas, para objetos muito distantes:

\[
m=\frac{f}{d}
\]

onde:

\[
d=1,50\times10^{11}\ m
\]

é a distância da Terra ao Sol.

O diâmetro do Sol é:

\[
D=1,39\times10^{9}\ m
\]

Então, o diâmetro da imagem será:

\[
D_i=mD
\]

\[
D_i=\frac{f}{d}D
\]

Convertendo a distância focal:

\[
f=20,0\ cm=0,200\ m
\]

Substituindo:

\[
D_i=
\frac{0,200}{1,50\times10^{11}}
(1,39\times10^{9})
\]

\[
D_i\approx1,85\times10^{-3}\ m
\]

\[
D_i\approx1,85\ mm
\]


\subsection{Questão 44 — Fábio Rodrigues}
\textbf{Dados}

A intensidade da luz após passar por três polarizadores, sendo a luz inicial não-polarizada, é dada por:

\[
I = \frac{1}{2} I_0 \cos^2(\theta_2 - \theta_1) \cos^2(\theta_3 - \theta_2)
\]

A intensidade é nula quando o segundo filtro está perpendicular ao primeiro ou ao terceiro, ou seja:

\[
\theta_2 = \theta_1 \pm 90^\circ \quad \text{ou} \quad \theta_2 = \theta_3 \pm 90^\circ
\]

Pela análise do gráfico, os mínimos de intensidade (zero) ocorrem em:

\[
\theta_2 = 60^\circ \quad \text{e} \quad \theta_2 = 140^\circ
\]

Podemos determinar os ângulos fixos (como a ordem no produto não importa, podemos associar livremente):

\[
\theta_1 = 60^\circ + 90^\circ = 150^\circ
\]

\[
\theta_3 = 140^\circ - 90^\circ = 50^\circ
\]

Substituindo \(\theta_1\), \(\theta_3\) e o valor de \(\theta_2 = 90^\circ\) pedido na questão:

\[
I = \frac{1}{2} I_0 \cos^2(90^\circ - 150^\circ) \cos^2(50^\circ - 90^\circ)
\]

\[
I = 0,5 I_0 \cos^2(-60^\circ) \cos^2(-40^\circ)
\]

Calculando os cossenos:

\[
I = 0,5 I_0 (0,5)^2 (\approx 0,7660)^2
\]

\[
I = 0,5 I_0 (0,25) (0,5868)
\]

\[
I \approx 0,07335 I_0
\]

A porcentagem da intensidade inicial transmitida é:

\[
\frac{I}{I_0} = 0,07335
\]

\[
\boxed{\frac{I}{I_0} \approx 7,34\%}
\]

\subsection{Questão 50 — Mateus Valentim}


\subsubsection{Resolução}



\textbf{(a) Determinação de \(n_1\)}



Pelo gráfico, para



\[

n_2=1{,}3,

\]



temos



\[

\theta_2=\theta_{2b}=40^\circ.

\]



Da figura, o ângulo de incidência é



\[

\theta_1 \approx 30^\circ.

\]



Aplicando a lei de Snell:



\[

n_1\sin\theta_1=n_2\sin\theta_2.

\]



Logo,



\[

n_1=\frac{n_2\sin\theta_2}{\sin\theta_1}

=\frac{(1{,}3)\sin40^\circ}{\sin30^\circ}.

\]



Assim,



\[

n_1\approx\frac{1{,}3(0{,}643)}{0{,}500}

\approx 1{,}67.

\]



Portanto,



\[

\boxed{n_1 \approx 1{,}67}.

\]



\bigskip



\textbf{(b) Novo ângulo de refração}



Para



\[

\theta_1=60^\circ

\qquad\text{e}\qquad

n_2=2{,}4,

\]



a lei de Snell fornece



\[

n_1\sin\theta_1=n_2\sin\theta_2.

\]



Substituindo os valores:



\[

(1{,}67)\sin60^\circ

=

(2{,}4)\sin\theta_2.

\]



\[

\sin\theta_2

=

\frac{1{,}67\sin60^\circ}{2{,}4}

\approx 0{,}602.

\]



Portanto,



\[

\theta_2=\sin^{-1}(0{,}602)

\approx 37^\circ.

\]



Logo,



\[

\boxed{\theta_2 \approx 37^\circ}.

\]



\subsection{Questão 51 — Cicero Rodrigues}

\textbf{51. }\normalfont Na Fig. 33.26, a luz incide, fazendo um ângulo $\theta_1 = 40,1^\circ$ com a normal, na interface de dois meios transparentes. Parte da luz atravessa as outras três camadas transparentes e parte é refletida para cima e escapa para o ar. Se $n_1 = 1,30$, $n_2 = 1,40$, $n_3 = 1,32$ e $n_4 = 1,45$, determine o valor \textbf{(a)} de $\theta_5$ e \textbf{(b)} de $\theta_4$.

\includegraphics[width=0.4\textwidth]{image_13.png}

\textbf{Resolução:}
\vspace{0.3cm}

\textbf{Dados:}
\[\theta_1 = 40,1^\circ\ \quad ; \quad n_{\text{ar}} = 1,00\]
\[n_1 = 1,30 \quad ; \quad n_3 = 1,32\]
\[n_2 = 1,40 \quad ; \quad n_4 = 1,45\]

\vspace{0.2cm}
\textbf{(a) Valor do ângulo \(\theta_5\)}
Pela lei da reflexão, o raio refletido na interface \(n_1/n_2\) retorna através do meio \(n_1\) com o mesmo ângulo em relação à normal (\(\theta_{\text{refl}} = \theta_1\)). Como as interfaces horizontais são paralelas, esse raio incide na interface superior (\(n_1/\text{Ar}\)) também com o ângulo \(\theta_1\).
Aplicando a Lei de Snell para a refração do meio \(n_1\) para o ar:
\[n_1 \sin(\theta_1) = n_{\text{ar}} \sin(\theta_5)\]
\[1,30 \sin(40,1^\circ) = 1,00 \cdot \sin(\theta_5)\]
\[\sin(\theta_5) = 1,30 \cdot 0,6441\]
\[\sin(\theta_5) \approx 0,8373\]
\[\theta_5 = \arcsin(0,8373)\]
\[\underline{\theta_5 \approx 56,9^\circ}\]

\vspace{0.2cm}
\textbf{(b) Valor do ângulo \(\theta_4\)}
Para refrações sucessivas através de camadas com interfaces paralelas, podemos aplicar a Lei de Snell de forma encadeada:
\(n_1 \sin\theta_1 = n_2 \sin\theta_2 = n_3 \sin\theta_3 = n_4 \sin\theta_4\).
Isso significa que podemos igualar diretamente o meio inicial ao meio final, ignorando os índices de refração das camadas intermediárias (\(n_2\) e \(n_3\)):
\[n_1 \sin(\theta_1) = n_4 \sin(\theta_4)\]
\[1,30 \sin(40,1^\circ) = 1,45 \sin(\theta_4)\]
\[1,30 \cdot 0,6441 = 1,45 \sin(\theta_4)\]
\[0,8373 = 1,45 \sin(\theta_4)\]
\[\sin(\theta_4) = \frac{0,8373}{1,45}\]
\[\sin(\theta_4) \approx 0,5774\]
\[\theta_4 = \arcsin(0,5774)\]
\[\underline{\theta_4 \approx 35,3^\circ}\]

\subsection{Questão 52 — Guilherme Araújo Floriano}

(a) De acordo com a lei de Snell:

$n_{ar}sen(50^{\circ}) = n_a\sin \theta_a$ e $n_{ar}sen(50^{\circ}) = n_v\sin\theta_v$, na qual osíndices a e v são usados para indicar os raios azul e vermelho. Para $n_{ar} \approx 1,0$, obtemos:

\[
\theta_a = \sin^{-1}{(\frac{\sin50^{\circ}}{1,524})} = 30,176^{\circ} \text{\space\space e \space\space}  
\theta_v = \sin^{-1}{(\frac{\sin50^{\circ}}{1,509})} = 30,507^{\circ} \
\Rightarrow \Delta\theta = 0,33^{\circ}
\]

\[
\]

(b) Como as duas interfaces do vidro com o ar são paralelas, os raios refratados saem do vidro com um ângulo igual ao ângulo de incidência ($50^{\circ}$) independentemente do índice de refração, de modo que a dispersão é $0^{\circ}$.

\subsection{Questão 53 — Arthur Roberto da Silva}
\subsubsection*{(a)}
De acordo com a lei de Snell, temos:
\begin{equation*}
n_1 \operatorname{sen} \theta_1 = n_2 \operatorname{sen} \theta_2
\end{equation*}

Isolando o ângulo de refração $\theta_2$, obtemos:
\begin{equation*}
\theta_2 = \operatorname{sen}^{-1}\left( \frac{n_1}{n_2} \operatorname{sen} \theta_1 \right)
\end{equation*}

Substituindo os valores do problema:
\begin{equation*}
\theta_2 = \operatorname{sen}^{-1}\left[ \frac{1,30}{1,50} \operatorname{sen} 40,0^\circ \right] = \operatorname{sen}^{-1}[ (1,30)(0,4284) ] = 56,8^\circ
\end{equation*}

\subsubsection*{(b)}
Aplicando várias vezes seguidas a lei de Snell para as interfaces das camadas, obtemos a relação de continuidade:
\begin{equation*}
n_1 \operatorname{sen} \theta_1 = n_2 \operatorname{sen} \theta_2 = n_3 \operatorname{sen} \theta_3 = n_4 \operatorname{sen} \theta_4
\end{equation*}

A partir da igualdade direta entre a primeira e a última camada ($n_1 \operatorname{sen} \theta_1 = n_4 \operatorname{sen} \theta_4$), o que nos dá:
\begin{equation*}
\theta_4 = \operatorname{sen}^{-1}\left( \frac{n_1}{n_4} \operatorname{sen} \theta_1 \right) = 35,3^\circ
\end{equation*}

\subsection{Questão 56 — Cicero Rodrigues}

\textbf{56. }\normalfont Arco-íris produzido por gotas quadradas. Suponha que, em um planeta exótico, as gotas de chuva tenham uma seção reta quadrada e caiam sempre com uma face paralela ao solo. A Fig. 33.31 mostra uma dessas gotas, na qual incide um feixe de luz branca com um ângulo de incidência $\theta = 70,0^\circ$ no ponto P. A parte da luz que penetra na gota se propaga até o ponto A, onde parte é refratada de volta para o ar e a outra parte é refletida. A luz refletida chega ao ponto B, onde novamente parte da luz é refratada de volta para o ar e parte é refletida. Qual é a diferença entre os ângulos dos raios de luz vermelha ($n = 1,331$) e de luz azul ($n = 1,343$) que deixam a gota \textbf{(a)} no ponto A e \textbf{(b)} no ponto B? (Se houver uma diferença, um observador externo verá um arco-íris ao observar a luz que sai da gota pelo ponto A ou pelo ponto B.)

\includegraphics[width=0.2\textwidth]{image_14.png}

\textbf{Resolução:}
\vspace{0.3cm}

\textbf{Dados:}
\[\theta = 70,0^\circ \quad \text{(Ângulo de incidência no ponto P)}\]
\[n_v = 1,331 \quad \text{(Índice de refração para luz vermelha)}\]
\[n_a = 1,343 \quad \text{(Índice de refração para luz azul)}\]
\[n_{\text{ar}} = 1,000\]

\vspace{0.2cm}
\textbf{Análise Geométrica Inicial}
A gota tem seção quadrada, portanto faces horizontais e verticais são perpendiculares entre si.
\begin{itemize}
    \item \textbf{No ponto P (face superior):} A normal é vertical. O raio sofre refração com ângulo \(\theta_1\):
    \[\sin(70,0^\circ) = n \sin(\theta_1) \implies \sin(\theta_1) = \frac{\sin(70,0^\circ)}{n}\]
    \item \textbf{No ponto A (face lateral direita):} A normal é horizontal. Por geometria, o ângulo de incidência \(\theta_2\) forma um triângulo retângulo com o ângulo refratado original, logo \(\theta_2 = 90^\circ - \theta_1\). O ângulo de saída no ar é \(\theta_3\):
    \[n \sin(\theta_2) = \sin(\theta_3)\]
    \item \textbf{No ponto B (face inferior):} O raio refletido em A atinge a face inferior, cuja normal volta a ser vertical. O novo ângulo de incidência é \(\theta_4 = 90^\circ - \theta_2\). Como \(\theta_2 = 90^\circ - \theta_1\), concluímos que \(\theta_4 = \theta_1\). O ângulo de saída no ar é \(\theta_5\):
    \[n \sin(\theta_4) = \sin(\theta_5) \implies n \sin(\theta_1) = \sin(\theta_5)\]
\end{itemize}

\vspace{0.2cm}
\textbf{(a) Diferença angular no ponto A (\(\Delta\theta_A\))}
Para a face lateral (Ponto A), deduzimos que:
\[\sin(\theta_3) = n \sin(\theta_2) = n \sin(90^\circ - \theta_1) = n \cos(\theta_1)\]
Usando a identidade \(\cos(\theta_1) = \sqrt{1 - \sin^2(\theta_1)}\) e substituindo \(\sin(\theta_1)\):
\[\sin(\theta_3) = n \sqrt{1 - \left(\frac{\sin 70,0^\circ}{n}\right)^2} = \sqrt{n^2 - \sin^2(70,0^\circ)}\]
Calculando para a luz vermelha (\(n_v = 1,331\)):
\[\sin(\theta_{3v}) = \sqrt{(1,331)^2 - \sin^2(70,0^\circ)} = \sqrt{1,7716 - 0,8830} \approx 0,9426\]
\[\theta_{3v} = \arcsin(0,9426) \approx 70,5^\circ\]
Calculando para a luz azul (\(n_a = 1,343\)):
\[\sin(\theta_{3a}) = \sqrt{(1,343)^2 - \sin^2(70,0^\circ)} = \sqrt{1,8036 - 0,8830} \approx 0,9595\]
\[\theta_{3a} = \arcsin(0,9595) \approx 73,6^\circ\]
A diferença angular no ponto A é:
\[\Delta\theta_A = 73,6^\circ - 70,5^\circ = 3,1^\circ\]
\[\underline{\Delta\theta_A = 3,1^\circ}\]

\vspace{0.2cm}
\textbf{(b) Diferença angular no ponto B (\(\Delta\theta_B\))}
Conforme deduzido na análise geométrica para o ponto B:
\[\sin(\theta_5) = n \sin(\theta_1)\]
Mas, da refração inicial no ponto P, sabemos exatamente que \(n \sin(\theta_1) = \sin(70,0^\circ)\). Portanto:
\[\sin(\theta_5) = \sin(70,0^\circ) \implies \theta_5 = 70,0^\circ\]
Essa relação independe do índice de refração \(n\). Logo, tanto a luz vermelha quanto a luz azul saem pelo ponto B com exatamente \(70,0^\circ\).
A diferença angular no ponto B é nula:
\[\underline{\Delta\theta_B = 0^\circ}\]

\subsection{Questão 59 — Mateus Valentim}
Inserir resolução aqui.

\subsection{Questão 60 — Fábio Rodrigues}
\textbf{Dados}

\[
h = 2,00\,\text{m}
\]

\[
n_{\text{água}} = 1,33 \quad \text{(índice de refração típico da água)}
\]

\[
n_{\text{ar}} = 1,00
\]

A região através da qual o peixe consegue ver o exterior do lago é um círculo na superfície, cujo raio é determinado pelo ângulo crítico da reflexão total interna. Esse fenômeno é conhecido como a "Janela de Snell".

Para ângulos de incidência maiores que o ângulo crítico (\(\theta_c\)), a luz vinda de baixo sofre reflexão total e não sai do lago. Pelo princípio da reversibilidade, a luz vinda de fora do lago só entra na água se refratando dentro de um cone de ângulo de abertura \(2\theta_c\).

O ângulo crítico é dado pela Lei de Snell:

\[
n_{\text{água}} \sin \theta_c = n_{\text{ar}} \sin(90^\circ)
\]

\[
\sin \theta_c = \frac{n_{\text{ar}}}{n_{\text{água}}} = \frac{1,00}{1,33}
\]

\[
\theta_c = \arcsin\left(\frac{1}{1,33}\right) \approx 48,75^\circ
\]

\textbf{(a) Diâmetro da circunferência}

O raio \(r\) da circunferência visível na superfície forma um triângulo retângulo com a profundidade \(h\) e o ângulo \(\theta_c\):

\[
\tan \theta_c = \frac{r}{h}
\]

\[
r = h \tan \theta_c
\]

Substituindo os valores:

\[
r = 2,00 \times \tan(48,75^\circ)
\]

\[
r \approx 2,00 \times 1,140
\]

\[
r \approx 2,28\,\text{m}
\]

O diâmetro \(D\) é o dobro do raio:

\[
D = 2r = 2(2,28)
\]

\[
\boxed{D = 4,56\,\text{m}}
\]

\textbf{(b) Comportamento do diâmetro com o aumento da profundidade}

Analisando a equação encontrada para o diâmetro:

\[
D = 2 h \tan \theta_c
\]

Como o ângulo crítico (\(\theta_c\)) depende apenas dos índices de refração do ar e da água (que são constantes), o termo \(\tan \theta_c\) é constante. Portanto, o diâmetro \(D\) é diretamente proporcional à profundidade \(h\). 

Logo, se o peixe desce para uma profundidade maior (\(h\) aumenta), o diâmetro da circunferência \textbf{aumenta}.

\[
\boxed{\text{Aumenta}}
\]

\subsection{Questão 61 — Rodrigo Albuquerque}

\textbf{ Lentes com raios dados. Um objeto \(O\) está sobre o eixo central de uma lente delgada. Para o problema 61, a Tabela 34-7 fornece:}

\[
p=+75\ cm
\]

\[
n=1,55
\]

\[
r_1=+30\ cm
\]

\[
r_2=-42\ cm
\]

\textbf{Determine:}

\begin{itemize}
\item[(a)] a distância da imagem \(i\);
\item[(b)] a ampliação lateral \(m\);
\item[(c)] se a imagem é real (R) ou virtual (V);
\item[(d)] se é invertida (I) ou não-invertida (NI);
\item[(e)] se está do mesmo lado da lente que o objeto (M) ou do lado oposto (O).
\end{itemize}

\includegraphics[width=0.5\textwidth]{image_8.png}


\vspace{0.5cm}

Primeiro usamos a equação do fabricante de lentes:

\[
\frac{1}{f}
=
(n-1)
\left(
\frac{1}{r_1}
-
\frac{1}{r_2}
\right)
\]

Substituindo:

\[
\frac{1}{f}
=
(1,55-1)
\left(
\frac{1}{30}
-
\frac{1}{-42}
\right)
\]

\[
\frac{1}{f}
=
0,55
\left(
\frac{1}{30}
+
\frac{1}{42}
\right)
\]

\[
\frac{1}{f}
=
0,55(0,0571)
\]

\[
\frac{1}{f}\approx0,0314
\]

\[
f\approx31,8\ cm
\]

\vspace{0.5cm}

Agora usamos a equação das lentes:

\[
\frac{1}{f}
=
\frac{1}{p}
+
\frac{1}{i}
\]

Substituindo:

\[
\frac{1}{31,8}
=
\frac{1}{75}
+
\frac{1}{i}
\]

\[
0,0314
=
0,0133
+
\frac{1}{i}
\]

\[
\frac{1}{i}=0,0181
\]

\[
i\approx55,2\ cm
\]

\vspace{0.5cm}

\textbf{(b) Ampliação lateral}

\[
m=-\frac{i}{p}
\]

\[
m=-\frac{55,2}{75}
\]

\[
m\approx-0,736
\]

\vspace{0.5cm}

\textbf{(c) Real ou virtual}

Como:

\[
i>0
\]

a imagem é real (R).

\vspace{0.5cm}

\textbf{(d) Invertida ou não-invertida}

Como:

\[
m<0
\]

a imagem é invertida (I).

\vspace{0.5cm}

\textbf{(e) Lado da imagem}

Imagens reais em lentes ficam do lado oposto da lente em relação ao objeto.

Logo:

\[
\text{lado} = O
\]

\subsection{Questão 62 — Dávyla Maria}
Inserir resolução aqui.

\subsection{Questão 63 — Kaylane Castro}


\subsection*{Enunciado}

A Fig. 33-62 mostra uma fibra ótica simplificada: um núcleo de plástico ($n_1 = 1,58$) envolvido por um revestimento de plástico com índice de refração menor ($n_2 = 1,53$).

Um raio luminoso incide em uma das extremidades da fibra com um ângulo $\theta$. O raio deve sofrer reflexão interna total no ponto $A$, onde atinge a interface núcleo-revestimento.

Determine o maior valor de $\theta$ para o qual existe reflexão interna total no ponto $A$.

% ---------------------------------------------------------

\subsection*{Dados}

\[
n_1 = 1,58
\]

\[
n_2 = 1,53
\]

\[
n_{\text{ar}} \approx 1,00
\]

% ---------------------------------------------------------

\subsection*{Análise Física}

O problema envolve dois fenômenos ópticos:

\begin{enumerate}
    \item Refração da luz na entrada da fibra;
    \item Reflexão interna total na interface núcleo-revestimento.
\end{enumerate}

Na entrada da fibra, aplica-se a Lei de Snell:

\[
n_{\text{ar}}\sin\theta = n_1\sin\theta_r
\]

onde:

\begin{itemize}
    \item $\theta$ é o ângulo de incidência no ar;
    \item $\theta_r$ é o ângulo refratado no núcleo.
\end{itemize}

Na parede interna da fibra, o raio sofre reflexão interna total.

O ângulo de incidência na interface é:

\[
\phi = 90^\circ - \theta_r
\]

A condição limite para reflexão interna total ocorre quando:

\[
\phi = \theta_c
\]

onde $\theta_c$ é o ângulo crítico.

% ---------------------------------------------------------

\subsection*{Ângulo Crítico}

Para a interface núcleo-revestimento:

\[
\sin\theta_c = \frac{n_2}{n_1}
\]

Substituindo os valores:

\[
\sin\theta_c =
\frac{1,53}{1,58}
\]

\[
\sin\theta_c \approx 0,96835
\]

Como:

\[
\phi = 90^\circ - \theta_r
\]

então:

\[
\sin\phi = \cos\theta_r
\]

Logo:

\[
\cos\theta_r = 0,96835
\]

% ---------------------------------------------------------

\subsection*{Cálculo de \texorpdfstring{$\sin\theta_r$}{sen(theta r)}}

Usando a identidade trigonométrica:

\[
\sin^2\theta_r + \cos^2\theta_r = 1
\]

temos:

\[
\sin\theta_r =
\sqrt{1 - \cos^2\theta_r}
\]

Substituindo:

\[
\sin\theta_r =
\sqrt{1 - (0,96835)^2}
\]

\[
\sin\theta_r =
\sqrt{1 - 0,93770}
\]

\[
\sin\theta_r =
\sqrt{0,06230}
\]

\[
\sin\theta_r \approx 0,2496
\]

% ---------------------------------------------------------

\subsection*{Aplicando a Lei de Snell}

Agora aplicamos a Lei de Snell na entrada da fibra:

\[
n_{\text{ar}}\sin\theta =
n_1\sin\theta_r
\]

Substituindo os valores:

\[
1,00\sin\theta =
1,58(0,2496)
\]

\[
\sin\theta \approx 0,3944
\]

Aplicando arco-seno:

\[
\theta = \arcsin(0,3944)
\]

\[
\theta \approx 23,2^\circ
\]

% ---------------------------------------------------------

\subsection*{Resposta}

\[
\boxed{
\theta_{\text{máx}} \approx 23,2^\circ
}
\]



\subsection{Questão 64 — Ana Alicy}
Inserir resolução aqui.

\subsection{Questão 65 — Ana Alicy}
Inserir resolução aqui.

\subsection{Questão 66 — Kaylane Castro}

\subsection*{Enunciado}

Suponha que o ângulo do vértice superior do prisma de vidro da Fig. 33-57 seja

\[
\phi = 60,0^\circ
\]

e que o índice de refração do vidro seja

\[
n = 1,60
\]

\begin{enumerate}
    \item[\textbf{(a)}] Qual é o menor ângulo de incidência $\theta$ para o qual um raio pode entrar na face esquerda do prisma e sair na face direita?
    
    \item[\textbf{(b)}] Qual deve ser o ângulo de incidência $\theta$ para que o raio saia do prisma com o mesmo ângulo $\theta$ com que entrou?
\end{enumerate}

% ---------------------------------------------------------

\subsection*{Dados}

\[
\phi = 60,0^\circ
\]

\[
n = 1,60
\]

\[
n_{\text{ar}} \approx 1,00
\]

% ---------------------------------------------------------

\subsection*{Item (a)}

O menor ângulo de incidência ocorre quando o raio atinge a segunda face exatamente na condição limite de reflexão interna total.

Nesse caso, o raio emerge tangenciando a superfície, isto é:

\[
\theta_{\text{saída}} = 90^\circ
\]

% ---------------------------------------------------------

\subsubsection*{1. Ângulo crítico na segunda face}

Aplicando a Lei de Snell:

\[
n\sin\theta_{i2}
=
n_{\text{ar}}\sin90^\circ
\]

\[
1,60\sin\theta_{i2}
=
1,00
\]

\[
\sin\theta_{i2}
=
\frac{1}{1,60}
\]

\[
\sin\theta_{i2}
=
0,625
\]

\[
\theta_{i2}
=
\arcsin(0,625)
\]

\[
\theta_{i2}
\approx 38,68^\circ
\]

% ---------------------------------------------------------

\subsubsection*{2. Relação geométrica do prisma}

Para um prisma:

\[
\phi = \theta_{r1} + \theta_{i2}
\]

Logo:

\[
\theta_{r1}
=
60,0^\circ - 38,68^\circ
\]

\[
\theta_{r1}
\approx 21,32^\circ
\]

% ---------------------------------------------------------

\subsubsection*{3. Aplicando a Lei de Snell na entrada}

Na entrada do prisma:

\[
n_{\text{ar}}\sin\theta
=
n\sin\theta_{r1}
\]

\[
1,00\sin\theta
=
1,60\sin(21,32^\circ)
\]

\[
\sin\theta
=
1,60(0,3636)
\]

\[
\sin\theta
\approx 0,5818
\]

\[
\theta
=
\arcsin(0,5818)
\]

\[
\theta
\approx 35,6^\circ
\]

% ---------------------------------------------------------

\subsection*{ Item (b)}

Para que o raio saia com o mesmo ângulo com que entrou, o caminho óptico deve ser simétrico.

Assim:

\[
\theta_{r1} = \theta_{i2}
\]

Usando a relação geométrica do prisma:

\[
\phi = \theta_{r1} + \theta_{i2}
\]

\[
60,0^\circ = 2\theta_{r1}
\]

\[
\theta_{r1}
=
30,0^\circ
\]

% ---------------------------------------------------------

\subsubsection*{Aplicando a Lei de Snell}

\[
n_{\text{ar}}\sin\theta
=
n\sin\theta_{r1}
\]

\[
1,00\sin\theta
=
1,60\sin30^\circ
\]

\[
\sin\theta
=
1,60(0,5)
\]

\[
\sin\theta
=
0,80
\]

\[
\theta
=
\arcsin(0,80)
\]

\[
\theta
\approx 53,1^\circ
\]

% ---------------------------------------------------------

\subsection*{Respostas}

\[
\text{\textbf{(a)}} \quad
\boxed{
\theta \approx 35,6^\circ
}
\]

\[
\text{\textbf{(b)}} \quad
\boxed{
\theta \approx 53,1^\circ
}
\]
```

\subsection{Questão 67 — Dávyla Maria}
Inserir resolução aqui.

\subsection{Questão 68 — }
Inserir resolução aqui.

\subsection{Questão 69 — Kaylane Castro}
Inserir resolução aqui.

\subsection{Questão 80 — Guilherme Araújo Floriano}

(a)

Fazendo $v = c$ na relação $kv = \omega = 2\pi f$, obtemos:

\[
f = \frac{kc}{2\pi} = \frac{4 m^{-1}*3*10^8 m/s}{6,28} = 1,91 * 10^8   
\]

(b)
\[
E_{rms} = \frac{E_m}{\sqrt2} = \frac{B_mc}{\sqrt2} = (85,8 * 10^9 T)(3*10^8 m/s)/(1,414) = 18,2 V/m
\]

(c)
\[
I = (E_{rms})^2/c\mu_{0} = 
(18,2 V/m)^2/(3 * 10^8 m/s)(4\pi * 10^{-7} H/n)
= 0,878 W/m^2
\]

\subsection{Questão 84 — Fábio Rodrigues}
\textbf{Dados}

\[
\text{Número de filtros: } 4
\]

\[
\text{Ângulo entre filtros vizinhos: } \theta = 30^\circ
\]

\[
\text{Intensidade inicial (luz não-polarizada): } I_0
\]

Quando a luz não-polarizada atravessa o primeiro filtro polarizador, sua intensidade é reduzida à metade:

\[
I_1 = \frac{1}{2} I_0
\]

Para os filtros seguintes, aplicamos a Lei de Malus. A intensidade transmitida por um filtro polarizador é dada por \(I = I_i \cos^2 \theta\), onde \(I_i\) é a intensidade incidente no filtro e \(\theta\) é o ângulo entre a direção de polarização da luz incidente e o eixo de transmissão do filtro.

Após o segundo filtro:

\[
I_2 = I_1 \cos^2(30^\circ) = \frac{1}{2} I_0 \cos^2(30^\circ)
\]

Após o terceiro filtro:

\[
I_3 = I_2 \cos^2(30^\circ) = \frac{1}{2} I_0 \cos^4(30^\circ)
\]

Após o quarto filtro:

\[
I_4 = I_3 \cos^2(30^\circ) = \frac{1}{2} I_0 \cos^6(30^\circ)
\]

A fração da luz incidente que é transmitida pelo conjunto é \(\frac{I_4}{I_0}\):

\[
\frac{I_4}{I_0} = \frac{1}{2} \cos^6(30^\circ)
\]

Sabendo que \(\cos(30^\circ) = \frac{\sqrt{3}}{2}\):

\[
\frac{I_4}{I_0} = \frac{1}{2} \left( \frac{\sqrt{3}}{2} \right)^6
\]

\[
\frac{I_4}{I_0} = \frac{1}{2} \left( \frac{3}{4} \right)^3
\]

\[
\frac{I_4}{I_0} = \frac{1}{2} \left( \frac{27}{64} \right)
\]

\[
\frac{I_4}{I_0} = \frac{27}{128}
\]

Em formato decimal, temos:

\[
\frac{I_4}{I_0} \approx 0,211
\]

\[
\boxed{\frac{I_4}{I_0} \approx 0,211}
\]

\subsection{Questão 85 — Mateus Valentim}
\subsubsection{Resolução}

Ao atravessar o primeiro polarizador, a luz não-polarizada tem sua intensidade
reduzida à metade:

\[
I_1=\frac{I_0}{2}.
\]

Ao atravessar o segundo polarizador, vale a lei de Malus:

\[
I_2=I_1\cos^2\theta.
\]

Queremos

\[
I_2=\frac{I_0}{3}.
\]

Assim,

\[
\frac{I_0}{3}
=
\frac{I_0}{2}\cos^2\theta.
\]

Cancelando \(I_0\),

\[
\cos^2\theta=\frac{2}{3}.
\]

Logo,

\[
\theta=\cos^{-1}\left(\sqrt{\frac{2}{3}}\right)
=35{,}3^\circ.
\]

Portanto, o ângulo entre as direções de polarização dos filtros deve ser

\[
\boxed{35{,}3^\circ}.
\]

\subsection{Questão 86 — Cicero Rodrigues}

\textbf{86. }\normalfont Um feixe de luz não polarizada atravessa um conjunto de quatro filtros polarizadores orientados de tal forma que o ângulo entre as direções de polarização de filtros vizinhos é 30$^\circ$. Que fração da luz incidente é transmitida pelo conjunto?

\textbf{Resolução:}
\vspace{0.3cm}

\textbf{Dados:}
\[I_0 \quad \text{(Intensidade inicial da luz não polarizada)}\]
\[\theta = 30^\circ \quad \text{(Ângulo entre filtros vizinhos)}\]
\[n = 4 \quad \text{(Número total de filtros)}\]

\vspace{0.2cm}
\textbf{Transmissão pelo primeiro filtro (\(I_1\))}
Quando a luz não polarizada atravessa o primeiro polarizador, exatamente metade da sua intensidade é absorvida, independentemente da orientação do filtro:
\[I_1 = \frac{1}{2}I_0\]

\vspace{0.2cm}
\textbf{Transmissão pelos filtros subsequentes (Lei de Malus)}
A luz que sai do primeiro filtro agora é linearmente polarizada. Para os próximos filtros, aplicamos a Lei de Malus: \(I = I_{\text{anterior}} \cos^2(\theta)\).
Como temos mais 3 filtros (do 1º para o 2º, do 2º para o 3º, e do 3º para o 4º), aplicamos a lei sucessivamente:
\[I_2 = I_1 \cos^2(30^\circ)\]
\[I_3 = I_2 \cos^2(30^\circ) = I_1 [\cos^2(30^\circ)]^2\]
\[I_4 = I_3 \cos^2(30^\circ) = I_1 [\cos^2(30^\circ)]^3\]

\vspace{0.2cm}
\textbf{Cálculo da fração final (\(I_4/I_0\))}
Substituindo \(I_1 = \frac{1}{2}I_0\) na equação final:
\[I_4 = \frac{1}{2}I_0 [\cos^2(30^\circ)]^3\]
Sabendo que \(\cos(30^\circ) = \frac{\sqrt{3}}{2}\):
\[I_4 = \frac{1}{2}I_0 \left[ \left( \frac{\sqrt{3}}{2} \right)^2 \right]^3\]
\[I_4 = \frac{1}{2}I_0 \left[ \frac{3}{4} \right]^3\]
\[I_4 = \frac{1}{2}I_0 \left( \frac{27}{64} \right)\]
\[I_4 = \frac{27}{128} I_0\]
A fração da luz transmitida em relação à incidente é:
\[\frac{I_4}{I_0} = \frac{27}{128} \approx 0,211\]
\[\underline{\text{Fração} = \frac{27}{128} \quad \text{ou} \quad 21,1\%}\]

\subsection{Questão 87 — Guilherme Araújo Floriano}

De acordo com a Figura \ref{fig:example_image}, $n_{max} = 1,470$ para $\lambda = 400 nm$ e $n_{min} = 1,456$ para
$\lambda = 700 nm$.

\begin{figure}[H]
    \centering
    \includegraphics[width=0.5\textwidth]{fig33-18.png}
    \caption{}
    \label{fig:example_image}
    
\end{figure}

a- Dado $\theta_B = \tan^{-1}\frac{n_2}{n_1}$, temos que: $\theta_{B,max} = \tan^{-1}n_{max}=\tan^{-1}1,470 = 55,8^{\circ}$

b- $\theta_{B,min} = \tan^{-1}n_{min}=\tan^{-1}1,456 = 55,5^{\circ}$

\subsection{Questão 89 — Arthur Roberto da Silva}
\noindent 81. (a) A direção de polarização é definida pelo campo elétrico, que é perpendicular ao campo magnético e à direção de propagação da onda. A função dada mostra que o campo magnético é paralelo ao eixo $x$ (por causa do índice da amplitude $B$) e que a onda está se propagando no sentido negativo do eixo $y$ (por causa do argumento da função seno). Assim, o campo elétrico é paralelo ao eixo $z$ e a direção da polarização da luz é a direção do eixo $z$.

\subsection{Questão 90 — Arthur Roberto da Silva}
De acordo com as equações de polarização para luz inicialmente não-polarizada que passa por uma sequência de filtros polarizadores (Lei de Malus), a razão entre a intensidade final ($I_{\text{final}}$) e a intensidade inicial ($I_0$) é dada por:

\begin{equation*}
\frac{I_{\text{final}}}{I_0} = \frac{(I_0 / 2)(\cos^2 45^\circ)^2}{I_0} = \frac{1}{8} = 0,125
\end{equation*}

\subsection*{Resposta Final}
A razão entre as intensidades é \textbf{$0,125$} (ou \textbf{$12,5\%$}).


\vspace{3cm} % Espaço aproximado equivalente à imagem

\noindent \textbf{24} \hspace{1cm} SOLUÇÕES DOS PROBLEMAS

\vspace{1cm}

\noindent (b) Como $k = 1,57 \times 10^7 \, \text{/m}$, $\lambda = 2\pi/k = 4,0 \times 10^{-7} \, \text{m}$, o que nos dá
\[
f = c/\lambda = 7,5 \times 10^{14} \, \text{Hz}.
\]

\noindent (c) De acordo com a Eq. 33-26, temos:
\[
I = \frac{E_{\text{rms}}^2}{c\mu_0} = \frac{E_m^2}{2c\mu_0} = \frac{(cB_m)^2}{2c\mu_0} = \frac{cB_m^2}{2\mu_0} = \frac{(3 \times 10^8 \, \text{m/s})(4,0 \times 10^{-6} \, \text{T})^2}{2(4\pi \times 10^{-7} \, \text{H/m})} = 1,9 \, \text{kW/m}^2.
\]

\subsection{Questão 91 — Rodrigo Albuquerque }

\textbf{91 Uma pessoa com um ponto próximo \(P_n\) de \(25\ cm\) observa um dedal através de uma lente de aumento simples com uma distância focal de \(10\ cm\), mantendo a lente perto do olho. Determine a ampliação angular do dedal quando ele é posicionado de tal forma que sua imagem aparece:}

\begin{itemize}
\item[(a)] em \(P_n\);
\item[(b)] no infinito.
\end{itemize}

\vspace{0.5cm}

Dados:

\[
P_n=25\ cm
\]

\[
f=10\ cm
\]

\vspace{0.5cm}

\textbf{(a) Imagem em \(P_n\)}

Para uma lente de aumento:

\[
M=1+\frac{P_n}{f}
\]

Substituindo:

\[
M=1+\frac{25}{10}
\]

\[
M=1+2,5
\]

\[
M=3,5
\]

\vspace{0.5cm}

\textbf{(b) Imagem no infinito}

Quando a imagem está no infinito:

\[
M=\frac{P_n}{f}
\]

Substituindo:

\[
M=\frac{25}{10}
\]

\[
M=2,5
\]

% ================= CAPÍTULO 34 =================

\section{Capítulo 34}

\subsection{Questão 3 — Arthur Roberto da Silva}
\noindent 4. No momento em que $S$ consegue ver $B$, os raios luminosos provenientes de $B$ estão sendo refletidos pela borda do espelho em direção a $S$. Nesse caso, o ângulo de reflexão é $45^\circ$, já que uma reta traçada de $S$ até a borda do espelho faz um ângulo de $45^\circ$ com a parede. De acordo com a lei de reflexão de espelhos planos,
\[
\frac{x}{d/2} = \tan 45^\circ \implies 1 = \frac{x}{d/2} \implies x = \frac{d}{2} = \frac{3,0 \, \text{m}}{2} = 1,5 \, \text{m}.
\]

\subsection{Questão 4 — Guilherme Araújo Floriano}

A intensidade da luz produzida por uma fonte pontual varia com o quadrado da distância da fonte. Antes da introdução do espelho, a intensidade da luz no centro da tela é dada por $I_P = A/d^2$, em que $A$ é uma constante. Depois que o espelho é introduzido, a intensidade da luz é a soma da luz que chega diretamente à tela, com a mesma intensidade $I_P$ de antes, com a luz refletida. Como a luz refletida parece ter sido produzida por uma fonte pontual situada a uma distância $d$ atrás do espelho, a distância entre a imagem da fonte e a tela é $3d$ e sua contribuiçãopara a intensidade da luz no centro da tela é:

$I_r = \frac{A}{(3d)^2} = \frac{A}{9d} = \frac{I_P}{9}$

A intensidade total da luz no centro da tela é, portanto, 

$I = I_P + I_r = I_P + \frac{I_P}{9} = \frac{10}{9}I_P$

e a razão entre a nova intensidade e a intensidade antiga é $I/I_P = 10/9 = 1,11$

\subsection{Questão 5 — Cicero Rodrigues}
\textbf{5. }\normalfont A Fig. 34.10 mostra uma lâmpada pendurada a uma distância $d_1 = 250$ cm acima da superfície da água de uma piscina na qual a profundidade da água é de $d_2 = 200$ cm. O fundo da piscina é um espelho. A que distância da superfície do espelho está a imagem da lâmpada? (Sugestão: Suponha que os raios não se desviam muito de uma reta vertical que passa pela lâmpada, e use a aproximação, válida para pequenos ângulos, de que $\text{sen}\,\theta \approx \tan\theta \approx \theta$.)

\includegraphics[width=0.2\textwidth]{image_15.png}

\textbf{Resolução:}
\vspace{0.3cm}

\textbf{Dados:}
\[d_1 = 250\;cm \quad \text{(Distância real da lâmpada à superfície)}\]
\[d_2 = 200\;cm \quad \text{(Profundidade da água)}\]
\[n_{\text{ar}} = 1,00\]
\[n_{\text{água}} \approx 1,33 \quad \text{(ou } 4/3 \text{)}\]

\vspace{0.2cm}
\textbf{(1) Posição aparente da lâmpada (\(d_1'\))}
A sugestão do enunciado para usar a aproximação de ângulos pequenos é a dedução matemática por trás da fórmula clássica de profundidade/altura aparente. Como a luz vai do ar para a água, um observador (ou o espelho) submerso "vê" a lâmpada mais afastada do que ela realmente está. A distância aparente acima da superfície é:
\[d_1' = d_1 \left( \frac{n_{\text{água}}}{n_{\text{ar}}} \right)\]
Substituindo os valores (usando o padrão \(1,33\)):
\[d_1' = 250 \left( \frac{1,33}{1,00} \right)\]
\[d_1' = 332,5\;cm\]
\textit{(Nota: Se usarmos a fração exata \(4/3\) para o índice da água, o valor será de \(333,33\;cm\)).}

\vspace{0.2cm}
\textbf{(2) Distância total do objeto aparente ao espelho (\(D\))}
Para o espelho localizado no fundo da piscina, a lâmpada parece estar emanando do ponto \(d_1'\). A distância total \(D\) desse "objeto virtual" até a superfície do espelho é a soma da altura aparente com a profundidade da água:
\[D = d_1' + d_2\]
\[D = 332,5 + 200\]
\[D = 532,5\;cm\]

\vspace{0.2cm}
\textbf{(3) Distância da imagem em relação ao espelho}
A propriedade fundamental de um espelho plano é formar uma imagem a uma distância (atrás dele) exatamente igual à distância em que o objeto se encontra (à sua frente). Portanto, a distância da imagem até a superfície do espelho é o próprio valor de \(D\). 
Arredondando para os três algarismos significativos convencionais do livro:
\[\underline{d_{\text{imagem}} = 533\;cm}\]

\subsection{Questão 17 — Mateus Valentim}
\subsubsection{Resolução}

Na dupla fenda, a separação angular entre franjas é proporcional ao comprimento
de onda:

\[
\Delta\theta=\frac{\lambda}{d}.
\]

Como \(d\) não muda, para que a separação angular seja \(10{,}0\%\) maior, o
comprimento de onda também deve ser \(10{,}0\%\) maior.

Assim,

\[
\lambda'=1{,}10\lambda.
\]

Logo,

\[
\lambda'=1{,}10(589\,\text{nm})
=648\,\text{nm}.
\]

Portanto, o comprimento de onda deve ser

\[
\boxed{648\,\text{nm}}.
\]

\subsection{Questão 18 — Fábio Rodrigues}
\textbf{Dados}

Pela análise do gráfico, o eixo horizontal possui 4 subdivisões entre \(0\) e \(p_s = 10,0\,\text{cm}\). Portanto, cada subdivisão vale:

\[
\frac{10,0}{4} = 2,5\,\text{cm}
\]

Um ponto exato no gráfico é a intersecção da curva com a linha de \(m = 0,5\), que ocorre na segunda subdivisão do eixo horizontal. Portanto:

Para \(m = 0,5\), temos \(p = 2 \times 2,5 = 5,0\,\text{cm}\).

\textbf{Passo 1: Determinar a distância focal (\(f\))}

Sabemos que a ampliação lateral de um espelho esférico é dada por:

\[
m = -\frac{i}{p} \implies i = -mp
\]

Substituindo essa relação na equação dos pontos conjugados de Gauss (\(\frac{1}{p} + \frac{1}{i} = \frac{1}{f}\)):

\[
\frac{1}{p} - \frac{1}{mp} = \frac{1}{f}
\]

\[
\frac{m - 1}{mp} = \frac{1}{f}
\]

\[
f = \frac{mp}{m - 1}
\]

Substituindo os valores do ponto conhecido (\(p = 5,0\,\text{cm}\) e \(m = 0,5\)):

\[
f = \frac{0,5 \times 5,0}{0,5 - 1}
\]

\[
f = \frac{2,5}{-0,5}
\]

\[
f = -5,0\,\text{cm}
\]

\textit{(O sinal negativo indica tratar-se de um espelho convexo).}

\textbf{Passo 2: Calcular a ampliação para \(p = 21\,\text{cm}\)}

Podemos reorganizar a equação deduzida anteriormente para isolar \(m\):

\[
\frac{m - 1}{m} = \frac{p}{f}
\]

\[
1 - \frac{1}{m} = \frac{p}{f}
\]

\[
\frac{1}{m} = 1 - \frac{p}{f} = \frac{f - p}{f}
\]

\[
m = \frac{f}{f - p}
\]

Substituindo \(f = -5,0\,\text{cm}\) e a nova distância \(p = 21\,\text{cm}\):

\[
m = \frac{-5,0}{-5,0 - 21}
\]

\[
m = \frac{-5,0}{-26,0}
\]

\[
m = \frac{5}{26}
\]

\[
m \approx 0,192
\]

\[
\boxed{m \approx 0,19}
\]

\subsection{Questão 19 — Rodrigo Albuquerque}

\textbf{ Mais espelhos. Um objeto \(O\) está sobre o eixo central de um espelho esférico. Para o problema 19, a Tabela 34-4 fornece:}

\[
\text{Espelho côncavo}
\]

\[
f=20\ cm
\]

\[
p=+10\ cm
\]

\textbf{Determine os dados que faltam:}

\begin{itemize}
\item[(c)] o raio de curvatura \(r\);
\item[(e)] a distância da imagem \(i\);
\item[(f)] a ampliação lateral \(m\);
\item[(g)] se a imagem é real (R) ou virtual (V);
\item[(h)] se é invertida (I) ou não-invertida (NI);
\item[(i)] se a imagem está do mesmo lado do espelho que o objeto (M) ou do lado oposto (O).
\end{itemize}

\includegraphics[width=0.5\textwidth]{image_9.png}

\vspace{0.5cm}

\textbf{(c) Raio de curvatura}

Para espelhos esféricos:

\[
r=2f
\]

\[
r=2(20)
\]

\[
r=40\ cm
\]

\vspace{0.5cm}

\textbf{(e) Distância da imagem}

Equação dos espelhos:

\[
\frac{1}{f}
=
\frac{1}{p}
+
\frac{1}{i}
\]

Substituindo:

\[
\frac{1}{20}
=
\frac{1}{10}
+
\frac{1}{i}
\]

\[
\frac{1}{i}
=
\frac{1}{20}-\frac{1}{10}
\]

\[
\frac{1}{i}
=
-\frac{1}{20}
\]

\[
i=-20\ cm
\]

\vspace{0.5cm}

\textbf{(f) Ampliação lateral}

\[
m=-\frac{i}{p}
\]

\[
m=-\frac{-20}{10}
\]

\[
m=2,0
\]

\vspace{0.5cm}

\textbf{(g) Real ou virtual}

Como:

\[
i<0
\]

a imagem é virtual (V).

\vspace{0.5cm}

\textbf{(h) Invertida ou não-invertida}

Como:

\[
m>0
\]

a imagem é não-invertida (NI).

\vspace{0.5cm}

\textbf{(i) Lado da imagem}

Imagens virtuais em espelhos ficam atrás do espelho.

Logo, a imagem está do lado oposto (O).

\subsection{Questão 20 — Dávyla Maria}
Inserir resolução aqui.

\subsection{Questão 21 — Kaylane Castro}


\subsection*{Enunciado}

Considere um espelho esférico com:

\[
r = -40\,\text{cm}
\]

e imagem formada em:

\[
i = -10\,\text{cm}
\]

Determine:

\begin{enumerate}
    \item[\textbf{(a)}] O tipo de espelho;
    \item[\textbf{(b)}] A distância focal;
    \item[\textbf{(d)}] A distância do objeto;
    \item[\textbf{(f)}] A ampliação lateral;
    \item[\textbf{(g)}] Se a imagem é real ou virtual;
    \item[\textbf{(h)}] Se a imagem é invertida ou não-invertida;
    \item[\textbf{(i)}] O lado em que a imagem se encontra.
\end{enumerate}

% ---------------------------------------------------------

\subsection*{Dados}

\[
r = -40\,\text{cm}
\]

\[
i = -10\,\text{cm}
\]

% ---------------------------------------------------------

\subsection*{Tipo de Espelho}

Como o raio de curvatura é negativo:

\[
r < 0
\]

o espelho é classificado como:

\[
\boxed{\text{Espelho Convexo}}
\]

% ---------------------------------------------------------

\subsection*{Distância Focal}

A distância focal de um espelho esférico é:

\[
f = \frac{r}{2}
\]

Substituindo os valores:

\[
f = \frac{-40}{2}
\]

\[
f = -20\,\text{cm}
\]

% ---------------------------------------------------------

\subsection*{Distância do Objeto}

Utilizando a equação dos espelhos esféricos:

\[
\frac{1}{f} = \frac{1}{p} + \frac{1}{i}
\]

Isolando $\frac{1}{p}$:

\[
\frac{1}{p} = \frac{1}{f} - \frac{1}{i}
\]

Substituindo os valores:

\[
\frac{1}{p}
=
\frac{1}{-20}
-
\frac{1}{-10}
\]

\[
\frac{1}{p}
=
-\frac{1}{20}
+
\frac{1}{10}
\]

\[
\frac{1}{p}
=
\frac{-1 + 2}{20}
\]

\[
\frac{1}{p}
=
\frac{1}{20}
\]

Logo:

\[
p = 20\,\text{cm}
\]

% ---------------------------------------------------------

\subsection*{Ampliação Lateral}

A ampliação lateral é dada por:

\[
m = -\frac{i}{p}
\]

Substituindo os valores:

\[
m
=
-\frac{-10}{20}
\]

\[
m = +0,50
\]

% ---------------------------------------------------------

\subsection*{Natureza da Imagem}

Como:

\[
i < 0
\]

a imagem é:

\[
\boxed{\text{Virtual}}
\]

% ---------------------------------------------------------

\subsection*{Orientação da Imagem}

Como:

\[
m > 0
\]

a imagem é:

\[
\boxed{\text{Não-Invertida}}
\]

% ---------------------------------------------------------

\subsection*{Lado da Imagem}

Para espelhos esféricos:

\begin{itemize}
    \item $i > 0$ $\rightarrow$ imagem real;
    \item $i < 0$ $\rightarrow$ imagem virtual.
\end{itemize}

Em ambos os casos, imagem e objeto encontram-se do mesmo lado da superfície refletora.

Como:

\[
i = -10\,\text{cm}
\]

a imagem está no lado:

\[
\boxed{\text{M}}
\]

% ---------------------------------------------------------

\subsection*{Tabela Final}

\begin{center}
\begin{tabular}{ccccccccc}
\hline

(a) Tipo & (b) $f$ & (c) $r$ & (d) $p$ & (e) $i$
& (f) $m$ & (g) R/V & (h) I/NI & (i) Lado \\

\hline

Convexo & $-20$ & $-40$ & $+20$ & $-10$
& $+0,50$ & V & NI & M \\

\hline
\end{tabular}
\end{center}

% ---------------------------------------------------------

\subsection*{Respostas}

\[
\boxed{
\text{Espelho Convexo}
}
\]

\[
\boxed{
f = -20\,\text{cm}
}
\]

\[
\boxed{
p = +20\,\text{cm}
}
\]

\[
\boxed{
m = +0,50
}
\]

\[
\boxed{
\text{Imagem Virtual e Não-Invertida}
}
\]

\[
\boxed{
\text{Lado M}
}
\]

\subsection{Questão 22 — Ana Alicy}
Inserir resolução aqui.

\subsection{Questão 23 — Ana Alicy}
Inserir resolução aqui.

\subsection{Questão 24 — Kaylane Castro}

\subsection*{Enunciado}

Um objeto está localizado sobre o eixo principal de um espelho esférico.

São fornecidos os seguintes dados:

\[
p = +24\,\text{cm}
\]

\[
|m| = 0,50
\]

A imagem formada é invertida.

Determine:

\begin{enumerate}
    \item[\textbf{(a)}] O tipo de espelho;
    \item[\textbf{(b)}] A distância focal;
    \item[\textbf{(c)}] O raio de curvatura;
    \item[\textbf{(e)}] A distância da imagem;
    \item[\textbf{(f)}] O sinal correto da ampliação lateral;
    \item[\textbf{(g)}] Se a imagem é real ou virtual;
    \item[\textbf{(i)}] O lado em que a imagem se encontra.
\end{enumerate}

% ---------------------------------------------------------

\subsection*{Dados}

\[
p = +24\,\text{cm}
\]

\[
|m| = 0,50
\]

Imagem invertida.

% ---------------------------------------------------------

\subsection*{Ampliação Lateral}

Imagens invertidas possuem ampliação negativa.

Assim:

\[
m = -0,50
\]

% ---------------------------------------------------------

\subsection*{Distância da Imagem}

A ampliação lateral é dada por:

\[
m = -\frac{i}{p}
\]

Substituindo os valores:

\[
-0,50 = -\frac{i}{24}
\]

Multiplicando ambos os lados por $24$:

\[
i = 0,50 \times 24
\]

\[
i = +12\,\text{cm}
\]

% ---------------------------------------------------------

\subsection*{Distância Focal}

Aplicando a equação dos espelhos esféricos:

\[
\frac{1}{f} = \frac{1}{p} + \frac{1}{i}
\]

Substituindo os valores:

\[
\frac{1}{f}
=
\frac{1}{24}
+
\frac{1}{12}
\]

\[
\frac{1}{f}
=
\frac{1 + 2}{24}
\]

\[
\frac{1}{f}
=
\frac{3}{24}
\]

\[
\frac{1}{f}
=
\frac{1}{8}
\]

Logo:

\[
f = +8,0\,\text{cm}
\]

% ---------------------------------------------------------

\subsection*{Raio de Curvatura}

O raio de curvatura é:

\[
r = 2f
\]

Substituindo o valor encontrado:

\[
r = 2(8,0)
\]

\[
r = +16\,\text{cm}
\]

% ---------------------------------------------------------

\subsection*{Tipo de Espelho}

Como:

\[
f > 0
\]

o espelho é:

\[
\boxed{\text{Côncavo}}
\]

% ---------------------------------------------------------

\subsection*{Natureza da Imagem}

Como:

\[
i > 0
\]

a imagem é:

\[
\boxed{\text{Real}}
\]

% ---------------------------------------------------------

\subsection*{Lado da Imagem}

Imagens reais em espelhos esféricos formam-se na frente do espelho, ou seja, no mesmo lado do objeto.

Portanto:

\[
\boxed{\text{Lado M}}
\]

% ---------------------------------------------------------

\subsection*{Tabela Final}

\begin{center}
\begin{tabular}{ccccccccc}
\hline

(a) Tipo & (b) $f$ & (c) $r$ & (d) $p$ & (e) $i$
& (f) $m$ & (g) R/V & (h) I/NI & (i) Lado \\

\hline

Côncavo & $+8,0$ & $+16$ & $+24$ & $+12$
& $-0,50$ & R & I & M \\

\hline
\end{tabular}
\end{center}

% ---------------------------------------------------------

\subsection*{Respostas}

\[
\boxed{
f = +8,0\,\text{cm}
}
\]

\[
\boxed{
r = +16\,\text{cm}
}
\]

\[
\boxed{
i = +12\,\text{cm}
}
\]

\[
\boxed{
m = -0,50
}
\]

\[
\boxed{
\text{Espelho Côncavo}
}
\]

\[
\boxed{
\text{Imagem Real}
}
\]

\[
\boxed{
\text{Lado M}
}
\]
```


\subsection{Questão 25 — Dávyla Maria}
Inserir resolução aqui.

\subsection{Questão 26 — Rodrigo Albuquerque}
\textbf{Mais espelhos. Um objeto \(O\) está sobre o eixo central de um espelho esférico. Para o problema 26, a Tabela 34-4 fornece:}

\[
f=20\ cm
\]

\[
m=+0,10
\]

\textbf{Determine os dados que faltam:}

\begin{itemize}
\item[(a)] o tipo de espelho;
\item[(c)] o raio de curvatura \(r\);
\item[(d)] a distância do objeto \(p\);
\item[(e)] a distância da imagem \(i\);
\item[(g)] se a imagem é real (R) ou virtual (V);
\item[(h)] se é invertida (I) ou não-invertida (NI);
\item[(i)] se a imagem está do mesmo lado do espelho que o objeto (M) ou do lado oposto (O).
\end{itemize}

\vspace{0.5cm}

A ampliação lateral é:

\[
m=-\frac{i}{p}
\]

Substituindo:

\[
0,10=-\frac{i}{p}
\]

\[
i=-0,10p
\]

\vspace{0.5cm}

Usando a equação dos espelhos:

\[
\frac{1}{f}
=
\frac{1}{p}
+
\frac{1}{i}
\]

Substituindo \(f=20\ cm\) e \(i=-0,10p\):

\[
\frac{1}{20}
=
\frac{1}{p}
+
\frac{1}{-0,10p}
\]

\[
\frac{1}{20}
=
\frac{1}{p}
-
\frac{10}{p}
\]

\[
\frac{1}{20}
=
-\frac{9}{p}
\]

\[
p=-180\ cm
\]

Então:

\[
i=-0,10(-180)
\]

\[
i=18\ cm
\]

\vspace{0.5cm}

\textbf{(a) Tipo de espelho}

Como:

\[
f>0
\]

o espelho é côncavo.

\vspace{0.5cm}

\textbf{(c) Raio de curvatura}

\[
r=2f
\]

\[
r=40\ cm
\]

\vspace{0.5cm}

\textbf{(g) Real ou virtual}

Como:

\[
i>0
\]

a imagem é real (R).

\vspace{0.5cm}

\textbf{(h) Invertida ou não-invertida}

Como:

\[
m>0
\]

a imagem é não-invertida (NI).

\vspace{0.5cm}

\textbf{(i) Lado da imagem}

Imagens reais em espelhos ficam do mesmo lado do espelho que o objeto.

Logo:

\[
\text{lado}=M
\]

\subsection{Questão 27 — Fábio Rodrigues}
\textbf{Dados}
\begin{itemize}
    \item \textbf{(a) Tipo:} Convexo
    \item \textbf{(c) Raio de curvatura (\(r\)):} A tabela informa \(40\,\text{cm}\). Como o espelho é convexo, o centro de curvatura fica atrás do espelho. Pela convenção de sinais, o raio é negativo: \(r = -40\,\text{cm}\).
    \item \textbf{(e) Distância da imagem (\(i\)):} A tabela informa \(4,0\,\text{cm}\). Um espelho convexo sempre forma imagens virtuais de objetos reais, o que significa que a imagem se forma atrás do espelho. Logo, o sinal deve ser negativo: \(i = -4,0\,\text{cm}\).
\end{itemize}

\textbf{(b) Distância focal (\(f\))}

A distância focal é a metade do raio de curvatura:
\[
f = \frac{r}{2}
\]
\[
f = \frac{-40}{2}
\]
\[
\boxed{f = -20\,\text{cm}}
\]

\textbf{(d) Distância do objeto (\(p\))}

Utilizando a equação dos pontos conjugados de Gauss:
\[
\frac{1}{f} = \frac{1}{p} + \frac{1}{i}
\]
\[
\frac{1}{-20} = \frac{1}{p} + \frac{1}{-4,0}
\]
Isolando \(\frac{1}{p}\):
\[
\frac{1}{p} = \frac{1}{4,0} - \frac{1}{20}
\]
Igualando os denominadores (multiplicando \(1/4,0\) por \(5/5\)):
\[
\frac{1}{p} = \frac{5}{20} - \frac{1}{20}
\]
\[
\frac{1}{p} = \frac{4}{20} = \frac{1}{5,0}
\]
\[
\boxed{p = +5,0\,\text{cm}}
\]

\textbf{(f) Ampliação lateral (\(m\))}

A ampliação lateral é dada por:
\[
m = -\frac{i}{p}
\]
Substituindo os valores encontrados:
\[
m = -\frac{(-4,0)}{5,0}
\]
\[
m = +\frac{4,0}{5,0}
\]
\[
\boxed{m = +0,80}
\]

\textbf{(g) Imagem Real ou Virtual (R/V)}

Como a distância da imagem é negativa (\(i < 0\)), a imagem forma-se atrás do espelho.
\[
\boxed{\text{V (Virtual)}}
\]

\textbf{(h) Imagem Invertida ou Não-invertida (I/NI)}

Como a ampliação lateral é positiva (\(m > 0\)), a orientação da imagem é a mesma do objeto.
\[
\boxed{\text{NI (Não-invertida)}}
\]

\textbf{(i) Lado da imagem}

Imagens virtuais se formam atrás do espelho, ou seja, no lado oposto ao objeto.
\[
\boxed{\text{O (Oposto)}}
\]

\subsection{Questão 28 — Mateus Valentim}
\subsubsection{Resolução}
\textbf{Dados:}

\[
\lambda = 400\,\text{nm} = 4\times10^{-7}\,\text{m}
\]

Pelo gráfico, para \(x=0\),

\[
\phi = 6\pi.
\]

A diferença de fase e a diferença de percurso estão relacionadas por

\[
\phi=\frac{2\pi}{\lambda}\Delta r.
\]

Assim,

\[
\Delta r(0)=\frac{\phi\lambda}{2\pi}
=\frac{6\pi}{2\pi}\lambda
=3\lambda.
\]

Logo, a distância entre as fontes é

\[
d=3\lambda
=3(4\times10^{-7})
=1.2\times10^{-6}\,\text{m}.
\]

Para que os raios cheguem ao detector em oposição de fase,

\[
\phi=(2m+1)\pi.
\]

Como \(\phi\) decresce com o aumento de \(x\), o maior valor de \(x\)
corresponde ao último valor ímpar possível:

\[
\phi=\pi.
\]

Portanto,

\[
\Delta r=\frac{\lambda}{2}
=\frac{4\times10^{-7}}{2}
=2\times10^{-7}\,\text{m}.
\]

Tomando \(S_1=(0,0)\), \(S_2=(0,d)\) e \(P=(x,0)\),

\[
r_1=x,
\qquad
r_2=\sqrt{x^2+d^2}.
\]

Logo,

\[
\sqrt{x^2+d^2}-x=\Delta r.
\]

Substituindo \(\Delta r=2\times10^{-7}\,\text{m}\),

\[
\sqrt{x^2+d^2}=x+\Delta r.
\]

Elevando ambos os lados ao quadrado,

\[
d^2=2x\Delta r+\Delta r^2.
\]

Assim,

\[
x=\frac{d^2-\Delta r^2}{2\Delta r}.
\]

Substituindo os valores numéricos,

\[
x=
\frac{(1.2\times10^{-6})^2-(2\times10^{-7})^2}
     {2(2\times10^{-7})}
=
3.5\times10^{-6}\,\text{m}.
\]

Portanto,

\[
\boxed{x=3.5\times10^{-6}\,\text{m}}
\]

ou, equivalentemente,

\[
\boxed{x=35\times10^{-7}\,\text{m}}.
\]

\subsection{Questão 29 — Cicero Rodrigues}

\textbf{29. }\normalfont Mais espelhos. Um objeto O está no eixo central de um espelho esférico ou plano. Para cada problema, a Tabela 34.2 mostra (a) o tipo de espelho, (b) a distância focal $f$, (c) o raio de curvatura $r$, (d) a distância do objeto $p$, (e) a distância da imagem $i$ e (f) a ampliação lateral $m$. A tabela também mostra (g) se a imagem é real (R) ou virtual (V), (h) se a imagem é invertida (I) ou não invertida (NI) e (i) se a imagem está do mesmo lado do espelho que o objeto O (M) ou do lado oposto (O). \textbf{Determine os dados que faltam para o caso 19 ($r = -40$, $i = -10$).}
\textbf{Resolução:}
\vspace{0.3cm}

\textbf{Dados:}
\[r = -40\;cm\]
\[i = -10\;cm\]

\vspace{0.2cm}
\textbf{(a) e (b) Tipo de espelho e distância focal (\(f\))}
A distância focal é sempre a metade do raio de curvatura:
\[f = \frac{r}{2} = \frac{-40}{2}\]
\[\underline{f = -20\;cm}\]
Pela convenção de sinais, um raio de curvatura negativo (\(r < 0\)) indica que o centro de curvatura está atrás do espelho. Logo, o espelho é \textbf{Convexo}.
\[\underline{\text{(a) Tipo: Convexo}}\]

\vspace{0.2cm}
\textbf{(d) Distância do objeto (\(p\))}
Aplicando a equação dos pontos conjugados (Equação de Gauss):
\[\frac{1}{f} = \frac{1}{p} + \frac{1}{i}\]
Isolando a distância do objeto (\(1/p\)):
\[\frac{1}{p} = \frac{1}{f} - \frac{1}{i}\]
\[\frac{1}{p} = \frac{1}{-20} - \frac{1}{-10}\]
\[\frac{1}{p} = -\frac{1}{20} + \frac{2}{20}\]
\[\frac{1}{p} = \frac{1}{20}\]
\[\underline{p = +20\;cm}\]

\vspace{0.2cm}
\textbf{(f) Ampliação lateral (\(m\))}
A ampliação é a razão negativa entre a distância da imagem e a do objeto:
\[m = -\frac{i}{p}\]
\[m = -\frac{-10}{20}\]
\[m = \frac{10}{20}\]
\[\underline{m = +0,5}\]

\vspace{0.2cm}
\textbf{(g), (h) e (i) Características da imagem}
Analisando os sinais matemáticos encontrados:
\begin{itemize}
    \item \textbf{(g) R/V:} Como \(i < 0\), a imagem se forma no lado virtual do espelho (atrás dele). Portanto, é \textbf{Virtual (V)}.
    \item \textbf{(h) I/NI:} Como a ampliação é positiva (\(m > 0\)), a imagem mantém a mesma orientação do objeto. Portanto, é \textbf{Não Invertida (NI)}.
    \item \textbf{(i) M/O:} Sendo virtual, a imagem está atrás da superfície refletora, enquanto o objeto real está na frente. Portanto, estão em lados opostos: \textbf{Lado Oposto (O)}.
\end{itemize}

\subsection{Questão 30 — Guilherme Araújo Floriano}

(a) Como informa a tabela, a imagem é formada do mesmo lado do espelho, o que significa
que a imagem é real (R), o espelho é côncavo e a distância focal é positiva.

(b) A distância focal é $f = +20 cm$.

(c) O raio de curvatura é $r = 2f = +40 cm$

(d) Como informa a tabela, p = +60 cm

(e) Dado $i = \frac{pf}{p-f} = +30 cm$

(f) $m = \frac{-i}{p} = -0,5$

(g) Como visto no item a, a imagem é Real (R)

(h) Como m < 0, a imagem é invertida (I)

(i) Como a imagem é real, é formada do memso lado do espelho

\subsection{Questão 31 — Arthur Roberto da Silva}
\noindent 25. (a) Como informa a tabela, a imagem é invertida (I), o que significa que o espelho é cônca\-vo, a imagem é real (R) e a ampliação lateral é negativa. De acordo com a Eq. 34-6, $i = -mp = -(-0,40)(+30 \, \text{cm}) = +12 \, \text{cm}$.

\vspace{0.3cm}

\noindent (b) $f = ip/(i + p) = +8,6 \, \text{cm}$.

\vspace{0.3cm}

\noindent (c) $r = 2f = +17,2 \, \text{cm} \approx +17 \, \text{cm}$.

\vspace{0.3cm}

\noindent (d) Como informa a tabela, $p = +30 \, \text{cm}$.

\vspace{0.3cm}

\noindent (e) Como foi visto no item (a), $i = +12 \, \text{cm}$.

\vspace{0.8cm}

\noindent (f) Como foi visto no item (a), $m = -0,40$.

\vspace{0.15cm}

\noindent (g) Como foi visto no item (a), a imagem é real (R).

\vspace{0.3cm}

\noindent (h) Como informa a tabela, a imagem é invertida (I).

\vspace{0.3cm}

\noindent (i) Como a imagem é real, é formada do mesmo lado do espelho (M).

\subsection{Questão 34 — Arthur Roberto da Silva}
Além de $n_1 = 1,0$, sabemos que (a) $n_2 = 1,5$, (c) $r = +30 \, \text{cm}$ e (d) $i = +600$.

\vspace{0.5cm}

\noindent (b) De acordo com a Eq. 34-8,
\[
p = \frac{n_1}{\frac{n_2 - n_1}{r} - \frac{n_2}{i}} = \frac{1,0}{\frac{1,5 - 1,0}{30 \, \text{cm}} - \frac{1,5}{600 \, \text{cm}}} = 71 \, \text{cm}.
\]

\vspace{0.5cm}

\noindent (e) Como $i > 0$, a imagem é real (R) e invertida.

\vspace{0.5cm}

\noindent (f) Como a imagem é real, é formada do outro lado da superfície esférica (O).

\vspace{0.5cm}

\noindent O diagrama de raios é semelhante ao da Fig. 34-12\textit{a} do livro.

\subsection{Questão 35 — Guilherme Araújo Floriano}

Além de $n_1 = 1,0$, sabemos que (a) $n_2 = 1,5$, (b) $p = +10 cm$ e (d) $i = -13 cm$

(c) 

\[
\frac{n_1}{p} + \frac{n_2}{i} = \frac{n_2 - n_1}{r}    
\]

\[
r = (n_2 - n_1)(\frac{n_1}{p} + \frac{n_2}{i})^{-1} = (1,5-1)(\frac{1}{10 cm} + \frac{1,5}{-13 cm})^{-1} = -32,5 cm
\]

(e) Como i < 0, a imagem é virtual (V)

(f) Como a imagem é virtual, é formada do mesmo lado da superfície esférica (M).

\subsection{Questão 36 — Cicero Rodrigues}

\textbf{36. }\normalfont Superfícies refratoras esféricas. Um objeto O está no eixo central de uma superfície refratora esférica. Para cada problema, a Tabela 34.3 mostra o índice de refração $n_1$ do meio em que se encontra o objeto, (a) o índice de refração $n_2$ do outro lado da superfície refratora, (b) a distância do objeto $p$, (c) o raio de curvatura $r$ da superfície e (d) a distância da imagem $i$. Determine os dados que faltam para o caso 36 ($n_1 = 1,5$; $n_2 = 1,0$; $r = -30$; $i = -7,5$), incluindo (e) se a imagem é real (R) ou virtual (V) e (f) se a imagem fica do mesmo lado da superfície que o objeto O (M) ou do lado oposto (O).
\textbf{Resolução:}
\vspace{0.3cm}

\textbf{Dados:}
\[n_1 = 1,5\]
\[n_2 = 1,0\]
\[r = -30\;cm\]
\[i = -7,5\;cm\]

\vspace{0.2cm}
\textbf{(b) Distância do objeto (\(p\))}
Para uma superfície refratora esférica, a relação entre as posições do objeto e da imagem é regida pela equação:
\[\frac{n_1}{p} + \frac{n_2}{i} = \frac{n_2 - n_1}{r}\]
Substituindo os valores numéricos fornecidos:
\[\frac{1,5}{p} + \frac{1,0}{-7,5} = \frac{1,0 - 1,5}{-30}\]
\[\frac{1,5}{p} - \frac{1}{7,5} = \frac{-0,5}{-30}\]
Simplificando o lado direito da equação:
\[\frac{1,5}{p} - \frac{1}{7,5} = \frac{1}{60}\]
Isolando o termo que contém a variável \(p\) e transformando a fração \(\frac{1}{7,5}\) em \(\frac{8}{60}\) para facilitar a soma:
\[\frac{1,5}{p} = \frac{1}{60} + \frac{1}{7,5}\]
\[\frac{1,5}{p} = \frac{1}{60} + \frac{8}{60}\]
\[\frac{1,5}{p} = \frac{9}{60}\]
Simplificando a fração (\(\frac{9}{60} = \frac{3}{20}\)):
\[\frac{1,5}{p} = \frac{3}{20}\]
Multiplicando cruzado:
\[3p = 1,5 \cdot 20\]
\[3p = 30\]
\[\underline{p = 10\;cm}\]

\vspace{0.2cm}
\textbf{(e) e (f) Características da imagem}
Analisando o sinal da distância da imagem (\(i = -7,5\;cm\)):
\begin{itemize}
    \item \textbf{(e) R/V:} Como \(i < 0\), os raios de luz refratados não convergem de fato para formar a imagem; eles divergem. Isso indica que a imagem formada é \textbf{Virtual (V)}.
    \item \textbf{(f) M/O:} Em sistemas refratores (onde a luz atravessa a fronteira), o lado "real" (positivo) é o lado oposto ao do objeto. Como a imagem é virtual (\(i < 0\)), ela se localiza antes de a luz cruzar a superfície, ou seja, fica do \textbf{Mesmo Lado (M)} em que o objeto está.
\end{itemize}

\subsection{Questão 37 — Mateus Valentim}
\subsubsection{Resolução}

Na reflexão em uma película de sabão suspensa no ar, ocorre inversão de fase na
reflexão na superfície ar-sabão, mas não ocorre inversão na reflexão na
superfície sabão-ar.

Assim, para interferência construtiva nas ondas refletidas,

\[
2nt=\left(m+\frac{1}{2}\right)\lambda,
\]

com \(m=0,1,2,\ldots\).

Logo,

\[
t=\frac{(2m+1)\lambda}{4n}.
\]

Para a menor espessura, usamos \(m=0\):

\[
t_1=\frac{\lambda}{4n}
=\frac{624\,\text{nm}}{4(1{,}33)}
=117\,\text{nm}.
\]

Para a segunda menor espessura, usamos \(m=1\):

\[
t_2=\frac{3\lambda}{4n}
=\frac{3(624\,\text{nm})}{4(1{,}33)}
=352\,\text{nm}.
\]

Portanto,

\[
\boxed{t_1=117\,\text{nm}}
\qquad \text{e} \qquad
\boxed{t_2=352\,\text{nm}}.
\]

\subsection{Questão 38 — Fábio Rodrigues}
\textbf{Dados}

\[
n_1 = 1,5
\]

\[
p = +100\,\text{cm}
\]

\[
i = +600\,\text{cm}
\]

\[
r = -30\,\text{cm}
\]

Para a refração em uma superfície esférica, utilizamos a equação:

\[
\frac{n_1}{p} + \frac{n_2}{i} = \frac{n_2 - n_1}{r}
\]

\textbf{(a) Índice de refração \(n_2\)}

Substituindo os valores conhecidos na equação:

\[
\frac{1,5}{100} + \frac{n_2}{600} = \frac{n_2 - 1,5}{-30}
\]

Para facilitar os cálculos, podemos multiplicar toda a equação por \(600\) (o mínimo múltiplo comum entre os denominadores):

\[
600 \left( \frac{1,5}{100} \right) + 600 \left( \frac{n_2}{600} \right) = 600 \left( \frac{n_2 - 1,5}{-30} \right)
\]

\[
6(1,5) + n_2 = -20(n_2 - 1,5)
\]

\[
9 + n_2 = -20n_2 + 30
\]

Agrupando os termos com \(n_2\):

\[
n_2 + 20n_2 = 30 - 9
\]

\[
21n_2 = 21
\]

\[
n_2 = \frac{21}{21}
\]

\[
\boxed{n_2 = 1,0}
\]

\textbf{(e) Imagem Real ou Virtual (R/V)}

Na refração por uma superfície esférica, o sinal da distância da imagem (\(i\)) determina a sua natureza. Como a tabela nos fornece um \(i = +600\,\text{cm}\) (valor positivo), isso significa que os raios de luz transmitidos efetivamente convergem para formar a imagem. Portanto, a imagem é real.

\[
\boxed{\text{R (Real)}}
\]

\textbf{(f) Lado da imagem}

Pela convenção de sinais para superfícies refratoras, uma distância de imagem positiva (\(i > 0\)) indica que a imagem se forma no lado para o qual a luz está se propagando, ou seja, no lado oposto ao do objeto.

\[
\boxed{\text{O (Oposto)}}
\]

\subsection{Questão 39 — Rodrigo Albuquerque}


\textbf{Superfícies refratoras esféricas. Um objeto \(O\) está sobre o eixo central de uma superfície refratora esférica. Para o problema 39, a Tabela 34-5 fornece:}

\[
n_1=1,5
\]

\[
n_2=1,0
\]

\[
p=+10\ cm
\]

\[
i=-6,0\ cm
\]

\textbf{Determine:}

\begin{itemize}
\item[(c)] o raio de curvatura \(r\);
\item[(e)] se a imagem é real (R) ou virtual (V);
\item[(f)] se a imagem fica do mesmo lado da superfície que o objeto (M) ou do lado oposto (O).
\end{itemize}

\vspace{0.5cm}

Para superfícies refratoras esféricas:

\[
\frac{n_1}{p}+\frac{n_2}{i}
=
\frac{n_2-n_1}{r}
\]

Substituindo os valores:

\[
\frac{1,5}{10}
+
\frac{1,0}{-6,0}
=
\frac{1,0-1,5}{r}
\]

\[
0,15-0,1667
=
\frac{-0,5}{r}
\]

\[
-0,0167
=
\frac{-0,5}{r}
\]

\[
r=\frac{-0,5}{-0,0167}
\]

\[
r\approx+30\ cm
\]

\vspace{0.5cm}

\textbf{(e) Real ou virtual}

Como:

\[
i<0
\]

a imagem é virtual (V).

\vspace{0.5cm}

\textbf{(f) Lado da imagem}

Imagens virtuais ficam do mesmo lado da superfície que o objeto.

Logo:

\[
\text{lado}=M
\]

\subsection{Questão 40 — Dávyla Maria}
Inserir resolução aqui.

\subsection{Questão 50 — Kaylane Castro}


\subsection*{Enunciado}

Um objeto $O$ está sobre o eixo central de uma lente delgada simétrica.

Dados:

\[
p = +10\,\text{cm}
\]

Lente divergente com módulo da distância focal:

\[
|f| = 6,0\,\text{cm}
\]

Logo:

\[
f = -6,0\,\text{cm}
\]

Determine:

\begin{itemize}
    \item[\textbf{(a)}] A distância da imagem $i$;
    \item[\textbf{(b)}] A ampliação lateral $m$;
    \item[\textbf{(c)}] Se a imagem é real (R) ou virtual (V);
    \item[\textbf{(d)}] Se a imagem é invertida (I) ou não-invertida (NI);
    \item[\textbf{(e)}] Se a imagem está do mesmo lado da lente que o objeto (M) ou do lado oposto (O).
\end{itemize}

% ---------------------------------------------------------

\subsection*{Dados}

\[
p = +10\,\text{cm}
\]

\[
f = -6,0\,\text{cm}
\]

% ---------------------------------------------------------

\subsection*{Distância da Imagem}

Aplicando a equação de Gauss para lentes delgadas:

\[
\frac{1}{f}
=
\frac{1}{p}
+
\frac{1}{i}
\]

Substituindo os valores:

\[
\frac{1}{-6,0}
=
\frac{1}{10}
+
\frac{1}{i}
\]

Isolando $\frac{1}{i}$:

\[
\frac{1}{i}
=
-\frac{1}{6,0}
-
\frac{1}{10}
\]

Calculando:

\[
\frac{1}{i}
=
\frac{-5 - 3}{30}
\]

\[
\frac{1}{i}
=
-\frac{8}{30}
\]

\[
i
=
-\frac{30}{8}
\]

\[
i = -3,75\,\text{cm}
\]

% ---------------------------------------------------------

\subsection*{Ampliação Lateral}

A ampliação lateral é dada por:

\[
m = -\frac{i}{p}
\]

Substituindo os valores:

\[
m
=
-\frac{-3,75}{10}
\]

\[
m = +0,375
\]

\[
m \approx +0,38
\]

% ---------------------------------------------------------

\subsection*{Natureza da Imagem}

Como:

\[
i < 0
\]

a imagem é:

\[
\boxed{\text{Virtual}}
\]

% ---------------------------------------------------------

\subsection*{Orientação da Imagem}

Como:

\[
m > 0
\]

a imagem é:

\[
\boxed{\text{Não-Invertida}}
\]

% ---------------------------------------------------------

\subsection*{Lado da Imagem}

Para lentes delgadas:

\begin{itemize}
    \item $i > 0$ $\rightarrow$ imagem real;
    \item $i < 0$ $\rightarrow$ imagem virtual.
\end{itemize}

Imagens virtuais formam-se do mesmo lado da lente que o objeto.

Logo:

\[
\boxed{\text{M}}
\]

% ---------------------------------------------------------

\subsection*{Tabela Final}

\begin{center}
\begin{tabular}{ccccccc}
\hline

$p$ & Lente & (a) $i$ & (b) $m$
& (c) R/V & (d) I/NI & (e) Lado \\

\hline

$+10$ & D, $6,0$
& $-3,8$
& $+0,38$
& V
& NI
& M \\

\hline
\end{tabular}
\end{center}

% ---------------------------------------------------------

\subsection*{Respostas}

\[
\boxed{
i = -3,75\,\text{cm}
}
\]

\[
\boxed{
m \approx +0,38
}
\]

\[
\boxed{
\text{Imagem Virtual e Não-Invertida}
}
\]

\[
\boxed{
\text{Mesmo lado da lente (M)}
}
\]
```


\subsection{Questão 51 — Ana Alicy}
Inserir resolução aqui.

\subsection{Questão 52 — Ana Alicy}
Inserir resolução aqui.

\subsection{Questão 53 — Kaylane Castro}


\subsection*{Enunciado}

Um objeto $O$ está sobre o eixo central de uma lente delgada simétrica.

Dados:

\[
p = +12\,\text{cm}
\]

Lente convergente com distância focal:

\[
f = +16\,\text{cm}
\]

Determine:

\begin{itemize}
    \item[\textbf{(a)}] A distância da imagem $i$;
    \item[\textbf{(b)}] A ampliação lateral $m$;
    \item[\textbf{(c)}] Se a imagem é real (R) ou virtual (V);
    \item[\textbf{(d)}] Se a imagem é invertida (I) ou não-invertida (NI);
    \item[\textbf{(e)}] Se a imagem está do mesmo lado da lente que o objeto (M) ou do lado oposto (O).
\end{itemize}

% ---------------------------------------------------------

\subsection*{Dados}

\[
p = +12\,\text{cm}
\]

\[
f = +16\,\text{cm}
\]

% ---------------------------------------------------------

\subsection*{Distância da Imagem}

Aplicando a equação de Gauss para lentes delgadas:

\[
\frac{1}{f}
=
\frac{1}{p}
+
\frac{1}{i}
\]

Substituindo os valores:

\[
\frac{1}{16}
=
\frac{1}{12}
+
\frac{1}{i}
\]

Isolando $\frac{1}{i}$:

\[
\frac{1}{i}
=
\frac{1}{16}
-
\frac{1}{12}
\]

\[
\frac{1}{i}
=
\frac{3 - 4}{48}
\]

\[
\frac{1}{i}
=
-\frac{1}{48}
\]

Logo:

\[
i = -48\,\text{cm}
\]

% ---------------------------------------------------------

\subsection*{Ampliação Lateral}

A ampliação lateral é dada por:

\[
m = -\frac{i}{p}
\]

Substituindo os valores:

\[
m
=
-\frac{-48}{12}
\]

\[
m = +4,0
\]

% ---------------------------------------------------------

\subsection*{Natureza da Imagem}

Como:

\[
i < 0
\]

a imagem é:

\[
\boxed{\text{Virtual}}
\]

% ---------------------------------------------------------

\subsection*{Orientação da Imagem}

Como:

\[
m > 0
\]

a imagem é:

\[
\boxed{\text{Não-Invertida}}
\]

% ---------------------------------------------------------

\subsection*{Lado da Imagem}

Para lentes delgadas:

\begin{itemize}
    \item $i > 0$ $\rightarrow$ imagem real;
    \item $i < 0$ $\rightarrow$ imagem virtual.
\end{itemize}

Imagens virtuais formam-se do mesmo lado da lente que o objeto.

Logo:

\[
\boxed{\text{M}}
\]

% ---------------------------------------------------------

\subsection*{Tabela Final}

\begin{center}
\begin{tabular}{ccccccc}
\hline

$p$ & Lente & (a) $i$ & (b) $m$
& (c) R/V & (d) I/NI & (e) Lado \\

\hline

$+12$ & C, $16$
& $-48$
& $+4,0$
& V
& NI
& M \\

\hline
\end{tabular}
\end{center}

% ---------------------------------------------------------

\subsection*{Respostas}

\[
\boxed{
i = -48\,\text{cm}
}
\]

\[
\boxed{
m = +4,0
}
\]

\[
\boxed{
\text{Imagem Virtual e Não-Invertida}
}
\]

\[
\boxed{
\text{Mesmo lado da lente (M)}
}
\]
```


\subsection{Questão 54 — Dávyla Maria}
Inserir resolução aqui.

\subsection{Questão 55 — Rodrigo Albuquerque}

Dados da tabela:

\[
p = 45\text{ cm}
\]

Lente convergente \(C,20\)

Como a lente é convergente:

\[
f = +20\text{ cm}
\]

Passo 1 — Fórmula da lente delgada

\[
\frac{1}{f} = \frac{1}{p} + \frac{1}{i}
\]

Substituindo os valores:

\[
\frac{1}{20} = \frac{1}{45} + \frac{1}{i}
\]

Isolando \(\frac{1}{i}\):

\[
\frac{1}{i} = \frac{1}{20} - \frac{1}{45}
\]

Passo 2 — Fazendo as frações

MMC entre \(20\) e \(45\):

\[
180
\]

Então:

\[
\frac{1}{20} = \frac{9}{180}
\]

\[
\frac{1}{45} = \frac{4}{180}
\]

Substituindo:

\[
\frac{1}{i} = \frac{9}{180} - \frac{4}{180}
\]

\[
\frac{1}{i} = \frac{5}{180}
\]

Simplificando:

\[
\frac{1}{i} = \frac{1}{36}
\]

Invertendo:

\[
i = 36\text{ cm}
\]

(a) Distância da imagem

\[
\boxed{i = +36\text{ cm}}
\]

O sinal positivo indica que a imagem é real.

Passo 3 — Ampliação lateral

\[
m = -\frac{i}{p}
\]

Substituindo:

\[
m = -\frac{36}{45}
\]

\[
m = -0,8
\]

(b) Ampliação

\[
\boxed{m = -0,8}
\]

O sinal negativo indica imagem invertida.

(c) Real ou virtual?

Como:

\[
i > 0
\]

A imagem é:

\[
\boxed{\text{Real (R)}}
\]

(d) Invertida ou não invertida?

Como:

\[
m < 0
\]

A imagem é:

\[
\boxed{\text{Invertida (I)}}
\]

(e) Mesmo lado ou lado oposto?

Imagem real em lentes delgadas fica no lado oposto da lente:

\[
\boxed{\text{Lado oposto (O)}}
\]
```


\subsection{Questão 56 — Fábio Rodrigues}
\textbf{Enunciado}

Lentes delgadas. Um objeto \(O\) está sobre o eixo central de uma lente delgada simétrica. A distância do objeto é \(p = +25\,\text{cm}\), e a lente é do tipo "C, 35" (onde C significa convergente e 35 é a distância focal em centímetros). Determine (a) a distância da imagem \(i\) e (b) a ampliação lateral \(m\) do objeto, incluindo os sinais. Determine também (c) se a imagem é real (R) ou virtual (V), (d) se é invertida (I) ou não-invertida (NI) e (e) se está do mesmo lado da lente que o objeto \(O\) (M) ou do lado oposto (O).

\vspace{0.5cm}

\textbf{Dados}

\[
p = +25\,\text{cm}
\]

Lente: C, 35. O "C" indica que é uma lente convergente, logo a distância focal é positiva:

\[
f = +35\,\text{cm}
\]

\textbf{(a) Distância da imagem (\(i\))}

Para lentes delgadas, utilizamos a equação de Gauss:

\[
\frac{1}{f} = \frac{1}{p} + \frac{1}{i}
\]

Substituindo os valores:

\[
\frac{1}{35} = \frac{1}{25} + \frac{1}{i}
\]

Isolando \(\frac{1}{i}\):

\[
\frac{1}{i} = \frac{1}{35} - \frac{1}{25}
\]

Para subtrair as frações, encontramos o mínimo múltiplo comum (MMC) entre 35 e 25, que é 175:

\[
\frac{1}{i} = \frac{5}{175} - \frac{7}{175}
\]

\[
\frac{1}{i} = -\frac{2}{175}
\]

Invertendo ambos os lados:

\[
i = -\frac{175}{2}
\]

\[
\boxed{i = -87,5\,\text{cm}}
\]

\textbf{(b) Ampliação lateral (\(m\))}

A ampliação lateral é dada pela fórmula:

\[
m = -\frac{i}{p}
\]

Substituindo os valores encontrados:

\[
m = -\frac{(-87,5)}{25}
\]

\[
m = +\frac{87,5}{25}
\]

\[
\boxed{m = +3,5}
\]

\textbf{(c) Imagem Real ou Virtual (R/V)}

Como a distância da imagem é negativa (\(i < 0\)), isso indica que a imagem formada pela lente é virtual.

\[
\boxed{\text{V (Virtual)}}
\]

\textbf{(d) Invertida ou Não-invertida (I/NI)}

Como a ampliação lateral é positiva (\(m > 0\)), a imagem possui a mesma orientação do objeto, ou seja, é direita (não-invertida).

\[
\boxed{\text{NI (Não-invertida)}}
\]

\textbf{(e) Lado da imagem}

Em lentes delgadas, imagens virtuais (\(i < 0\)) formam-se do mesmo lado de onde a luz incide, que é o lado onde o objeto está posicionado.

\[
\boxed{\text{M (Mesmo)}}
\]

\subsection{Questão 57 — Mateus Valentim}
Inserir resolução aqui.

\subsection{Questão 59 — Cicero Rodrigues}

\textbf{59. }\normalfont Lentes com raios dados. Um objeto O está no eixo central de uma lente delgada. Para cada problema, a Tabela 34.5 mostra a distância do objeto $p$, o índice de refração $n$ da lente, o raio $r_1$ da superfície da lente mais próxima do objeto e o raio $r_2$ da superfície da lente mais distante do objeto. Determine \textbf{(a)} a distância da imagem $i$ e \textbf{(b)} a ampliação lateral $m$ do objeto, incluindo o sinal. Determine também se \textbf{(c)} se a imagem é real (R) ou virtual (V), \textbf{(d)} se é invertida (I) ou não invertida (NI) e \textbf{(e)} se está do mesmo lado da lente que o objeto O (M) ou do lado oposto (O). \textbf{Caso 59:} $p = +75$; $n = 1,55$; $r_1 = +30$; $r_2 = -42$.

\textbf{Resolução:}
\vspace{0.3cm}

\textbf{Dados:}
\[p = +75\;cm\]
\[n = 1,55\]
\[r_1 = +30\;cm\]
\[r_2 = -42\;cm\]

\vspace{0.2cm}
\textbf{Cálculo prévio da Distância Focal (\(f\))}
Utilizamos a Equação dos Fabricantes de Lentes:
\[\frac{1}{f} = (n - 1) \left( \frac{1}{r_1} - \frac{1}{r_2} \right)\]
Substituindo os valores:
\[\frac{1}{f} = (1,55 - 1) \left( \frac{1}{30} - \frac{1}{-42} \right)\]
\[\frac{1}{f} = 0,55 \left( \frac{1}{30} + \frac{1}{42} \right)\]
Tirando o MMC entre 30 e 42 (que é 210):
\[\frac{1}{f} = 0,55 \left( \frac{7}{210} + \frac{5}{210} \right) = \frac{55}{100} \left( \frac{12}{210} \right)\]
\[\frac{1}{f} = \frac{11}{20} \cdot \frac{2}{35} = \frac{22}{700} = \frac{11}{350}\]
\[f = \frac{350}{11} \approx +31,8\;cm\]

\vspace{0.2cm}
\textbf{(a) Distância da imagem (\(i\))}
Aplicamos a Equação de Gauss para lentes delgadas:
\[\frac{1}{p} + \frac{1}{i} = \frac{1}{f}\]
Isolando \(\frac{1}{i}\):
\[\frac{1}{i} = \frac{1}{f} - \frac{1}{p}\]
Usando o valor exato em fração para \(1/f\):
\[\frac{1}{i} = \frac{11}{350} - \frac{1}{75}\]
O MMC entre 350 e 75 é 1050:
\[\frac{1}{i} = \frac{11 \cdot 3}{1050} - \frac{14}{1050}\]
\[\frac{1}{i} = \frac{33 - 14}{1050} = \frac{19}{1050}\]
\[i = \frac{1050}{19} \approx +55,263\;cm\]
\[\underline{i \approx +55,3\;cm}\]

\vspace{0.2cm}
\textbf{(b) Ampliação lateral (\(m\))}
\[m = -\frac{i}{p}\]
Usando a fração exata para evitar erros de arredondamento (\(i = 1050/19\)):
\[m = -\frac{1050/19}{75} = -\frac{1050}{19 \cdot 75}\]
\[m = -\frac{14}{19} \approx -0,7368\]
\[\underline{m \approx -0,737}\]

\vspace{0.2cm}
\textbf{(c), (d) e (e) Características da imagem}
Analisando os sinais encontrados:
\begin{itemize}
    \item \textbf{(c) R/V:} Como \(i > 0\), a luz cruza a lente e converge de fato no espaço imagem. Portanto, a imagem é \textbf{Real (R)}.
    \item \textbf{(d) I/NI:} Como \(m < 0\), a imagem aponta para a direção oposta à do objeto. Portanto, é \textbf{Invertida (I)}.
    \item \textbf{(e) M/O:} Para lentes, a imagem real se forma após a luz atravessar o dispositivo ótico. Portanto, a imagem está no \textbf{Lado Oposto (O)} ao objeto.
\end{itemize}

\subsection{Questão 60 — Guilherme Araújo Floriano}

(a) 

\[   
\frac{1}{p} + \frac{1}{i} = (n-1)*(\frac{1}{r_1} - \frac{1}{r_2})
\]

\[
\frac{1}{29} + \frac{1}{i} = (1,65-1)*(\frac{1}{35} - \frac{1}{\infty})
\]

\[
i = (\frac{0,65}{35} - \frac{1}{29})^{-1} = -63cm
\]

(b)

\[
m = -\frac{i}{p} = +2,2  
\]

(c) Como i < 0, a imagem é virtual (V)

(d) Como m > 0, a imagem é não invertida (NI)

(e) Como a imagem é virtual, é formada do mesmo lado da lente (M)

\subsection{Questão 62 — Arthur Roberto da Silva}
\noindent 66. (a) Combinando a Eq. 34-9 com a Eq. 34-10, obtemos $i = -9,7 \, \text{cm}$.

\vspace{0.4cm}

\noindent (b) De acordo com a Eq. 34-7, $m = -i/p = +0,54$.

\vspace{0.4cm}

\noindent (c) Como $i < 0$, a imagem é virtual (V).

\vspace{0.4cm}

\noindent (d) Como $m > 0$, a imagem é não invertida (NI).

\vspace{0.4cm}

\noindent (e) Como a imagem é virtual, é formada do mesmo lado da lente (M).

\subsection{Questão 63 — Guilherme Araújo Floriano}

(a) 

\[   
\frac{1}{p} + \frac{1}{i} = (n-1)*(\frac{1}{r_1} - \frac{1}{r_2})
\]

\[
\frac{1}{60} + \frac{1}{i} = (1,5-1)*(\frac{1}{35} - \frac{1}{-35})
\]

\[
i = (\frac{0,5*2}{35} - \frac{1}{60})^{-1} = +84cm
\]

(b)

\[
m = -\frac{i}{p} = - 1,4 
\]

(c) Como i > 0, a imagem é real (R)

(d) Como m < 0, a imagem é invertida (I)

(e) Como a imagem é real, é formada do lado oposto da lente (O)

\subsection{Questão 64 — Cicero Rodrigues}

\textbf{64. }\normalfont Lentes com raios dados. Um objeto O está no eixo central de uma lente delgada. Para cada problema, a Tabela 34.5 mostra a distância do objeto $p$, o índice de refração $n$ da lente, o raio $r_1$ da superfície da lente mais próxima do objeto e o raio $r_2$ da superfície da lente mais distante do objeto. Determine \textbf{(a)} a distância da imagem $i$ e \textbf{(b)} a ampliação lateral $m$ do objeto, incluindo o sinal. Determine também se \textbf{(c)} se a imagem é real (R) ou virtual (V), \textbf{(d)} se é invertida (I) ou não invertida (NI) e \textbf{(e)} se está do mesmo lado da lente que o objeto O (M) ou do lado oposto (O). \textbf{Caso 64:} $p = +10$; $n = 1,50$; $r_1 = -30$; $r_2 = -60$.
\textbf{Resolução:}
\vspace{0.3cm}

\textbf{Dados:}
\[p = +10\;cm\]
\[n = 1,50\]
\[r_1 = -30\;cm\]
\[r_2 = -60\;cm\]

\vspace{0.2cm}
\textbf{Cálculo prévio da Distância Focal (\(f\))}
Utilizamos a Equação dos Fabricantes de Lentes:
\[\frac{1}{f} = (n - 1) \left( \frac{1}{r_1} - \frac{1}{r_2} \right)\]
Substituindo os valores:
\[\frac{1}{f} = (1,50 - 1) \left( \frac{1}{-30} - \frac{1}{-60} \right)\]
\[\frac{1}{f} = 0,50 \left( -\frac{2}{60} + \frac{1}{60} \right)\]
\[\frac{1}{f} = 0,50 \left( -\frac{1}{60} \right)\]
\[\frac{1}{f} = -\frac{0,50}{60} = -\frac{1}{120}\]
\[f = -120\;cm\]

\vspace{0.2cm}
\textbf{(a) Distância da imagem (\(i\))}
Aplicamos a Equação de Gauss para lentes delgadas:
\[\frac{1}{p} + \frac{1}{i} = \frac{1}{f}\]
Isolando \(\frac{1}{i}\):
\[\frac{1}{i} = \frac{1}{f} - \frac{1}{p}\]
\[\frac{1}{i} = \frac{1}{-120} - \frac{1}{10}\]
Tirando o MMC entre 120 e 10 (que é 120):
\[\frac{1}{i} = -\frac{1}{120} - \frac{12}{120}\]
\[\frac{1}{i} = -\frac{13}{120}\]
\[i = -\frac{120}{13} \approx -9,2307\;cm\]
\[\underline{i \approx -9,23\;cm}\]

\vspace{0.2cm}
\textbf{(b) Ampliação lateral (\(m\))}
\[m = -\frac{i}{p}\]
Usando a fração exata para evitar erros de arredondamento (\(i = -120/13\)):
\[m = -\frac{-120/13}{10}\]
\[m = \frac{120/13}{10} = \frac{120}{130} = \frac{12}{13}\]
\[m \approx +0,9230\]
\[\underline{m \approx +0,923}\]

\vspace{0.2cm}
\textbf{(c), (d) e (e) Características da imagem}
Analisando os sinais matemáticos encontrados:
\begin{itemize}
    \item \textbf{(c) R/V:} Como \(i < 0\), a luz diverge ao passar pela lente e a imagem é formada pelo prolongamento dos raios. Portanto, é uma imagem \textbf{Virtual (V)}.
    \item \textbf{(d) I/NI:} Como \(m > 0\), a imagem mantém a mesma orientação do objeto original. Portanto, é \textbf{Não Invertida (NI)}.
    \item \textbf{(e) M/O:} Por ser virtual, a imagem se forma no espaço objeto, antes da luz atravessar a lente. Logo, está do \textbf{Mesmo Lado (M)} da lente que o objeto O.
\end{itemize}

\subsection{Questão 65 — Mateus Valentim}
Inserir resolução aqui.

\subsection{Questão 66 — Fábio Rodrigues}
\textbf{Enunciado}

Lentes com raios dados. Um objeto \(O\) está sobre o eixo central de uma lente delgada. A distância do objeto é \(p = +10\,\text{cm}\), o índice de refração da lente é \(n = 1,50\), o raio da superfície da lente mais próxima do objeto é \(r_1 = -30\,\text{cm}\) e o raio da superfície da lente mais distante do objeto é \(r_2 = -60\,\text{cm}\). Determine (a) a distância da imagem \(i\) e (b) a ampliação lateral \(m\) do objeto, incluindo o sinal. Determine também (c) se a imagem é real (R) ou virtual (V), (d) se é invertida (I) ou não-invertida (NI) e (e) se está do mesmo lado da lente que o objeto \(O\) (M) ou do lado oposto (O).

\vspace{0.5cm}

\textbf{Dados}

\[
p = +10\,\text{cm}
\]

\[
n = 1,50
\]

\[
r_1 = -30\,\text{cm}
\]

\[
r_2 = -60\,\text{cm}
\]

\textbf{(a) Distância da imagem (\(i\))}

Primeiro, calculamos a distância focal (\(f\)) da lente utilizando a equação dos fabricantes de lentes:

\[
\frac{1}{f} = (n - 1) \left( \frac{1}{r_1} - \frac{1}{r_2} \right)
\]

Substituindo os valores:

\[
\frac{1}{f} = (1,50 - 1) \left( \frac{1}{-30} - \frac{1}{-60} \right)
\]

\[
\frac{1}{f} = 0,50 \left( -\frac{2}{60} + \frac{1}{60} \right)
\]

\[
\frac{1}{f} = 0,50 \left( -\frac{1}{60} \right)
\]

\[
\frac{1}{f} = -\frac{1}{120}
\]

\[
f = -120\,\text{cm}
\]

Agora, usamos a equação de Gauss para encontrar a distância da imagem (\(i\)):

\[
\frac{1}{f} = \frac{1}{p} + \frac{1}{i}
\]

\[
-\frac{1}{120} = \frac{1}{10} + \frac{1}{i}
\]

Isolando \(\frac{1}{i}\):

\[
\frac{1}{i} = -\frac{1}{120} - \frac{1}{10}
\]

Encontrando o mínimo múltiplo comum (120):

\[
\frac{1}{i} = -\frac{1}{120} - \frac{12}{120}
\]

\[
\frac{1}{i} = -\frac{13}{120}
\]

Invertendo ambos os lados:

\[
i = -\frac{120}{13} \approx -9,23\,\text{cm}
\]

\[
\boxed{i \approx -9,23\,\text{cm}}
\]

\textbf{(b) Ampliação lateral (\(m\))}

A ampliação lateral é dada pela fórmula:

\[
m = -\frac{i}{p}
\]

Substituindo usando o valor em fração para maior precisão:

\[
m = -\frac{\left(-\frac{120}{13}\right)}{10}
\]

\[
m = +\frac{120}{130} = +\frac{12}{13} \approx +0,923
\]

\[
\boxed{m \approx +0,92}
\]

\textbf{(c) Imagem Real ou Virtual (R/V)}

Como a distância da imagem é negativa (\(i < 0\)), a imagem não pode ser projetada e é formada pelo prolongamento dos raios.

\[
\boxed{\text{V (Virtual)}}
\]

\textbf{(d) Invertida ou Não-invertida (I/NI)}

Como a ampliação lateral é positiva (\(m > 0\)), a imagem possui a mesma orientação do objeto (é direita).

\[
\boxed{\text{NI (Não-invertida)}}
\]

\textbf{(e) Lado da imagem}

Em lentes delgadas, as imagens virtuais (\(i < 0\)) sempre se formam do mesmo lado de onde a luz incide (o lado do objeto).

\[
\boxed{\text{M (Mesmo)}}
\]


\subsection{Questão 68 — Dávyla Maria}
Inserir resolução aqui.

\subsection{Questão 69 — Kaylane Castro}


\subsection*{Enunciado}

Um objeto $O$ está sobre o eixo central de uma lente delgada simétrica.

Dados:

\[
p = +16\,\text{cm}
\]

\[
m = +0,25
\]

Determine:

\begin{itemize}
    \item[\textbf{(a)}] O tipo de lente, convergente (C) ou divergente (D);
    \item[\textbf{(b)}] A distância focal $f$;
    \item[\textbf{(d)}] A distância da imagem $i$;
    \item[\textbf{(f)}] Se a imagem é real (R) ou virtual (V);
    \item[\textbf{(g)}] Se a imagem é invertida (I) ou não-invertida (NI);
    \item[\textbf{(h)}] Se a imagem está do mesmo lado da lente que o objeto (M) ou do lado oposto (O).
\end{itemize}

% ---------------------------------------------------------

\subsection*{Dados}

\[
p = +16\,\text{cm}
\]

\[
m = +0,25
\]

% ---------------------------------------------------------

\subsection*{Orientação da Imagem}

Como:

\[
m > 0
\]

a imagem é:

\[
\boxed{\text{Não-Invertida}}
\]

% ---------------------------------------------------------

\subsection*{Distância da Imagem}

A ampliação lateral é dada por:

\[
m = -\frac{i}{p}
\]

Substituindo os valores:

\[
0,25 = -\frac{i}{16}
\]

Multiplicando ambos os lados por $16$:

\[
i = -0,25 \times 16
\]

\[
i = -4,0\,\text{cm}
\]

% ---------------------------------------------------------

\subsection*{Natureza da Imagem}

Como:

\[
i < 0
\]

a imagem é:

\[
\boxed{\text{Virtual}}
\]

% ---------------------------------------------------------

\subsection*{Distância Focal}

Aplicando a equação de Gauss para lentes delgadas:

\[
\frac{1}{f}
=
\frac{1}{p}
+
\frac{1}{i}
\]

Substituindo os valores:

\[
\frac{1}{f}
=
\frac{1}{16}
+
\frac{1}{-4,0}
\]

\[
\frac{1}{f}
=
\frac{1}{16}
-
\frac{4}{16}
\]

\[
\frac{1}{f}
=
-\frac{3}{16}
\]

Logo:

\[
f = -\frac{16}{3}
\]

\[
f \approx -5,33\,\text{cm}
\]

% ---------------------------------------------------------

\subsection*{Tipo de Lente}

Como:

\[
f < 0
\]

a lente é:

\[
\boxed{\text{Divergente}}
\]

% ---------------------------------------------------------

\subsection*{Lado da Imagem}

Imagens virtuais formam-se do mesmo lado da lente que o objeto.

Logo:

\[
\boxed{\text{M}}
\]

% ---------------------------------------------------------

\subsection*{Tabela Final}

\begin{center}
\begin{tabular}{cccccccc}
\hline

(a) Tipo & (b) $f$ & (c) $p$ & (d) $i$
& (e) $m$ & (f) R/V & (g) I/NI & (h) Lado \\

\hline

D & $-5,3$ & $+16$
& $-4,0$
& $+0,25$
& V
& NI
& M \\

\hline
\end{tabular}
\end{center}

% ---------------------------------------------------------

\subsection*{Respostas}

\[
\boxed{
f \approx -5,33\,\text{cm}
}
\]

\[
\boxed{
i = -4,0\,\text{cm}
}
\]

\[
\boxed{
\text{Lente Divergente}
}
\]

\[
\boxed{
\text{Imagem Virtual e Não-Invertida}
}
\]
```


\subsection{Questão 70 — Ana Alicy}
Inserir resolução aqui.

\subsection{Questão 88 — Arthur Roberto }
\noindent De acordo com a Eq. 34-15, o diâmetro mínimo da ocular é
\[
d_{\text{oc}} = \frac{d_{\text{ob}}}{m_\theta} = \frac{75 \, \text{mm}}{36} = 2,1 \, \text{mm}.
\]

\subsection{Questão 90 — Rodrigo Albuquerque}

Dados:

\[
p_o = 10,0\ \text{mm}
\]

(distância do objeto até a objetiva)

\[
L = 300\ \text{mm}
\]

(distância entre as lentes)

A imagem intermediária forma-se a:

\[
50,0\ \text{mm}
\]

da ocular.

---

A distância da imagem produzida pela objetiva é:

\[
i_o = 300 - 50
\]

\[
i_o = 250\ \text{mm}
\]

---

A ampliação linear da objetiva é:

\[
m_o = -\frac{i_o}{p_o}
\]

Substituindo:

\[
m_o = -\frac{250}{10}
\]

\[
m_o = -25
\]

---

Para a ocular, usamos a ampliação angular:

\[
M_e = \frac{25\ \text{cm}}{f_e}
\]

Como a imagem intermediária está a \(50,0\ \text{mm}\) da ocular, tomamos:

\[
f_e = 50,0\ \text{mm} = 5,0\ \text{cm}
\]

Então:

\[
M_e = \frac{25}{5}
\]

\[
M_e = 5
\]

---

A ampliação total do microscópio é:

\[
M = m_o \cdot M_e
\]

\[
M = (-25)(5)
\]

\[
M = -125
\]

---

Logo, a ampliação total do microscópio é:

\[
\boxed{125}
\]

(o sinal negativo apenas indica que a imagem final é invertida)


\end{document}
